\chapter{Redes Neuronales Artificiales}\label{ch:neural_networks}

Las redes neuronales artificiales (\gls{ANN}) representan una de las primeras y más exitosas aproximaciones al cómputo inspirado biológicamente. Su desarrollo se inició en 1943 cuando Warren McCulloch y Walter Pitts propusieron el primer modelo matemático formal de una neurona artificial~\citep{McCullochPitts1943}, estableciendo los fundamentos conceptuales que perduran hasta hoy.

La motivación central de las \gls{ANN} es emular, de forma simplificada, el procesamiento de información de las neuronas biológicas. Estas células especializadas reciben señales de entrada, las integran mediante procesos electroquímicos complejos, y generan respuestas de salida que se propagan a través de redes interconectadas. Crucialmente, las neuronas biológicas exhiben plasticidad sináptica, adaptando sus respuestas en función de la experiencia y permitiendo el aprendizaje a partir de ejemplos.

Esta analogía con los mecanismos de procesamiento del cerebro humano ha motivado décadas de investigación en inteligencia artificial, buscando desarrollar sistemas computacionales capaces de aprender y adaptarse de manera similar a los sistemas biológicos. En este capítulo examinaremos los fundamentos matemáticos de las \gls{ANN}, desde la neurona individual hasta las arquitecturas profundas modernas.

\section{Neurona Artificial}\label{sec:neurona_artificial}

\begin{figure}[h]
	\centering
	\includegraphics[width=0.6\textwidth]{figures/neuron_biological}
	\caption[Esquema de una neurona biológica]{Esquema de una neurona biológica mostrando dendritas (entradas), soma (cuerpo celular), y axón (salida). Esta estructura inspiró el modelo matemático de McCulloch-Pitts.}
	\label{fig:neuron}
\end{figure}

El modelo de McCulloch-Pitts estableció las bases conceptuales al abstraer el comportamiento de una neurona biológica (\autoref{fig:neuron}) en términos matemáticos simples. En el sistema nervioso, las neuronas reciben señales electroquímicas de otras células a través de las dendritas y conexiones sinápticas. Estas señales se integran en el soma celular, y si la suma ponderada supera un umbral crítico de despolarización, la neurona genera un potencial de acción que se propaga a través del axón hacia otras neuronas.

Inspirándose en este mecanismo, una neurona artificial constituye un modelo matemático que:
\begin{enumerate}
	\item Recibe múltiples señales de entrada numéricas
	\item Las pondera mediante parámetros ajustables (pesos sinápticos)
	\item Aplica una función de activación no lineal que modela el umbral de disparo
	\item Produce una señal de salida que puede servir como entrada a otras neuronas
\end{enumerate}

Es importante reconocer que las neuronas biológicas reales exhiben una complejidad que va mucho más allá de esta abstracción inicial. Los procesos neuronales involucran dinámicas temporales complejas, neurotransmisores químicos específicos, plasticidad sináptica dependiente del tiempo, y arquitecturas dendríticas tridimensionales intrincadas. Sin embargo, la simplificación del modelo de McCulloch-Pitts ha demostrado ser extraordinariamente fructífera, proporcionando la base conceptual para el desarrollo de las \gls{ANN} modernas.

\subsection{Formulación matemática}

El modelo matemático estándar de una neurona artificial puede formularse como una función que mapea un vector de entradas \(\mathbf{x} = (x_1, x_2, \ldots, x_n)^T\) a una salida escalar \(y\). La operación fundamental consiste en:

\begin{equation}
	y = \sigma\left(\sum_{i=1}^{n} w_i x_i + b\right)
	\label{eq:neuron}
\end{equation}

donde:
\begin{itemize}
	\item \(\mathbf{w} = (w_1, w_2, \ldots, w_n)^T\) son los pesos sinápticos que modulan la importancia de cada entrada
	\item \(b\) es el sesgo o umbral de activación (bias)
	\item \(\sigma\) es la función de activación no lineal
\end{itemize}

La función de activación \(\sigma\) desempeña un papel fundamental al introducir la no linealidad esencial que permite a las redes neuronales aproximar funciones complejas y capturar relaciones no lineales en los datos. Los pesos sinápticos \(\mathbf{w}\) y el sesgo \(b\) constituyen los parámetros ajustables del modelo. Estos se actualizan durante el proceso de entrenamiento para aproximar la función objetivo deseada.

\subsubsection{Redes Neuronales de Pulsos}

Es importante mencionar que esta formulación clásica, aunque extremadamente exitosa, representa una abstracción simplificada del comportamiento neuronal. Existen modelos alternativos más cercanos a la biología, como las \textbf{Redes Neuronales de Pulsos} o \gls{SNN}, consideradas la ``tercera generación'' de redes neuronales \citep{Maass1997spiking}. A diferencia del modelo clásico donde la salida es un valor continuo, las \glspl{SNN} procesan información mediante pulsos discretos (spikes) que se propagan en el tiempo, capturando así la dinámica temporal inherente a las neuronas biológicas. Estas redes utilizan modelos como el \textit{Integrate-and-Fire} o el \textit{Hodgkin-Huxley}, donde la neurona acumula potencial de membrana hasta alcanzar un umbral y emitir un pulso. Las \glspl{SNN} ofrecen ventajas computacionales significativas en términos de eficiencia energética y son especialmente prometedoras para hardware neuromórfico \citep{Tavanaei2018deep}. Sin embargo, para los propósitos de este libro, nos centraremos en el modelo clásico de neurona artificial descrito anteriormente.

\subsubsection{Definición formal}

Para establecer un marco matemático riguroso para este libro, formalizamos primero el concepto de función de activación y, posteriormente, el de neurona artificial.

\begin{definition}[Función de activación]\label{def:funcion_activacion}
	Una \textit{función de activación} es una función $\sigma: \mathbb{R} \to \mathbb{R}$ que satisface las siguientes propiedades:
	\begin{itemize}
		\item \textbf{No linealidad}: $\sigma$ no es una aplicación afín,
		\item \textbf{Continuidad}: $\sigma$ es continua,
		\item \textbf{Diferenciabilidad casi por doquier}: $\sigma$ es diferenciable excepto posiblemente en un conjunto de medida de Lebesgue nula\footnote{Un conjunto $A \subset \mathbb{R}$ tiene medida de Lebesgue nula si puede ser cubierto por intervalos cuya suma de longitudes es arbitrariamente pequeña. Véase las Definiciones~\ref{def:medida_lebesgue} y \ref{def:casi_por_doquier} para la definición formal.}.
	\end{itemize}
\end{definition}

Estas propiedades aseguran que $\sigma$ pueda introducir no linealidades necesarias para la capacidad expresiva de la red, mientras permite el cálculo de gradientes para el entrenamiento mediante retropropagación. La \autoref{tab:funciones_activacion} resume las principales funciones de activación utilizadas en la práctica junto con sus propiedades más relevantes, y en la \autoref{subsubsec:funciones_activacion_capas} se describen en detalle cada una de ellas.

\begin{definition}[Neurona artificial]\label{def:neurona_artificial}
	Sea $n \in \mathbb{N}$ y sea $\sigma: \mathbb{R} \to \mathbb{R}$ una función de activación. Para cada $\theta = (\mathbf{w}, b) \in \mathbb{R}^{n+1}$ con $\mathbf{w} \in \mathbb{R}^n$ (pesos) y $b \in \mathbb{R}$ (sesgo), una \textit{neurona artificial} es una función $f_{\theta}: \mathbb{R}^n \to \mathbb{R}$ que se expresa como la composición de una aplicación afín con una función de activación:
	\begin{equation}
		f_{\theta}(\mathbf{x}) = \sigma \circ T_{\theta}(\mathbf{x}) = \sigma\Big( \langle \mathbf{w}, \mathbf{x} \rangle + b \Big)
		\label{eq:neuron_formal}
	\end{equation}
	donde $T_{\theta}: \mathbb{R}^n \to \mathbb{R}$ es la \textit{aplicación afín} definida por
	\[
	T_{\theta}(\mathbf{x}) = \langle \mathbf{w}, \mathbf{x} \rangle + b.
	\]
\end{definition}

El conjunto de parámetros $\theta$ es entrenable, lo que significa que durante el proceso de aprendizaje se ajusta para minimizar una función de pérdida. Así, una neurona artificial define una \textit{familia de funciones} $\{f_{\theta} : \theta \in \mathbb{R}^{n+1}\}$ parametrizada por los pesos y el sesgo.

La \autoref{fig:neurona_artificial} ilustra la estructura de una neurona artificial, mostrando el flujo desde las entradas ponderadas hasta la salida activada.

\begin{figure}[h]
\centering
\begin{tikzpicture}[
    node distance=0.8cm,
    input/.style={circle, draw, minimum size=0.6cm, fill=gdlblue!15},
    weight/.style={rectangle, draw, minimum width=0.6cm, minimum height=0.4cm, fill=gdlgray!15, font=\scriptsize},
    sum/.style={circle, draw, minimum size=0.8cm, fill=gdlyellow!20},
    activation/.style={rectangle, draw, rounded corners, minimum width=0.8cm, minimum height=0.6cm, fill=gdlgreen!20},
    output/.style={circle, draw, minimum size=0.6cm, fill=gdlred!15},
    arrow/.style={->, thick, >=stealth}
]

% Entradas
\node[input] (x1) at (0, 1.5) {$x_1$};
\node[input] (x2) at (0, 0) {$x_2$};
\node[input] (xn) at (0, -1.5) {$x_n$};
\node at (0, -0.75) {$\vdots$};

% Pesos
\node[weight] (w1) at (1.5, 1.5) {$w_1$};
\node[weight] (w2) at (1.5, 0) {$w_2$};
\node[weight] (wn) at (1.5, -1.5) {$w_n$};

% Nodo sumador
\node[sum] (suma) at (3.5, 0) {$\sum$};
\node[below=0.1cm of suma, font=\scriptsize] {$+b$};

% Función de activación
\node[activation] (sigma) at (5.5, 0) {$\sigma$};

% Salida
\node[output] (y) at (7, 0) {$y$};

% Conexiones
\draw[arrow] (x1) -- (w1);
\draw[arrow] (x2) -- (w2);
\draw[arrow] (xn) -- (wn);
\draw[arrow] (w1) -- (suma);
\draw[arrow] (w2) -- (suma);
\draw[arrow] (wn) -- (suma);
\draw[arrow] (suma) -- node[above, font=\scriptsize] {$z$} (sigma);
\draw[arrow] (sigma) -- (y);

% Etiquetas de fases
\node[above=1.5cm of x1, font=\small\bfseries] {Entradas};
\node[above=0.3cm of w1, font=\small\bfseries] {Pesos};
\node[above=0.5cm of suma, font=\small\bfseries] {Suma};
\node[above=0.5cm of sigma, font=\small\bfseries] {Activación};

\end{tikzpicture}
\caption[Estructura de una neurona artificial]{Estructura de una neurona artificial. Las entradas $x_i$ se multiplican por pesos $w_i$, se suman junto con el sesgo $b$, y el resultado $z = \sum w_i x_i + b$ se transforma mediante la función de activación $\sigma$ para producir la salida $y = \sigma(z)$.}
\label{fig:neurona_artificial}
\end{figure}

La no linealidad de $\sigma$ es crucial: la composición de múltiples aplicaciones afines sigue siendo una aplicación afín, por lo que sin funciones de activación no lineales, una red neuronal profunda con múltiples capas sería equivalente a una única transformación lineal, perdiendo toda su capacidad expresiva. Es esta no linealidad la que permite a las redes neuronales aproximar funciones arbitrariamente complejas, como establece el \autoref{thm:aproximacion_universal} de aproximación universal (véase la \autoref{sec:aproximacion_universal}).

\subsubsection{Propiedades deseables de funciones de activación}

Aunque la \autoref{def:funcion_activacion} establece las propiedades mínimas necesarias para que una función pueda servir como función de activación, en la práctica es deseable que cumpla propiedades adicionales que facilitan el entrenamiento y mejoran el comportamiento de las redes neuronales. A continuación presentamos un análisis detallado de estas propiedades deseables.

\begin{definition}[Acotación o crecimiento controlado]\label{prop:acotacion_activacion}
	Una función de activación $\sigma: \mathbb{R} \to \mathbb{R}$ tiene \textit{crecimiento controlado} si satisface una de las siguientes condiciones:
	\begin{enumerate}
		\item \textbf{Acotación}: Existe $M > 0$ tal que $|\sigma(z)| \leq M$ para todo $z \in \mathbb{R}$.
		\item \textbf{Crecimiento lineal}: Existen constantes $C_1, C_2 > 0$ tales que $|\sigma(z)| \leq C_1|z| + C_2$ para todo $z \in \mathbb{R}$.
	\end{enumerate}
\end{definition}

Esta propiedad es crucial para la estabilidad numérica durante el entrenamiento. La primera condición (acotación) garantiza que las activaciones permanecen en un rango finito independientemente de la magnitud de la entrada, impidiendo todo desbordamiento numérico. La segunda condición (crecimiento lineal) es más permisiva: permite que las activaciones crezcan, pero a una tasa no mayor que la de la propia entrada. En ambos casos, se previene la explosión de valores durante la propagación hacia adelante, que de otro modo podría causar desbordamientos numéricos (overflow) o gradientes excesivamente grandes. La \autoref{tab:funciones_activacion} muestra que las funciones de activación habituales satisfacen una u otra condición; para descripciones detalladas de cada función, véase la \autoref{subsubsec:funciones_activacion_capas}.

\begin{definition}[Monotonía]\label{prop:monotonia_activacion}
	Una función $\sigma: \mathbb{R} \to \mathbb{R}$ es \textit{monótona} (creciente) si para cualesquiera $z_1, z_2 \in \mathbb{R}$ con $z_1 < z_2$ se cumple:
	\[
	\sigma(z_1) \leq \sigma(z_2)
	\]
\end{definition}

La monotonía proporciona tres ventajas fundamentales:
\begin{enumerate}
	\item \textbf{Preservación del orden}: Las activaciones mantienen el orden relativo de las entradas, facilitando la interpretabilidad.
	\item \textbf{Unicidad de óptimos locales}: En combinación con convexidad de la función de pérdida, simplifica el paisaje de optimización.
	\item \textbf{Suavidad de la superficie de pérdida}: Reduce la presencia de mínimos locales espurios en el espacio de parámetros.
\end{enumerate}

La mayoría de funciones de activación clásicas son monótonas (véase la \autoref{subsubsec:funciones_activacion_capas}). Sin embargo, la monotonía no es condición necesaria para un buen desempeño empírico: existen funciones de activación no monótonas que exhiben resultados competitivos en la práctica, lo que sugiere que las ventajas teóricas de la monotonía pueden ser compensadas por otras propiedades favorables.

\begin{definition}[Centrada en cero]\label{prop:zero_centered}
	Una función de activación $\sigma: \mathbb{R} \to \mathbb{R}$ está \textit{centrada en cero} si satisface:
	\begin{enumerate}
		\item $\sigma(0) = 0$ o $\sigma(0) \approx 0$
		\item Es simétrica respecto al origen: $\sigma(-z) = -\sigma(z)$ para todo $z \in \mathbb{R}$ (función impar)
	\end{enumerate}
\end{definition}

Esta propiedad acelera significativamente la convergencia del entrenamiento. Cuando las activaciones están centradas en cero, los gradientes tienen media cercana a cero, lo que evita sesgos sistemáticos en las actualizaciones de pesos.

Cuando una función de activación no cumple esta propiedad (por ejemplo, si su rango no está centrado en el origen, si $\sigma(0) \neq 0$), las salidas de cada capa presentan un sesgo sistemático. Esto provoca que todos los gradientes respecto a los pesos de una misma capa tengan el mismo signo, causando el fenómeno de ``zigzagging'' durante el descenso de gradiente: las actualizaciones de parámetros oscilan entre direcciones ortogonales en lugar de converger directamente hacia el óptimo.

\begin{definition}[Gradiente no saturante]\label{prop:gradiente_no_saturante}
	Una función de activación $\sigma: \mathbb{R} \to \mathbb{R}$ tiene \textit{gradiente no saturante} si existe $\delta > 0$ y un conjunto $S \subseteq \mathbb{R}$ de medida infinita tal que:
	\[
	|\sigma'(z)| \geq \delta \quad \text{para todo } z \in S
	\]
\end{definition}

Esta propiedad es fundamental para evitar el problema del \textit{desvanecimiento del gradiente} (vanishing gradient), uno de los obstáculos principales en el entrenamiento de redes profundas (véase el análisis formal en la \autoref{subsubsec:backpropagation}).

\begin{corollary}
Toda función de activación acotada con asíntotas horizontales necesariamente satura.
\end{corollary}

\begin{proof}
Si $\lim_{z \to \pm\infty} \sigma(z) = c_{\pm}$ con $c_+, c_-$ finitos, entonces $\lim_{z \to \pm\infty} \sigma'(z) = 0$. Para cualquier $\delta > 0$, existe $M > 0$ tal que $|\sigma'(z)| < \delta$ para todo $|z| > M$, de modo que el conjunto $\{z : |\sigma'(z)| \geq \delta\} \subseteq [-M, M]$ tiene medida finita. Por la \autoref{prop:gradiente_no_saturante}, $\sigma$ no puede tener gradiente no saturante y, por tanto, necesariamente satura.
\end{proof}

Cuando $|z|$ es grande, $|\sigma'(z)| \approx 0$, lo que causa que los gradientes se vuelvan exponencialmente pequeños en capas profundas durante la retropropagación. Matemáticamente, si consideramos una red con $L$ capas donde cada capa tiene activación saturada, el gradiente de la pérdida respecto a los pesos de las primeras capas será proporcional a $\prod_{\ell=1}^{L} \sigma'(z^{(\ell)})$, que puede volverse numéricamente cero.

Una función que mantiene $|\sigma'(z)| \geq \delta > 0$ en un conjunto de medida infinita permite que el gradiente fluya sin atenuación a través de múltiples capas. El caso óptimo corresponde a funciones cuya derivada es constante igual a 1 en una semirrecta, lo que elimina completamente la saturación en esa región activa y facilita el entrenamiento de redes muy profundas. Más aún, si se exige $|\sigma'(z)| > 0$ para todo $z \in \mathbb{R}$ (es decir, gradiente estrictamente no nulo en todo el dominio), se elimina el riesgo de que ciertas neuronas dejen de contribuir al aprendizaje (fenómeno conocido como ``neuronas muertas'').

La \autoref{fig:vanishing_gradient} (véase la \autoref{subsubsec:backpropagation}) ilustra visualmente cómo el gradiente se desvanece exponencialmente en redes profundas con activaciones saturantes, comparado con el comportamiento estable de activaciones no saturantes.

\begin{definition}[Eficiencia computacional]\label{prop:eficiencia_computacional}
	Una función de activación es \textit{computacionalmente eficiente} si tanto $\sigma(z)$ como su derivada $\sigma'(z)$ pueden evaluarse en tiempo $O(1)$ con operaciones aritméticas básicas.
\end{definition}

En redes neuronales profundas con millones de neuronas, la función de activación se evalúa miles de millones de veces durante el entrenamiento. Por tanto, su eficiencia computacional tiene un impacto directo en el tiempo total de entrenamiento.

Taxonomía por tipo de operación:
\begin{itemize}
	\item \textbf{Comparaciones y selecciones} (funciones definidas a trozos mediante $\max$): Requieren únicamente operaciones de comparación y selección, lo que las hace extremadamente eficientes en hardware moderno.
	\item \textbf{Funciones trascendentes} (exponenciales, logaritmos): Se evalúan en $O(1)$ pero con una constante significativamente mayor que las anteriores, al requerir aproximaciones polinómicas internas.
	\item \textbf{Funciones vectoriales} (normalización sobre todo el vector de salida): Tienen complejidad $O(n)$ en la dimensión del vector, ya que requieren calcular estadísticas globales antes de normalizar cada componente.
\end{itemize}

Aunque todas estas categorías son polinomiales en complejidad, las diferencias de constantes son significativas en la práctica. Las funciones de activación más simples computacionalmente permiten entrenar redes más profundas en menos tiempo, lo que ha sido un factor relevante en la adopción generalizada de activaciones basadas en comparaciones.

\begin{definition}[Continuidad de Lipschitz]\label{prop:lipschitz_continuidad}
	Una función $\sigma: \mathbb{R} \to \mathbb{R}$ es \textit{Lipschitz continua} si existe una constante $L > 0$ (llamada \textit{constante de Lipschitz}) tal que para cualesquiera $z_1, z_2 \in \mathbb{R}$:
	\[
	|\sigma(z_1) - \sigma(z_2)| \leq L |z_1 - z_2|
	\]
\end{definition}

La continuidad de Lipschitz es más fuerte que la continuidad uniforme y proporciona control cuantitativo sobre cómo las perturbaciones en la entrada se propagan a la salida. Esto tiene ciertas ventajas en el entrenamiento de redes de neuronas:

\begin{enumerate}
	\item \textbf{Estabilidad numérica}: Pequeños cambios en los pesos producen cambios acotados en las activaciones.
	\item \textbf{Control de gradientes}: Si $\sigma$ es Lipschitz con constante $L$, entonces $|\sigma'(z)| \leq L$ casi por doquier (cuando la derivada existe).
	\item \textbf{Garantías de convergencia}: Muchos teoremas de convergencia de algoritmos de optimización asumen Lipschitz continuidad de la función objetivo.
\end{enumerate}

Para funciones de activación diferenciables casi por doquier, la constante de Lipschitz se puede calcular como $L = \sup_{z} |\sigma'(z)|$. Funciones acotadas con asíntotas horizontales tienden a tener constantes de Lipschitz pequeñas (ya que su derivada se anula en los extremos y alcanza un máximo moderado), mientras que funciones con crecimiento lineal típicamente presentan $L = 1$. La \autoref{tab:funciones_activacion} recoge los valores concretos para las funciones de activación más comunes.

Cabe observar que una función puede ser Lipschitz continua sin ser de clase $C^1$: basta con que sea continua y tenga derivada acotada casi por doquier. En los puntos de no diferenciabilidad se utilizan subgradientes, lo que resulta suficiente para los algoritmos de optimización basados en gradiente empleados en el entrenamiento de redes neuronales.

Las propiedades presentadas no son mutuamente excluyentes, y las funciones de activación más exitosas en la práctica suelen combinar varias de ellas. La \autoref{tab:funciones_activacion} resume las propiedades de las funciones de activación más utilizadas.

\begin{table}[h]
\centering
\caption{Comparación de funciones de activación comunes}
\label{tab:funciones_activacion}
\begin{tabular}{@{}lccccc@{}}
\toprule
\textbf{Función} & \textbf{Fórmula} & \textbf{Rango} & \textbf{Centrada} & \textbf{No satura} & \textbf{Lipschitz} \\
\midrule
Sigmoide & $\frac{1}{1+e^{-z}}$ & $(0,1)$ & No & No & $L=\frac{1}{4}$ \\
\addlinespace
Tanh & $\frac{e^z-e^{-z}}{e^z+e^{-z}}$ & $(-1,1)$ & Sí & No & $L=1$ \\
\addlinespace
ReLU & $\max(0,z)$ & $[0,\infty)$ & No & Sí ($z>0$) & $L=1$ \\
\addlinespace
Leaky ReLU & $\max(\alpha z,z)$ & $\mathbb{R}$ & No & Sí & $L=1$ \\
\addlinespace
Swish & $z \cdot \sigma(z)$ & $\approx(-0.3,\infty)$ & No & Parcial & $L\approx 1.1$ \\
\bottomrule
\end{tabular}
\end{table}

La \autoref{fig:activation_functions} muestra gráficamente las funciones de activación más comunes junto con sus derivadas, permitiendo visualizar las zonas de saturación y el comportamiento del gradiente.

\begin{figure}[htbp]
\centering
\includegraphics[width=0.95\textwidth]{figures/activation_functions}
\caption[Gráficos de funciones de activación]{Funciones de activación comunes y sus derivadas. Las zonas sombreadas indican regiones de saturación donde el gradiente tiende a cero. Nótese cómo ReLU mantiene gradiente constante para $z > 0$, mientras que Sigmoide y Tanh saturan en los extremos.}
\label{fig:activation_functions}
\end{figure}

Un resultado teórico fundamental que justifica el uso de funciones de activación con estas propiedades es el \textbf{teorema de aproximación universal} (véase la \autoref{sec:aproximacion_universal}), que establece que redes neuronales con una sola capa oculta y funciones de activación apropiadas pueden aproximar cualquier función continua en un compacto con precisión arbitraria. Este teorema será presentado y demostrado en detalle más adelante en este capítulo, junto con los preliminares de análisis funcional necesarios para su comprensión completa.

La elección de la función de activación apropiada depende del problema específico, la arquitectura de la red, y consideraciones prácticas de entrenamiento. En las secciones siguientes veremos cómo estas propiedades influyen en el diseño de arquitectura.

\section{Datos y su estructura}\label{sec:datos_estructura}

En la sección anterior definimos la neurona artificial como una función que procesa datos de entrada y genera datos de salida. Sin embargo, en esta formulación no hemos hablado de la información o los datos que procesa la red de neuronas, no nos hemos parado a pensar qué representa realmente el espacio $\mathbb{R}^n$. Es necesario entender si los datos del espacio $\mathbb{R}^n$ son simplemente números o si hay algún tipo de relación entre ellos que contiene información relevante.

Esta distinción es crucial. Cuando representamos una imagen de $28 \times 28$ píxeles como un vector en $\mathbb{R}^{784}$, no estamos tratando con $784$ valores independientes: cada componente corresponde a un píxel específico en una grilla bidimensional, y los píxeles espacialmente vecinos están naturalmente relacionados. La estructura de grilla no es arbitraria, sino intrínseca a la naturaleza de la imagen.

En esta sección formalizaremos esta distinción mediante dos marcos conceptuales complementarios que nos permitirán unificar el tratamiento de datos euclidianos clásicos (vectores, imágenes, secuencias) con datos geométricos complejos (grafos, variedades, nubes de puntos). En los siguientes apartados presentaremos dos marcos matemáticamente equivalentes pero con énfasis conceptuales distintos: el triple \textit{dominio-estructura-señal} $(\Omega, \mathcal{S}, f)$ y el par \textit{geometría-atributos} $(\mathcal{D}, f)$.

\subsection{Señales sobre dominios estructurados}\label{sec:senales_dominios}

El primer marco conceptual se basa en tres componentes fundamentales que separan claramente las nociones de posiciones, relaciones y valores. Esta descomposición tripartita constituye la base del paradigma unificador del \gls{GDL} desarrollado en el \autoref{ch:gdl}, donde las ``cinco G's'' de \citet{bronstein2017geometricdeeplearning} (geometría, grupos, grafos, geodésicas y gauges) extienden y enriquecen estos conceptos fundamentales.

\begin{definition}[Dominio de datos]\label{def:dominio_datos}
Un \textit{dominio de datos} es un conjunto $\Omega$ cuyos elementos representan las posiciones o ubicaciones donde se observan los datos. Puede ser:
\begin{itemize}
	\item \textbf{Discreto y finito}: $\Omega = \{p_1, \ldots, p_n\}$ (vértices de un grafo, píxeles de una imagen, índices de un vector)
	\item \textbf{Discreto e infinito}: $\Omega = \mathbb{Z}^d$ (grilla infinita)
	\item \textbf{Continuo}: $\Omega \subseteq \mathbb{R}^d$ o una variedad diferenciable $\mathcal{M}$
\end{itemize}
\end{definition}

\begin{definition}[Estructura del dominio]\label{def:estructura_dominio}
La \textit{estructura} $\mathcal{S}$ sobre un dominio $\Omega$ define las relaciones intrínsecas entre sus elementos. Puede ser de naturaleza:
\begin{itemize}
	\item \textbf{Topológica}: Una topología $\tau$ que define conjuntos abiertos y nociones de vecindad
	\item \textbf{Métrica}: Una función distancia $d: \Omega \times \Omega \to \mathbb{R}_{\geq 0}$ que cuantifica separaciones entre puntos
	\item \textbf{Algebraica}: Simetrías representadas mediante un grupo $G$ que actúa sobre $\Omega$
	\item \textbf{Combinatoria}: Un grafo $(V, E)$ que define conexiones explícitas entre puntos
	\item \textbf{Geométrica}: Una variedad diferenciable con atlas, métrica riemanniana, conexión, etc.
\end{itemize}

El par $(\Omega, \mathcal{S})$ constituye un \textit{dominio estructurado}.
\end{definition}

Cada tipo de estructura sobre el dominio se corresponde con una rama de la matemática que proporciona herramientas específicas para su análisis. En el \autoref{ch:gdl} profundizaremos en estas estructuras y en cómo las simetrías del dominio (formalizadas mediante grupos) determinan las restricciones de equivarianza que deben satisfacer las arquitecturas neuronales.

\begin{definition}[Señal]\label{def:senal}
Una \textit{señal} sobre un dominio estructurado $(\Omega, \mathcal{S})$ es una función
\[
f: \Omega \to \mathcal{C}
\]
que asigna a cada punto $p \in \Omega$ un vector de características $f(p) \in \mathcal{C}$, donde $\mathcal{C}$ es un espacio vectorial denominado \textit{espacio de canales}. Típicamente $\mathcal{C} = \mathbb{R}^d$, donde $d$ es el número de canales o atributos.

Cuando $\Omega$ es finito, $\Omega = \{p_1, \ldots, p_n\}$, y $\mathcal{C} = \mathbb{R}^d$, podemos representar la señal mediante la matriz:
\[
F = \begin{bmatrix}
	f(p_1)^T \\
	f(p_2)^T \\
	\vdots \\
	f(p_n)^T
\end{bmatrix} \in \mathbb{R}^{n \times d}
\]
donde cada fila $f(p_i) \in \mathbb{R}^d$ contiene las características observadas en la posición $p_i$.
\end{definition}

\begin{remark}\label{rem:espacio_canales_general}
La elección de $\mathcal{C} = \mathbb{R}^d$ es suficiente para la mayoría de las aplicaciones que trataremos en este capítulo. Sin embargo, en capítulos posteriores el espacio de canales adquiere una estructura más rica: en el \autoref{ch:gdl}, $\mathcal{C}$ puede ser un espacio de representación de un grupo de simetría $G$ (por ejemplo, armónicos esféricos para $SO(3)$), y en el contexto de redes gauge-equivariantes, $\mathcal{C}$ corresponde a la fibra de un fibrado vectorial asociado sobre la variedad base.
\end{remark}

\begin{definition}[Espacio de señales]\label{def:espacio_senales}
Dado un dominio estructurado $(\Omega, \mathcal{S})$ y un espacio de canales $\mathcal{C}$, el \textit{espacio de señales} $\mathcal{F}(\Omega, \mathcal{C})$ es el conjunto de todas las señales sobre $\Omega$ con valores en $\mathcal{C}$:
\[
\mathcal{F}(\Omega, \mathcal{C}) = \{f: \Omega \to \mathcal{C}\}.
\]
Cuando $\Omega$ es finito con $|\Omega| = n$ y $\mathcal{C} = \mathbb{R}^d$, el espacio de señales es isomorfo a un espacio matricial: $\mathcal{F}(\Omega, \mathbb{R}^d) \cong \mathbb{R}^{n \times d}$, donde cada señal se identifica con su representación matricial $F$. Cuando $\Omega$ es continuo, el espacio $\mathcal{F}(\Omega, \mathcal{C})$ se restringe a un espacio funcional con condiciones de regularidad apropiadas (por ejemplo, $L^2(\Omega, \mathcal{C})$ o $C^k(\Omega, \mathcal{C})$).
\end{definition}

Una red neuronal puede verse entonces como un operador entre espacios de señales. Una capa transforma señales de entrada en señales de salida, posiblemente sobre dominios y espacios de canales diferentes:
\[
T: \mathcal{F}(\Omega, \mathcal{C}) \to \mathcal{F}(\Omega', \mathcal{C}').
\]
Esta perspectiva funcional será fundamental en los capítulos sobre redes convolucionales (\autoref{ch:cnn}), donde la convolución es un operador lineal equivariante entre espacios de señales, y sobre mecanismos de atención (\autoref{ch:attention}), donde la atención implementa un operador no lineal adaptativo.

En resumen:
\begin{itemize}
	\item El \textbf{dominio} $\Omega$: define \textit{dónde} están los datos (las posiciones)
	\item La \textbf{estructura} $\mathcal{S}$: define \textit{cómo se relacionan} las posiciones entre sí
	\item La \textbf{señal} $f$: define los \textit{valores observados} en cada posición
\end{itemize}

La \autoref{fig:dominio_estructura_senal} ilustra visualmente este marco conceptual y cómo los tres componentes se combinan para formar datos completos.

\begin{figure}[h]
\centering
\begin{tikzpicture}[
    scale=0.75,
    transform shape,
    box/.style={rectangle, draw, rounded corners, minimum width=3cm, minimum height=1.2cm, align=center},
    arrow/.style={->, thick, >=stealth},
    label/.style={font=\footnotesize, align=center}
]

% Dominio
\node[box, fill=gdlblue!15] (dom) at (0, 2) {\textbf{Dominio} $\Omega$\\{\scriptsize Posiciones/ubicaciones}};
\node[label, left=0.1cm of dom] {Píxeles, vértices, puntos};

% Estructura
\node[box, fill=gdlgreen!15] (str) at (6, 2) {\textbf{Estructura} $\mathcal{S}$\\{\scriptsize Relaciones entre posiciones}};
\node[label, right=0.1cm of str] {Grilla, grafo, métrica};

% Geometría
\node[box, fill=gdlyellow!20] (geo) at (3, 0) {\textbf{Geometría} $(\Omega, \mathcal{S})$\\{\scriptsize Dominio estructurado}};

% Señal
\node[box, fill=gdlorange!18] (sig) at (8, 0) {\textbf{Señal} $f: \Omega \to \mathcal{C}$\\{\scriptsize Valores observados}};
\node[label, right=0.1cm of sig] {RGB, características, atributos};

% Datos completos
\node[box, fill=gdlred!18, minimum width=4cm] (dat) at (5.5, -2) {\textbf{Datos completos}\\{\scriptsize $(\Omega, \mathcal{S}, f)$}};

% Flechas
\draw[arrow] (dom) -- (geo);
\draw[arrow] (str) -- (geo);
\draw[arrow] (geo) -- (dat);
\draw[arrow] (sig) -- (dat);

\end{tikzpicture}
\caption[Marco conceptual dominio-estructura-señal]{Marco conceptual para la representación de datos. El dominio $\Omega$ define las posiciones, la estructura $\mathcal{S}$ define las relaciones entre posiciones, y la señal $f$ asigna valores a cada posición. Juntos forman los datos completos que procesa una red neuronal.}
\label{fig:dominio_estructura_senal}
\end{figure}

Un vector $\mathbf{x} \in \mathbb{R}^n$ no es meramente una colección de $n$ números reales: representa una señal $f$ sobre un dominio $\Omega$ con estructura $\mathcal{S}$ específica. Las relaciones entre componentes $x_i$ y $x_j$ no son arbitrarias, sino que vienen determinadas por la estructura del dominio subyacente.

\begin{example}[Señal de audio (temporal)]\label{ex:senal_audio}
Consideremos una señal de audio representada formalmente como:
\begin{itemize}
	\item \textbf{Dominio}: $\Omega = [0, T] \subset \mathbb{R}$ (intervalo temporal continuo)
	\item \textbf{Estructura}: $\mathcal{S} = $ orden total de $\mathbb{R}$ con topología estándar
	\item \textbf{Señal}: $f(t) \in \mathbb{R}$ (amplitud sonora en el instante $t$)
\end{itemize}

En cada instante temporal $t$, observamos una amplitud $f(t)$. La estructura es el orden temporal: $t_1 < t_2$ tiene significado causal. Las señales de audio respetan localidad temporal: valores en instantes cercanos están típicamente correlacionados.
\end{example}

\begin{example}[Imagen digital en grilla]\label{ex:imagen_grilla}
Una imagen digital se estructura como:
\begin{itemize}
	\item \textbf{Dominio}: $\Omega = \{1,\ldots,H\} \times \{1,\ldots,W\}$
	\item \textbf{Estructura}: $\mathcal{S} = $ grilla 2D rectangular con vecindad local (4 u 8 conexos)
	\item \textbf{Señal}: $f(i,j) \in \mathbb{R}^3$ (intensidades RGB del píxel $(i,j)$)
\end{itemize}

Si permutamos arbitrariamente el orden de píxeles al vectorizar la imagen, la representación vectorial cambia, pero la imagen (como señal sobre grilla 2D) permanece invariante. La estructura de grilla es intrínseca al objeto geométrico, no a su representación vectorial particular.
\end{example}

\begin{example}[Red social (grafo)]\label{ex:red_social}
Una red social se modela mediante:
\begin{itemize}
	\item \textbf{Dominio}: $\Omega = V = \{u_1, \ldots, u_n\}$ (conjunto de usuarios)
	\item \textbf{Estructura}: $\mathcal{S} = (V, E)$ donde $E \subseteq V \times V$ representa conexiones o amistades
	\item \textbf{Señal}: $f(u) \in \mathbb{R}^d$ (perfil del usuario $u$: edad, intereses, actividad, etc.)
\end{itemize}

No existe orden canónico de usuarios. La ``posición'' de un usuario no es un índice numérico, sino su identidad abstracta en el grafo. Las relaciones están definidas por las aristas $E$, no por proximidad de índices en alguna vectorización arbitraria. Dos usuarios $u_i, u_j$ son ``vecinos'' si y solo si $(u_i, u_j) \in E$.
\end{example}

\begin{example}[Molécula (grafo geométrico 3D)]\label{ex:molecula}
Una molécula combina estructura discreta y geometría continua:
\begin{itemize}
	\item \textbf{Dominio}: $\Omega = \{a_1, \ldots, a_n\}$ (conjunto de átomos)
	\item \textbf{Estructura}: $\mathcal{S} = $ (grafo de enlaces químicos) $+$ (posiciones en $\mathbb{R}^3$) $+$ (simetrías bajo grupo euclidiano $E(3)$)
	\item \textbf{Señal}: $f(a) \in \mathbb{R}^d$ (tipo atómico, carga parcial, masa atómica, hibridación, etc.)
\end{itemize}

La molécula combina conectividad discreta (enlaces químicos covalentes) con geometría continua (posiciones tridimensionales de átomos). Las propiedades físico-químicas son invariantes bajo transformaciones del grupo euclidiano $E(3)$: traslaciones, rotaciones y reflexiones no cambian la identidad molecular.
\end{example}

\begin{example}[Nube de puntos 3D]\label{ex:nube_puntos}
Una nube de puntos tridimensional se caracteriza por:
\begin{itemize}
	\item \textbf{Dominio}: $\Omega = \{p_1, \ldots, p_n\} \subset \mathbb{R}^3$ (conjunto finito de puntos en el espacio)
	\item \textbf{Estructura}: $\mathcal{S} = $ métrica euclidiana $d(p_i, p_j) = \|p_i - p_j\|_2$
	\item \textbf{Señal}: $f(p) \in \mathbb{R}^d$ (normales de superficie, curvaturas, colores, reflectancia, intensidad LIDAR, etc.)
\end{itemize}

No hay orden natural de los puntos. Permutar la lista $\{p_1,\ldots,p_n\}$ no cambia la nube de puntos como objeto geométrico tridimensional. La estructura es puramente métrica: solo importan las distancias relativas entre puntos, no su orden de enumeración.
\end{example}

\begin{example}[Superficie 3D (variedad riemanniana)]\label{ex:superficie_riemanniana}
Una superficie 3D como variedad riemanniana se define mediante:
\begin{itemize}
	\item \textbf{Dominio}: $\Omega = \mathcal{M}$ (variedad diferenciable 2D embebida en $\mathbb{R}^3$, por ejemplo, superficie de una proteína o un objeto 3D)
	\item \textbf{Estructura}: $\mathcal{S} = $ métrica riemanniana $g$ inducida por el embedding en $\mathbb{R}^3$
	\item \textbf{Señal}: $f: \mathcal{M} \to \mathbb{R}^d$ (campos escalares o vectoriales sobre la superficie: temperatura, hidrofilicidad, curvatura gaussiana, etc.)
\end{itemize}

Las distancias geodésicas sobre $\mathcal{M}$ y la curvatura son propiedades intrínsecas de la variedad, independientes del embedding particular en $\mathbb{R}^3$. Dos parametrizaciones diferentes de la misma superficie definen la misma geometría intrínseca.
\end{example}

\begin{example}[Secuencia de texto (NLP)]\label{ex:secuencia_texto}
Una secuencia de texto para procesamiento del lenguaje natural se representa como:
\begin{itemize}
	\item \textbf{Dominio}: $\Omega = \{1, 2, \ldots, L\}$ (posiciones discretas en una secuencia de longitud $L$)
	\item \textbf{Estructura}: $\mathcal{S} = $ orden total con estructura secuencial bidireccional (cada posición $i$ puede atender a posiciones anteriores y posteriores)
	\item \textbf{Señal}: $f(i) \in \mathbb{R}^d$ (embedding vectorial del token en la posición $i$: representación semántica de palabras, subpalabras o caracteres)
\end{itemize}

A diferencia del audio unidimensional con causalidad estricta, el texto permite dependencias bidireccionales: el significado de una palabra depende tanto del contexto previo como del posterior. Las arquitecturas \gls{Transformer} explotan esta estructura mediante mecanismos de atención que permiten que cada posición interactúe con todas las demás, respetando la invarianza permutacional salvo por los positional encodings que preservan el orden secuencial.
\end{example}

\begin{remark}[Invarianza por permutación]\label{rem:invarianza_permutacion}
Los ejemplos anteriores revelan una distinción fundamental. En algunos dominios (audio, texto, imagen) existe un orden canónico o parcial de las posiciones, de modo que la representación matricial $F \in \mathbb{R}^{n \times d}$ es esencialmente única. En otros (grafos, nubes de puntos), no existe tal orden: si $P \in \{0,1\}^{n \times n}$ es una matriz de permutación, entonces $F$ y $PF$ representan la misma señal sobre el mismo dominio, simplemente con las posiciones enumeradas en distinto orden. Formalmente, la señal $f: \Omega \to \mathcal{C}$ es intrínseca al dominio, pero su representación matricial depende de una elección arbitraria de enumeración $\Omega = \{p_1, \ldots, p_n\}$.

Esta observación tiene consecuencias profundas para el diseño de arquitecturas neuronales: una red que procese datos sin orden canónico debe producir resultados que sean invariantes (para tareas de clasificación) o equivariantes (para tareas de predicción por nodo) bajo permutaciones de la entrada. Esta es una de las ideas centrales del \gls{GDL} que formalizaremos en el \autoref{ch:gdl}.
\end{remark}

La \autoref{tab:resumen_ejemplos_senales} resume las características de los ejemplos anteriores, destacando la diversidad de dominios y la presencia o ausencia de orden canónico.

\begin{table}[htbp]
\centering
\caption[Resumen de ejemplos de señales sobre dominios estructurados]{Resumen comparativo de los ejemplos de señales sobre dominios estructurados. La columna ``¿Orden canónico?'' indica si las posiciones del dominio tienen un ordenamiento natural, lo cual determina si la representación matricial es única salvo permutaciones.}
\label{tab:resumen_ejemplos_senales}
\begin{tabular}{lllcl}
\toprule
\textbf{Ejemplo} & $\Omega$ & \textbf{Tipo de} $\mathcal{S}$ & $\dim(\mathcal{C})$ & \textbf{¿Orden?} \\
\midrule
Audio & $[0,T]$ & Orden total & 1 & Sí \\
Imagen & Grilla 2D & Grilla rectangular & 3 (RGB) & Sí (parcial) \\
Red social & Vértices & Grafo & $d$ & No \\
Molécula & Átomos & Grafo $+ E(3)$ & $d$ & No \\
Nube 3D & Puntos $\subset \mathbb{R}^3$ & Métrica euclidiana & $d$ & No \\
Superficie & $\mathcal{M}$ & Riemanniana & $d$ & No \\
Texto & $\{1,\ldots,L\}$ & Secuencial & $d$ & Sí \\
\bottomrule
\end{tabular}
\end{table}

\subsection{Geometría y atributos}\label{sec:geometria_atributos}

El marco anterior separa dominio, estructura y señal en tres componentes distintas. Sin embargo, puede ser útil conceptualmente unificar el dominio y su estructura en una sola entidad: la \textit{geometría} del problema.

\begin{definition}[Geometría del dominio]\label{def:geometria_dominio}
La \textit{geometría} de un problema de datos es el par $\mathcal{D} = (\Omega, \mathcal{S})$ que combina el dominio $\Omega$ con su estructura $\mathcal{S}$ en una entidad única que caracteriza completamente las relaciones espaciales, topológicas o algebraicas entre posiciones.
\end{definition}

\begin{definition}[Atributos]\label{def:atributos}
Los \textit{atributos} (o características, features) es la imagen, la señal $f: \Omega \to \mathcal{C}$ que asigna valores observados a cada posición del dominio.
\end{definition}

Los dos marcos conceptuales son matemáticamente equivalentes\label{rem:equivalencia_marcos}. La equivalencia es directa: $(\Omega, \mathcal{S}, f) \leftrightarrow (\mathcal{D}, f)$ donde $\mathcal{D} = (\Omega, \mathcal{S})$.

El marco geometría/atributos ofrece una ventaja conceptual importante al distinguir entre estructura intrínseca y observaciones variables:

\begin{itemize}
	\item La \textbf{geometría} $\mathcal{D} = (\Omega, \mathcal{S})$ es típicamente \textit{fija} o \textit{intrínseca al dato individual}: define la ``forma'' del espacio de datos.

	\item Los \textbf{atributos} $f$ son \textit{variables}: diferentes instancias de datos comparten la misma geometría pero tienen atributos distintos.
\end{itemize}

Esta distinción es fundamental para el diseño de arquitecturas neuronales: las operaciones deben respetar la geometría (que es intrínseca) mientras transforman los atributos (que son variables). Ignorar la geometría, como hace una neurona simple con producto interno global, desperdicia información estructural valiosa.

\subsection{Procesamiento neuronal de señales}\label{sec:neurona_senales}

Habiendo establecido el marco de señales sobre dominios estructurados, podemos ahora reinterpretar la neurona artificial definida anteriormente y analizar sus limitaciones desde una perspectiva geométrica.

Una neurona artificial $f_\theta: \mathbb{R}^n \to \mathbb{R}$ puede reinterpretarse como un operador que procesa señales sobre un dominio discreto:

\begin{itemize}
	\item \textbf{Dominio de entrada}: $\Omega = \{1, 2, \ldots, n\}$ (conjunto finito de índices sin significado geométrico particular)
	\item \textbf{Estructura}: Ninguna significativa (conjunto sin relaciones geométricas, topológicas o algebraicas)
	\item \textbf{Señal de entrada}: $x: \Omega \to \mathbb{R}$ donde $x(i) = x_i$ representa el valor en el índice $i$
	\item \textbf{Procesamiento}: $y = \sigma(\langle \mathbf{w}, \mathbf{x} \rangle + b) = \sigma\left(\sum_{i=1}^n w_i \cdot x(i) + b\right)$
\end{itemize}

La transformación lineal $\langle \mathbf{w}, \mathbf{x} \rangle = \sum_{i=1}^n w_i x_i$ combina todas las componentes de forma global mediante una suma ponderada. No existe noción de:

\begin{itemize}
	\item \textbf{Localidad}: Todas las posiciones $i \in \Omega$ contribuyen por igual a la salida; no hay concepto de ``procesamiento local'' basado en vecindad

	\item \textbf{Vecindad estructurada}: No se distingue entre índices ``cercanos'' y ``lejanos''; la conexión $w_1 x_1$ tiene el mismo carácter estructural que $w_{100} x_{100}$

	\item \textbf{Invarianza geométrica}: Permutar los índices mediante $\pi: \{1,\ldots,n\} \to \{1,\ldots,n\}$ produce una señal diferente $x_\pi$ donde $x_\pi(i) = x(\pi(i))$, y típicamente $f_\theta(x_\pi) \neq f_\theta(x)$, aunque semánticamente la permutación puede ser irrelevante (por ejemplo, en grafos)
\end{itemize}

La neurona trata el dominio $\Omega = \{1,\ldots,n\}$ como un conjunto sin estructura, ignorando cualquier geometría que pueda existir implícitamente en la naturaleza de los datos.

Las \glspl{ANN} se centran en transformaciones de la señal o atributos entre las capas, dejando la geometría en un segundo plano, o utilizándola únicamente para optimizar los cálculos computacionales. En el \gls{GDL} veremos que la geometría de los datos toma mayor relevancia y es uno de los pilares fundamentales. Los ejemplos concretos de diferentes tipos de geometrías (euclidiana, de grafos, variedades) y su relación con las arquitecturas neuronales se desarrollan en profundidad en el capítulo de Aprendizaje Profundo Geométrico (\autoref{ch:gdl}), donde se presenta el marco matemático completo.

\section{Capas de una red neuronal}

En la \autoref{def:neurona_artificial}, hemos presentado una \textit{neurona artificial} como una función con una única salida escalar:
\[
f(x) = \sigma(w^T x + b): \mathbb{R}^n \to \mathbb{R}
\]
Esta elección no es arbitraria, sino que responde a varios criterios:

\begin{enumerate}
    \item \textbf{Analogía biológica}: Las neuronas biológicas emiten una única señal de salida (frecuencia de disparo), que se propaga a través del axón hacia otras neuronas. La definición escalar captura esta característica fundamental.

    \item \textbf{Simplicidad matemática}: Una función $\mathbb{R}^n \to \mathbb{R}$ es conceptualmente más simple que una función $\mathbb{R}^n \to \mathbb{R}^m$, facilitando el análisis inicial de las propiedades de una neurona individual (separación por hiperplano, activación no lineal, etc.).
\end{enumerate}

Matemáticamente, podríamos haber definido una neurona como una función vectorial $f: \mathbb{R}^n \to \mathbb{R}^m$ con múltiples salidas. Sin embargo, esta generalización no aporta poder expresivo adicional: una neurona con $m$ salidas es funcionalmente equivalente a $m$ neuronas escalares operando en paralelo, cada una con sus propios pesos y sesgo. Por tanto, adoptamos la definición escalar como primitiva, y construiremos arquitecturas más complejas mediante la composición y paralelización de estas unidades básicas.

Para obtener múltiples valores de salida, utilizamos varias neuronas que comparten el mismo dominio de entrada. Cada neurona transforma independientemente la señal de entrada en un único valor escalar, y el conjunto de las $m$ salidas conforma una nueva señal que, en general, reside sobre un dominio distinto al de entrada que puede tener una estructura diferente. Esta organización nos lleva a formalizar la noción de \textit{capa neuronal}.

\begin{definition}[Capa neuronal]
\label{def:capa_neuronal}
Una \textit{capa neuronal} es un conjunto de $m$ neuronas que comparten el mismo dominio de entrada y operan en paralelo sobre la señal de entrada. Dada una señal de entrada $\mathbf{x} \in \mathbb{R}^n$ (donde $n$ es la dimensión del dominio de entrada), la capa produce una señal de salida $\mathbf{y} \in \mathbb{R}^m$ mediante:
\[
L(\mathbf{x}) = \sigma(\mathbf{W}\mathbf{x} + \mathbf{b})
\]
donde:
\begin{itemize}
    \item $\mathbf{W} \in \mathbb{R}^{m \times n}$ es la \textit{matriz de pesos}, cuya fila $i$-ésima contiene los pesos $w_i^T$ de la neurona $i$
    \item $\mathbf{b} \in \mathbb{R}^m$ es el \textit{vector de sesgos}
    \item $\sigma$ es la función de activación, aplicada componente a componente
\end{itemize}
Cada neurona $i$ computa $y_i = \sigma(w_i^T \mathbf{x} + b_i)$, transformando los atributos de la señal de entrada en un único atributo de salida.
\end{definition}

\paragraph{Conectividad y matrices dispersas.}

 La representación matricial $\mathbf{W}\mathbf{x} + \mathbf{b}$ es completamente general. Si una neurona $i$ solo procesa un subconjunto de las posiciones del dominio de entrada, esto se expresa mediante coeficientes nulos: $W_{ij} = 0$ cuando la neurona $i$ no recibe señal de la posición $j$. Aunque en implementaciones prácticas se utilizan optimizaciones para matrices dispersas, matemáticamente toda capa neuronal admite esta representación unificada. Esta observación será relevante cuando estudiemos arquitecturas con conectividad estructurada, como las redes convolucionales.

\paragraph{Pesos compartidos.}

La \autoref{def:capa_neuronal} asume que cada neurona tiene sus propios parámetros independientes $(w_i, b_i)$. Sin embargo, en muchas arquitecturas los pesos se \textit{comparten} entre neuronas: múltiples filas de la matriz $\mathbf{W}$ contienen los mismos valores. Matemáticamente, esto impone restricciones de la forma $W_{i,:} = W_{j,:}$ para ciertos pares $(i,j)$, reduciendo drásticamente el número de parámetros libres.

\subsection{Composición de capas}

Una neurona por sí misma no es capaz de resolver problemas complejos, pero al combinar muchas de ellas en una red se pueden aprender funciones más complejas. Como hemos visto en la \autoref{def:neurona_artificial}, una neurona es función de la forma $y = \sigma\left(\sum_{i=1}^n w_i x_i + b\right)$, donde $\sigma$ es una función de activación no lineal. Geométricamente, la parte lineal $\sum_{i=1}^n w_i x_i + b$ define un hiperplano en $\mathbb{R}^n$, y la función de activación introduce una no linealidad que transforma este hiperplano. Una neurona individual, por tanto, solo puede separar el espacio de entrada mediante un único hiperplano transformado no linealmente. Esto limita drásticamente su \textit{capacidad expresiva}: no puede aproximar funciones con múltiples regiones de decisión complejas, patrones jerárquicos, o estructuras altamente no lineales. Por ejemplo, una neurona individual no puede resolver el problema clásico del XOR, que requiere una partición no linealmente separable del espacio.

Cuando conectamos múltiples neuronas organizadas en capas, obtenemos una \textit{\gls{composicion_funcional}} de transformaciones no lineales. Denotaremos por $L_i$ la $i$-ésima capa de la red, donde cada capa $L_i: \mathbb{R}^{n_{i-1}} \to \mathbb{R}^{n_i}$ transforma una señal de dimensión $n_{i-1}$ en una señal de dimensión $n_i$. Una red neuronal de $k$ capas computa la composición:
\[
\mathcal{N} = L_k \circ L_{k-1} \circ \cdots \circ L_2 \circ L_1: \mathbb{R}^{n_0} \to \mathbb{R}^{n_k}
\]
donde $n_0$ es la dimensión del dominio de entrada y $n_k$ la del dominio de salida.

Esta composición tiene propiedades matemáticas fundamentales:
\begin{enumerate}
    \item \textbf{Representaciones jerárquicas}: Cada composición $L_i \circ L_{i-1} \circ \cdots \circ L_1$ puede interpretarse como una transformación que extrae características cada vez más abstractas de los datos. Las primeras capas detectan patrones simples (bordes, texturas), mientras que capas posteriores combinan estos patrones en representaciones de mayor nivel semántico (formas, objetos).

    \item \textbf{Particionamiento complejo del espacio}: Múltiples neuronas en una capa generan múltiples hiperplanos, permitiendo particionar $\mathbb{R}^n$ en regiones poliédricas arbitrariamente complejas. La composición de capas refina recursivamente estas particiones, creando fronteras de decisión altamente no lineales.
\end{enumerate}

% ============ Versión TikZ con colores ============
\begin{figure}[h]
\centering
\begin{tikzpicture}[
    x=2.1cm, y=0.85cm,
    neuron/.style={circle, minimum size=5.5mm, line width=0.6pt},
    input/.style={neuron, draw=gdlblue!70, fill=gdlblue!15},
    hidden1/.style={neuron, draw=gdlgreen!60!black, fill=gdlgreen!15},
    hidden2/.style={neuron, draw=gdlgreen!70, fill=gdlgreen!18},
    hidden3/.style={neuron, draw=gdlpurple!70, fill=gdlpurple!18},
    output/.style={neuron, draw=gdlorange!70, fill=gdlorange!18},
    conn/.style={->, >=stealth, draw=gdlgray!35, line width=0.2pt},
    lbl/.style={font=\small\bfseries}
]

% Capa de entrada (3 neuronas)
\node[input] (I1) at (0, 2) {};
\node[input] (I2) at (0, 1) {};
\node[input] (I3) at (0, 0) {};
\node[left=1mm of I1, font=\small] {$x_1$};
\node[left=1mm of I2, font=\small] {$x_2$};
\node[left=1mm of I3, font=\small] {$x_3$};
\node[lbl] at (0, 3.5) {Entrada};

% Primera capa oculta (4 neuronas)
\node[hidden1] (H1-1) at (1, 2.5) {};
\node[hidden1] (H1-2) at (1, 1.5) {};
\node[hidden1] (H1-3) at (1, 0.5) {};
\node[hidden1] (H1-4) at (1, -0.5) {};
\node[lbl] at (1, 3.5) {Oculta 1};

% Segunda capa oculta (4 neuronas)
\node[hidden2] (H2-1) at (2, 2.5) {};
\node[hidden2] (H2-2) at (2, 1.5) {};
\node[hidden2] (H2-3) at (2, 0.5) {};
\node[hidden2] (H2-4) at (2, -0.5) {};
\node[lbl] at (2, 3.5) {Oculta 2};

% Tercera capa oculta (4 neuronas)
\node[hidden3] (H3-1) at (3, 2.5) {};
\node[hidden3] (H3-2) at (3, 1.5) {};
\node[hidden3] (H3-3) at (3, 0.5) {};
\node[hidden3] (H3-4) at (3, -0.5) {};
\node[lbl] at (3, 3.5) {Oculta 3};

% Capa de salida (2 neuronas)
\node[output] (O1) at (4, 1.5) {};
\node[output] (O2) at (4, 0.5) {};
\node[right=1mm of O1, font=\small] {$y_1$};
\node[right=1mm of O2, font=\small] {$y_2$};
\node[lbl] at (4, 3.5) {Salida};

% Conexiones: entrada -> oculta 1
\foreach \i in {1,2,3}
    \foreach \j in {1,...,4}
        \draw[conn] (I\i) -- (H1-\j);

% Conexiones: oculta 1 -> oculta 2
\foreach \i in {1,...,4}
    \foreach \j in {1,...,4}
        \draw[conn] (H1-\i) -- (H2-\j);

% Conexiones: oculta 2 -> oculta 3
\foreach \i in {1,...,4}
    \foreach \j in {1,...,4}
        \draw[conn] (H2-\i) -- (H3-\j);

% Conexiones: oculta 3 -> salida
\foreach \i in {1,...,4}
    \foreach \j in {1,2}
        \draw[conn] (H3-\i) -- (O\j);

% Peso ejemplo (conexión destacada)
\draw[->, >=stealth, draw=gdlblue!60, line width=0.5pt] (I1) -- node[above, font=\scriptsize, sloped, pos=0.4] {$w_{11}^{(1)}$} (H1-1);

% Corchete para capas ocultas
\draw[decorate, decoration={brace, amplitude=4pt, raise=2pt}, gdlgray!60, thick]
    (1, 3.7) -- (3, 3.7) node[midway, above=5pt, font=\small\itshape] {Capas ocultas};

\end{tikzpicture}
\caption[Estructura de una red neuronal en capas]{Estructura de una red neuronal feedforward organizada en capas. Los nodos representan neuronas individuales con colores diferenciados por tipo de capa: entrada (azul), ocultas (verde/teal/violeta), y salida (naranja). Las flechas representan conexiones ponderadas $w_{ij}^{(\ell)}$.}
\label{fig:neural_network}
\end{figure}

En una \gls{ANN} tenemos una primera capa de la red denominada capa de entrada que recibe los datos crudos, es decir, la señal del dominio del problema; y la última capa es la capa de salida que produce la respuesta final de la red. Las capas intermedias se conocen como capas ocultas, ya que no se ven directamente en la entrada o la salida de la red neuronal.

Veamos por qué es importante tener más de una capa y más de una neurona en alguna capa con más detalle en el problema XOR.

\subsubsection{El problema XOR}

El problema XOR consiste en computar la función lógica:
\[
\text{XOR}(x_1, x_2) = \begin{cases}
1 & \text{si exactamente una entrada es 1} \\
0 & \text{en otro caso}
\end{cases}
\]
Este problema es \textit{no linealmente separable}: no existe un único hiperplano que separe las clases $(0,0), (1,1)$ (salida 0) de $(0,1), (1,0)$ (salida 1) en $\mathbb{R}^2$.

La \autoref{fig:xor_geometrico} ilustra geométricamente cómo dos neuronas en la capa oculta, cada una definiendo un hiperplano, permiten resolver el problema XOR.

\begin{figure}[h]
\centering
\begin{tikzpicture}[scale=2.5]

% Ejes
\draw[->] (-0.15, 0) -- (1.25, 0) node[right, font=\small] {$x_1$};
\draw[->] (0, -0.15) -- (0, 1.25) node[above, font=\small] {$x_2$};

% Puntos clase 0 (azul) - XOR=0
\fill[gdlblue!70] (0, 0) circle (0.06) node[below left, font=\small] {$(0,0)$};
\fill[gdlblue!70] (1, 1) circle (0.06) node[above right, font=\small] {$(1,1)$};

% Puntos clase 1 (rojo) - XOR=1
\fill[gdlred!70] (0, 1) circle (0.06) node[above left, font=\small] {$(0,1)$};
\fill[gdlred!70] (1, 0) circle (0.06) node[below right, font=\small] {$(1,0)$};

% Hiperplano 1 (neurona 1): x1 + x2 = 0.5
\draw[gdlgreen!60!black, thick] (-0.1, 0.6) -- (0.6, -0.1);
\node[gdlgreen!60!black, font=\small, rotate=-45] at (0.05, 0.15) {$h_1$};

% Hiperplano 2 (neurona 2): x1 + x2 = 1.5
\draw[gdlgreen!70!black, thick] (0.4, 1.1) -- (1.1, 0.4);
\node[gdlgreen!70!black, font=\small, rotate=-45] at (0.95, 0.85) {$h_2$};

% Leyenda
\node[anchor=west, font=\small] at (1.4, 0.9) {$h_1$: neurona 1};
\node[anchor=west, font=\small] at (1.4, 0.7) {$h_2$: neurona 2};
\node[anchor=west, font=\small, gdlblue!70] at (1.4, 0.4) {XOR = 0};
\node[anchor=west, font=\small, gdlred!70] at (1.4, 0.2) {XOR = 1};

\end{tikzpicture}
\caption[Solución geométrica del problema XOR]{Solución geométrica del problema XOR mediante dos neuronas. Cada neurona de la capa oculta define un hiperplano ($h_1$ y $h_2$) que divide el espacio. La región entre ambos hiperplanos corresponde a XOR=1 (puntos rojos), mientras que las regiones exteriores corresponden a XOR=0 (puntos azules).}
\label{fig:xor_geometrico}
\end{figure}

¿Qué arquitecturas pueden resolver XOR?
\begin{itemize}
    \item \textbf{1 capa con 1 neurona}: \textit{Imposible}. Una neurona computa $\sigma(w_1 x_1 + w_2 x_2 + b)$, que define un único hiperplano. XOR requiere una frontera de decisión no convexa.

    \item \textbf{Muchas capas con 1 neurona cada una}: También \textit{imposible}. Si cada capa aplica una función escalar $f_i: \mathbb{R} \to \mathbb{R}$, la composición $f_k \circ \cdots \circ f_1: \mathbb{R}^2 \to \mathbb{R}$ colapsa toda la información bidimensional a un escalar inmediatamente, perdiendo la capacidad de distinguir las cuatro regiones del espacio de entrada.

    \item \textbf{2 capas con $\{2, 1\}$ neuronas}: \textit{Suficiente}. La capa oculta con 2 neuronas crea dos hiperplanos que particionan $\mathbb{R}^2$ en regiones, y la neurona de salida combina linealmente estas particiones para generar la frontera XOR.
\end{itemize}

Este contraejemplo revela que \textit{ni solo anchura ni solo profundidad bastan}: necesitamos $m \geq 2$ neuronas en al menos una capa \textit{y} al menos $k \geq 2$ capas. Las arquitecturas extremas (profundas-estrechas o shallow-anchas con $m=1$) fallan en problemas fundamentales.


\subsection{Anchura vs profundidad}\label{subsec:anchura_profundidad}

El análisis del problema XOR nos lleva a formalizar dos conceptos fundamentales que caracterizan la arquitectura de una red neuronal: la \textit{anchura} y la \textit{profundidad}. Estos parámetros determinan la capacidad expresiva de la red y tienen profundas implicaciones tanto teóricas como prácticas.

\begin{definition}[Anchura de una capa]\label{def:anchura}
La \textit{anchura} de una capa $L_i: \mathbb{R}^{n_{i-1}} \to \mathbb{R}^{n_i}$ es la dimensión $n_i$ de su espacio de salida, es decir, el número de neuronas que contiene.
\end{definition}

\begin{definition}[Profundidad de una red]\label{def:profundidad}
La \textit{profundidad} de una red neuronal $\mathcal{N} = L_k \circ \cdots \circ L_1$ es el número $k$ de capas en la composición.
\end{definition}

Estas definiciones formales nos permiten analizar rigurosamente el balance entre anchura y profundidad.

La anchura proporciona \textit{diversidad de características}: cada neurona captura un aspecto diferente de la señal de entrada, y la señal de salida $h \in \mathbb{R}^m$ codifica simultáneamente múltiples propiedades geométricas del dato.

La profundidad proporciona \textit{abstracción jerárquica}: en lugar de intentar resolver el problema completo de una vez, cada capa realiza una transformación intermedia, delegando el refinamiento final a las capas posteriores. Las capas implementan una jerarquía de representaciones donde cada nivel extrae características de complejidad creciente:
\[
x \xrightarrow{L_1} h^{(1)} \xrightarrow{L_2} h^{(2)} \xrightarrow{L_3} \cdots \xrightarrow{L_k} y
\]
Empíricamente, se observa que:
\begin{itemize}
    \item \textbf{Capas tempranas} ($L_1, L_2$): Detectan características de bajo nivel (bordes, texturas, colores básicos en imágenes; fonemas en audio).
    \item \textbf{Capas intermedias} ($L_3, \ldots, L_{k-2}$): Combinan características de bajo nivel en patrones más complejos (formas geométricas, motivos texturales).
    \item \textbf{Capas finales} ($L_{k-1}, L_k$): Representan conceptos de alto nivel específicos de la tarea (objetos, escenas; palabras, conceptos semánticos).
\end{itemize}
Esta jerarquía refleja la estructura composicional de muchos problemas del mundo real y permite que la red aprenda representaciones cada vez más abstractas y semánticamente significativas.

Como veremos en el Teorema de aproximación universal (\autoref{sec:aproximacion_universal}), las redes neuronales con una sola capa oculta suficientemente ancha pueden aproximar cualquier función continua en un compacto con precisión arbitraria. Sin embargo, esta garantía existencial no especifica la eficiencia de dicha aproximación en términos del número de neuronas necesarias.

\paragraph{Teoría de aproximación y profundidad.}

Resultados teóricos recientes demuestran que la profundidad (número de capas) proporciona ventajas exponenciales en capacidad expresiva. Para ciertas clases de funciones, una red profunda de ancho moderado puede aproximar funciones que requieren un número exponencialmente mayor de neuronas si se limita a una sola capa~\citep{telgarsky2016benefitsdepthneuralnetworks}. Matemáticamente, esto se debe a que la composición de funciones permite construir funciones con complejidad combinatoria que crece exponencialmente con la profundidad.

Hemos visto que, aunque una sola capa oculta basta teóricamente para aproximar funciones continuas, las arquitecturas profundas ofrecen ventajas exponenciales en eficiencia. En la práctica, se prefieren redes profundas con diferentes tipos de capas especializadas, cada una adaptada a aspectos específicos del problema, para optimizar tanto la capacidad expresiva como el costo computacional. Esta especialización refleja una distinción fundamental: mientras algunas capas transforman únicamente los valores de la señal preservando su estructura espacial, otras modifican la organización geométrica misma del dominio. Esta taxonomía funcional nos permitirá comprender mejor la diferencia entre las arquitecturas tradicionales de deep learning y el enfoque de \gls{GDL}.

\subsection{Tipos de capas}\label{subsec:tipos_de_capas}

Las capas neuronales pueden transformar la señal, el dominio, o ambos simultáneamente. Esta distinción es fundamental para entender la diferencia filosófica y matemática entre las \glspl{ANN} tradicionales y \gls{GDL}: mientras que las primeras típicamente ignoran o destruyen la geometría original del dato, las segundas la preservan y explotan sistemáticamente.

Formalmente, clasificamos las capas en tres categorías según qué dimensión transforman:

\begin{definition}[Transformación de señal pura]\label{def:signal_transform}
	Una capa $L$ es una \textit{transformación de señal pura} si preserva la estructura del dominio de entrada. Formalmente, si la entrada es una señal $f: \Omega \to \mathbb{R}^{c_{\text{in}}}$, la salida es $g: \Omega \to \mathbb{R}^{c_{\text{out}}}$ donde el dominio $\Omega$ permanece invariante.

	\textbf{Características:}
	\begin{itemize}
		\item Operaciones punto a punto (element-wise): $g(x) = \phi(f(x))$ para alguna función $\phi: \mathbb{R}^{c_{\text{in}}} \to \mathbb{R}^{c_{\text{out}}}$
		\item Preservan la forma (shape) del tensor de entrada
		\item No requieren información sobre vecindades, conectividad o estructura geométrica
	\end{itemize}
\end{definition}

\begin{definition}[Transformación geométrica pura]\label{def:geometry_transform}
	Una capa $L$ es una \textit{transformación geométrica pura} si modifica la estructura geométrica del dominio sin procesar los valores de la señal de forma no trivial. Formalmente, si $f: (\Omega, \mathcal{G}) \to \mathbb{R}^c$ es una señal sobre un dominio $\Omega$ con estructura geométrica $\mathcal{G}$ (métrica, conexión, topología), entonces $L(f) = g: (\Omega', \mathcal{G}') \to \mathbb{R}^c$ donde la geometría $\mathcal{G}' \neq \mathcal{G}$, pero los valores de la señal se preservan o reindexan según una correspondencia $\pi: \Omega' \to \Omega$.
\end{definition}

En \textbf{redes neuronales tradicionales}, las transformaciones geométricas puras se reducen típicamente a reindexaciones (reshape, transpose, flatten), ya que el dominio subyacente es un espacio euclidiano discretizado sin estructura geométrica rica. Sin embargo, en \gls{GDL}, estas transformaciones incluyen operaciones más sofisticadas: cambios de sistema de coordenadas, proyecciones entre variedades, modificaciones de la métrica o conexión, y transformaciones conformes o isométricas. La distinción es fundamental: mientras que en \gls{ANN} tradicionales la geometría es ``plana'' y las transformaciones puramente estructurales son triviales, en \gls{GDL} la geometría es intrínseca al problema y las transformaciones geométricas pueden ser altamente no triviales.

\begin{definition}[Transformación acoplada señal-geometría]\label{def:coupled_transform}
	Una capa $L$ es una \textit{transformación acoplada} si modifica simultánea e intrínsecamente tanto la señal como el dominio. Formalmente, $L: (f: \Omega \to \mathbb{R}^{c_{\text{in}}}) \mapsto (g: \Omega' \to \mathbb{R}^{c_{\text{out}}})$ donde típicamente $\Omega' \neq \Omega$ y $g$ no puede expresarse como reindexación de $f$.
\end{definition}

Las \textbf{redes neuronales tradicionales} (perceptrones multicapa, redes densas) se basan principalmente en capas que \textit{destruyen} o ignoran completamente la geometría original del dato. Es muy típico aplanar la geometría y transformarla en un vector, perdiendo toda información sobre vecindades espaciales y estructura local, o no utilizar las propiedades intrínsecas de la geometría como simetrías para mejorar las redes de neuronas. En contraste, como veremos en el \autoref{ch:gdl}, el \gls{GDL} preserva y explota sistemáticamente la geometría subyacente.

\subsection{Capas tradicionales}

A modo de introducción vamos a presentar aquí brevemente algunas de las capas utilizadas tradicionalmente en Deep Learning.

\subsubsection{Funciones de activación}\label{subsubsec:funciones_activacion_capas}

Las funciones de activación introducen no linealidades esenciales en las \gls{ANN}, permitiendo que aproximen funciones complejas. Sin estas no linealidades, cualquier red neuronal se reduciría a una simple transformación lineal, limitando drásticamente su capacidad expresiva. Anteriormente hemos visto que la función de activación forma parte intrínseca de la neurona, pero por razones prácticas se suele implementar como una capa independiente que transforma la salida de cada neurona componente a componente. Consideramos este tipo de capas dentro de las que únicamente realizan transformaciones de la señal.

A continuación presentamos las funciones de activación más utilizadas, analizando en cada caso las propiedades deseables establecidas en la \autoref{prop:acotacion_activacion} (véase también la \autoref{tab:funciones_activacion} para una comparación resumida):

\paragraph{Sigmoide}

La función sigmoide es una de las funciones de activación clásicas. Su rango de valores es $(0, 1)$, útil para interpretaciones probabilísticas:
\[
\sigma(x) = \frac{1}{1 + e^{-x}}
\]

Es acotada (\autoref{prop:acotacion_activacion}), monótona (\autoref{prop:monotonia_activacion}) y Lipschitz continua con constante $L = 1/4$ (\autoref{prop:lipschitz_continuidad}). Sin embargo, no está centrada en cero (\autoref{prop:zero_centered}), ya que $\sigma(z) \in (0,1)$ y $\sigma(0) = 1/2 \neq 0$, lo que puede provocar el fenómeno de zigzagging en el descenso de gradiente. Además, sus asíntotas horizontales ($\lim_{z\to\pm\infty}\sigma(z) \in \{0,1\}$) implican saturación del gradiente en ambos extremos (\autoref{prop:gradiente_no_saturante}).

\paragraph{Tangente Hiperbólica (tanh)}

La función tanh tiene un rango $(-1, 1)$:
\[
\tanh(x) = \frac{e^x - e^{-x}}{e^x + e^{-x}} = \frac{e^{2x} - 1}{e^{2x} + 1}
\]

A diferencia de la sigmoide, la tangente hiperbólica está centrada en cero (\autoref{prop:zero_centered}): es función impar ($\tanh(-z) = -\tanh(z)$) y $\tanh(0) = 0$. Comparte con la sigmoide la acotación (\autoref{prop:acotacion_activacion}), la monotonía (\autoref{prop:monotonia_activacion}) y la continuidad de Lipschitz con $L = 1$ (\autoref{prop:lipschitz_continuidad}).

Sin embargo, tanto la sigmoide como la tangente hiperbólica son funciones acotadas con asíntotas horizontales, lo que las hace saturantes (\autoref{prop:gradiente_no_saturante}). Como se analiza en la \autoref{subsubsec:backpropagation}, esto provoca el problema del \textit{desvanecimiento del gradiente}: en redes profundas, el gradiente se atenúa exponencialmente en las primeras capas, que apenas ``aprenden''.

\paragraph{ReLU (Rectified Linear Unit)}

La función \gls{ReLU} es una de las funciones de activación más utilizadas en capas intermedias:
\[
\text{ReLU}(x) = \max(0, x) = \begin{cases}
x & \text{si } x > 0 \\
0 & \text{si } x \leq 0
\end{cases}
\]

\gls{ReLU} resuelve el problema del desvanecimiento del gradiente: para $z > 0$, $\text{ReLU}'(z) = 1$, satisfaciendo la \autoref{prop:gradiente_no_saturante} en la semirrecta positiva. Tiene crecimiento lineal (\autoref{prop:acotacion_activacion}, condición 2), es monótona (\autoref{prop:monotonia_activacion}) y Lipschitz continua con $L = 1$ (\autoref{prop:lipschitz_continuidad}). Además, al requerir solo una comparación y una selección, es computacionalmente muy eficiente (\autoref{prop:eficiencia_computacional}). No obstante, no está centrada en cero (\autoref{prop:zero_centered}), y su gradiente nulo para $z \leq 0$ puede causar el problema de ``neuronas muertas'': neuronas cuya entrada permanece negativa dejan de contribuir al aprendizaje de forma permanente.

\paragraph{LeakyReLU}

Una variante de \gls{ReLU} que permite un pequeño gradiente para valores negativos:
\[
\text{LeakyReLU}(x) = \begin{cases}
x & \text{si } x > 0 \\
\alpha x & \text{si } x \leq 0
\end{cases}
\]

donde $\alpha > 0$ es un parámetro pequeño (típicamente $\alpha = 0.01$). LeakyReLU satisface plenamente la \autoref{prop:gradiente_no_saturante}: $|\text{LeakyReLU}'(z)| \geq \alpha > 0$ para todo $z \neq 0$, eliminando tanto la saturación como el problema de neuronas muertas de ReLU. Es monótona (\autoref{prop:monotonia_activacion}), tiene crecimiento lineal (\autoref{prop:acotacion_activacion}), es Lipschitz con $L = 1$ (\autoref{prop:lipschitz_continuidad}) y comparte la eficiencia computacional de ReLU (\autoref{prop:eficiencia_computacional}).

\paragraph{Softmax}

Esta función se utiliza normalmente para obtener valores probabilísticos en la última capa de las redes de neuronas, y es fundamental para los mecanismos de atención que estudiaremos en el \autoref{ch:attention}.

\begin{definition}[Función Softmax]\label{def:softmax}
La \textit{función softmax} $\sigma: \mathbb{R}^n \to (0,1)^n$ transforma un vector de valores reales en una distribución de probabilidad. Para un vector $\mathbf{z} = (z_1, \ldots, z_n) \in \mathbb{R}^n$:
\begin{equation}\label{eq:softmax}
\sigma(\mathbf{z})_i = \frac{e^{z_i}}{\sum_{j=1}^{n} e^{z_j}}, \quad i = 1, \ldots, n
\end{equation}
\end{definition}

\begin{proposition}[Propiedades de la función Softmax]\label{prop:softmax_propiedades}
La función softmax satisface las siguientes propiedades:
\begin{enumerate}
	\item \textbf{Normalización}: $\sum_{i=1}^n \sigma(\mathbf{z})_i = 1$ para todo $\mathbf{z} \in \mathbb{R}^n$.
	\item \textbf{Positividad}: $\sigma(\mathbf{z})_i > 0$ para todo $i$ y todo $\mathbf{z} \in \mathbb{R}^n$.
	\item \textbf{Invarianza a traslación}: $\sigma(\mathbf{z} + c\mathbf{1}) = \sigma(\mathbf{z})$ para todo $c \in \mathbb{R}$.
	\item \textbf{Diferenciabilidad}: El jacobiano tiene la forma $\frac{\partial \sigma_i}{\partial z_j} = \sigma_i(\delta_{ij} - \sigma_j)$.
	\item \textbf{Equivarianza a permutaciones}: Para cualquier matriz de permutación $P$, $\sigma(P\mathbf{z}) = P\sigma(\mathbf{z})$.
\end{enumerate}
\end{proposition}

\begin{proof}
Las propiedades (1) y (2) se verifican directamente de la definición. La propiedad (3) resulta de que $e^{z_i + c}/\sum_j e^{z_j + c} = e^{z_i}e^c/(e^c\sum_j e^{z_j}) = e^{z_i}/\sum_j e^{z_j}$. Para (4), aplicando la regla del cociente:
\[
\frac{\partial \sigma_i}{\partial z_j} = \frac{\delta_{ij}e^{z_i}\sum_k e^{z_k} - e^{z_i}e^{z_j}}{(\sum_k e^{z_k})^2} = \sigma_i\delta_{ij} - \sigma_i\sigma_j = \sigma_i(\delta_{ij} - \sigma_j)
\]
La propiedad (5) se sigue de que la suma en el denominador es invariante bajo permutación de índices.
\end{proof}

\subsubsection{Normalización}

Las técnicas de normalización han revolucionado el entrenamiento de redes neuronales profundas al abordar problemas fundamentales como el desplazamiento de covariables internas (\textit{internal covariate shift}) y la inestabilidad de gradientes. Estas técnicas normalizan las distribuciones de activaciones en diferentes puntos de la red. Es un tipo de capa de transformación de la señal pura, en particular, la normalización aplica estadísticas sobre la señal (media, varianza) sin modificar la estructura del dominio.

\begin{definition}[Normalización en redes neuronales]\label{def:normalizacion}
	Una técnica de normalización es una transformación que estandariza las activaciones de una capa aplicando una función de la forma:
	\[
	\text{Norm}(\mathbf{x}) = \gamma \odot \frac{\mathbf{x} - \boldsymbol{\mu}}{\boldsymbol{\sigma}} + \boldsymbol{\beta}
	\]
	donde $\boldsymbol{\mu}$ y $\boldsymbol{\sigma}$ son estadísticas de centralización y escala, y $\gamma$, $\boldsymbol{\beta}$ son parámetros aprendibles que permiten recuperar la transformación identidad si es necesario\footnote{En la práctica, se utiliza $\sqrt{\boldsymbol{\sigma}^2 + \epsilon}$ con $\epsilon \approx 10^{-8}$ en lugar de $\boldsymbol{\sigma}$ para evitar divisiones por cero, como se observa en las especializaciones que siguen.}.
\end{definition}

Hay diferentes capas de normalización que se usan en función de como se quiera agrupar las neuronas de una capa.

\paragraph{\gls{BN}}

Introducida por Ioffe y Szegedy~\citep{ioffe2015batch}, la normalización por lotes calcula estadísticas sobre el mini-batch:

\[
\text{BN}(\mathbf{x}_i) = \gamma \frac{\mathbf{x}_i - \boldsymbol{\mu}_{\mathcal{B}}}{\sqrt{\boldsymbol{\sigma}^2_{\mathcal{B}} + \epsilon}} + \boldsymbol{\beta}
\]

donde para un mini-batch $\mathcal{B} = \{\mathbf{x}_1, \ldots, \mathbf{x}_m\}$:
\begin{align}
\boldsymbol{\mu}_{\mathcal{B}} &= \frac{1}{m} \sum_{i=1}^m \mathbf{x}_i \\
\boldsymbol{\sigma}^2_{\mathcal{B}} &= \frac{1}{m} \sum_{i=1}^m (\mathbf{x}_i - \boldsymbol{\mu}_{\mathcal{B}})^2
\end{align}

\paragraph{Layer Normalization}

Propuesta por Ba, Kiros y Hinton~\citep{ba2016layer}, normaliza a través de las características de cada muestra individual:

\[
\text{LN}(\mathbf{x}) = \gamma \frac{\mathbf{x} - \boldsymbol{\mu}_L}{\sqrt{\boldsymbol{\sigma}^2_L + \epsilon}} + \boldsymbol{\beta}
\]

donde para una muestra $\mathbf{x} \in \mathbb{R}^d$:
\begin{align}
\mu_L &= \frac{1}{d} \sum_{i=1}^d x_i \\
\sigma^2_L &= \frac{1}{d} \sum_{i=1}^d (x_i - \mu_L)^2
\end{align}

\paragraph{Group Normalization}

Introducida por Wu y He~\citep{wu2018group}, divide los canales en grupos y normaliza dentro de cada grupo:

Para un tensor de entrada con forma $(N, C, H, W)$, divide $C$ canales en $G$ grupos de $C/G$ canales cada uno, y aplica normalización dentro de cada grupo.

\[
\text{GN}(\mathbf{x}) = \gamma \frac{\mathbf{x} - \boldsymbol{\mu}_G}{\sqrt{\boldsymbol{\sigma}^2_G + \epsilon}} + \boldsymbol{\beta}
\]

\paragraph{Instance Normalization}

Normaliza cada canal de cada muestra independientemente, común en tareas de transferencia de estilo~\citep{ulyanov2016instance}:

\[
\text{IN}(\mathbf{x}) = \gamma \frac{\mathbf{x} - \boldsymbol{\mu}_{IN}}{\sqrt{\boldsymbol{\sigma}^2_{IN} + \epsilon}} + \boldsymbol{\beta}
\]

donde las estadísticas se calculan por canal y por muestra.

\subsubsection{Dropout}

El dropout~\citep{srivastava2014dropout}, es una técnica de regularización fundamental que aborda el problema del sobreajuste en redes neuronales profundas de manera elegante y computacionalmente eficiente. Es un tipo de capa que la podemos considerar como de señal pura, ya que únicamente aplica un enmascaramiento estocástico punto a punto sin alterar la geometría del dominio.

\begin{definition}[Dropout]\label{def:dropout}
	El dropout es una técnica de regularización que durante el entrenamiento desactiva aleatoriamente neuronas con probabilidad $p_{\text{drop}}$, estableciendo sus salidas a cero. Formalmente, para una capa con entrada $\mathbf{x}$:
	\[
	\mathbf{y} = \frac{1}{1-p_{\text{drop}}} \mathbf{r} \odot \mathbf{x}
	\]
	donde $\mathbf{r} \sim \text{Bernoulli}(1-p_{\text{drop}})$ es un vector de máscaras binarias aleatorias, y $\odot$ denota el producto elemento a elemento.
\end{definition}

El dropout proporciona varios beneficios de regularización:
\begin{itemize}
	\item \textbf{Reduce co-adaptación}: Previene que las neuronas desarrollen dependencias complejas entre sí
	\item \textbf{Mejora generalización}: Fuerza a la red a no depender excesivamente de neuronas específicas
	\item \textbf{Robustez}: Hace que las representaciones sean menos frágiles a perturbaciones
\end{itemize}

\subsubsection{Capas densas}\label{def:capa_densa}

Una \textit{capa densa} (también llamada capa completamente conectada o \textit{fully connected layer}) es precisamente una capa neuronal según la \autoref{def:capa_neuronal}, sin ninguna restricción estructural sobre la matriz de pesos: cada neurona de salida recibe señal de todas las posiciones de entrada, es decir, $W$ puede tener todos sus coeficientes no nulos. Es la arquitectura más básica y constituye la base sobre la cual se construyen arquitecturas más especializadas.

Desde una perspectiva geométrica, una capa densa $L(\mathbf{x}) = \sigma(W\mathbf{x} + \mathbf{b})$ realiza tres operaciones sucesivas:
\begin{enumerate}
	\item \textbf{Transformación lineal}: $W\mathbf{x}$ representa un cambio de base y escalamiento en $\mathbb{R}^n \to \mathbb{R}^m$
	\item \textbf{Traslación}: $+ \mathbf{b}$ desplaza el resultado en $\mathbb{R}^m$
	\item \textbf{Deformación no lineal}: $\sigma(\cdot)$ introduce curvatura en el espacio de características
\end{enumerate}

Al no imponer restricciones de conectividad dispersa ni pesos compartidos, la capa densa destruye completamente la geometría original del dominio de entrada mediante una proyección afín a un espacio de dimensión arbitraria. Esta es una transformación acoplada señal-geometría (\autoref{def:coupled_transform}).

\subsubsection{Embedding}

Las capas de embedding constituyen un componente fundamental para el procesamiento de datos categóricos y discretos en redes neuronales. Permiten la transformación de variables categóricas de alta dimensionalidad y dispersas en representaciones densas de dimensionalidad más baja que capturan relaciones semánticas. Los embeddings mapean tokens discretos a espacios continuos aprendidos, cambiando fundamentalmente tanto el dominio (discreto $\to$ continuo) como la representación de la señal, es decir, son un tipo de capa que transforma de forma acoplada la señal y la geometría.

\begin{definition}[Capa de embedding]\label{def:capa_embedding}
	Una capa de embedding es una transformación que mapea índices discretos a vectores densos. Formalmente, para un vocabulario de tamaño $V_{\text{vocab}}$ y dimensión de embedding $d_{\text{emb}}$:
	\[
	\text{Embedding}: \{1, 2, \ldots, V_{\text{vocab}}\} \to \mathbb{R}^{d_{\text{emb}}}
	\]
	implementada mediante una matriz de lookup $E_{\text{emb}} \in \mathbb{R}^{V_{\text{vocab}} \times d_{\text{emb}}}$, donde cada fila $(E_{\text{emb}})_i$ es el vector de embedding para el token $i$.
\end{definition}

La aplicación más popular de los embeddings es transformar texto en vectores para su posterior procesamiento en la red de neuronas.

Los embeddings representan un caso fundamental de aprendizaje de representaciones geométricas: transforman un espacio discreto y desconectado (categorías) en una variedad continua donde la geometría euclidiana captura relaciones semánticas. Esta transformación es esencial para que las redes neuronales procesen efectivamente datos simbólicos y categóricos.

\subsubsection{Transformaciones de estructura}

Las transformaciones de estructura modifican la geometría y organización de los datos sin cambiar fundamentalmente su contenido informativo. Estas operaciones son cruciales para adaptar los datos a diferentes arquitecturas de red, explotar patrones específicos del dominio y optimizar la computación.

\paragraph{Reshape}

\begin{definition}[Operaciones de reshape]\label{def:reshape}
	Una operación de reshape es una transformación que reorganiza un tensor cambiando sus dimensiones mientras preserva el número total de elementos:
	\[
	\text{reshape}: \mathbb{R}^{d_1 \times d_2 \times \cdots \times d_n} \to \mathbb{R}^{d_1' \times d_2' \times \cdots \times d_m'}
	\]
	donde $\prod_{i=1}^n d_i = \prod_{j=1}^m d_j'$.
\end{definition}

\paragraph{Flatten} es un caso especial de reshape que convierte tensores multidimensionales en vectores unidimensionales:
\[
\text{flatten}: \mathbb{R}^{d_1 \times d_2 \times \cdots \times d_n} \to \mathbb{R}^{\prod_{i=1}^n d_i}
\]

Una aplicación muy común de estas capas es servir de transición entre las capas convolucionales y las capas densas.

\paragraph{Pooling}

El pooling reduce la dimensionalidad espacial combinando información local, proporcionando invarianza traslacional limitada, reduciendo el ruido y el costo computacional.

\begin{definition}[Operación de pooling]\label{def:pooling}
	Una operación de pooling agrega valores en ventanas locales usando una función de agregación $f$:
	\[
	y_{i,j} = f\left(\{x_{si+k, sj+l} : k,l \in W_{\text{pool}}\}\right)
	\]
	donde $s$ es el stride y $W_{\text{pool}}$ define el tamaño de la ventana de pooling.
\end{definition}

\textbf{Max Pooling}: Selecciona el valor máximo en cada ventana:
\[
\text{MaxPool}(X) = \max_{(i,j) \in W_{\text{pool}}} X_{i,j}
\]
\textbf{Average Pooling}: Calcula la media aritmética en cada ventana:
\[
\text{AvgPool}(X) = \frac{1}{|W_{\text{pool}}|} \sum_{(i,j) \in W_{\text{pool}}} X_{i,j}
\]

Estas transformaciones son fundamentales para construir arquitecturas jerárquicas que procesen información a múltiples escalas y resoluciones.

\subsubsection{Convoluciones}

Las capas convolucionales se explican con mayor detalle en el \autoref{ch:cnn}. Se utilizan fundamentalmente para datos en rejilla como imágenes o secuencias como audio. En las \glspl{ANN} tradicionales, las convoluciones se introducen principalmente como optimización computacional que aprovecha la estructura regular de grilla del dominio. Como veremos en el \autoref{ch:cnn}, existen variantes más sofisticadas de redes convolucionales que explotan más profundamente las propiedades geométricas y simetrías del dominio.

\begin{definition}[Convolución discreta]\label{def:convolucion_discreta}
	Para una función de entrada $f: \mathbb{Z}^d \to \mathbb{R}$ y un kernel $g: \mathbb{Z}^d \to \mathbb{R}$, la convolución discreta se define como:
	\[
	(f * g)[n] = \sum_{m \in \mathbb{Z}^d} f[m] g[n - m]
	\]
\end{definition}

\subsubsection{Recurrentes}\label{subsubsec:rnn}

Las redes neuronales recurrentes (\glspl{RNN}) extienden las arquitecturas feedforward incorporando conexiones cíclicas que permiten el procesamiento de datos secuenciales. Esta capacidad de mantener \textit{memoria} a través del tiempo las hace fundamentales para tareas que involucran secuencias temporales. Las \glspl{RNN} procesan geometría temporal (cadenas secuenciales) mientras agregan información con estado (memoria) que transforma progresivamente la señal.

\begin{definition}[Red neuronal recurrente]\label{def:rnn}
	Una \textit{\gls{RNN}} es una red neuronal con al menos un ciclo en su grafo de conectividad. Para el procesamiento de secuencias, el estado oculto evoluciona según:
	\[
	\mathbf{h}_t = f(\mathbf{h}_{t-1}, \mathbf{x}_t; \boldsymbol{\theta})
	\]
	donde $\mathbf{h}_t$ es el estado oculto en el tiempo $t$, $\mathbf{x}_t$ es la entrada actual, y $\boldsymbol{\theta}$ son los parámetros compartidos a través del tiempo.
\end{definition}

La forma más simple de una \gls{RNN} utiliza:
\begin{gather*}
    \mathbf{h}_t = \tanh(W_h \mathbf{h}_{t-1} + W_x \mathbf{x}_t + \mathbf{b}_h)\\
    \mathbf{y}_t = W_y \mathbf{h}_t + \mathbf{b}_y\\
\end{gather*}

donde:
\begin{itemize}
	\item $W_h \in \mathbb{R}^{d_h \times d_h}$ son los pesos recurrentes
	\item $W_x \in \mathbb{R}^{d_h \times d_x}$ son los pesos de entrada
	\item $W_y \in \mathbb{R}^{d_y \times d_h}$ son los pesos de salida
	\item $\mathbf{b}_h, \mathbf{b}_y$ son los vectores de sesgo
\end{itemize}

Un principio clave de las \gls{RNN} es que los mismos parámetros $\boldsymbol{\theta} = \{W_h, W_x, W_y, \mathbf{b}_h, \mathbf{b}_y\}$ se aplican en cada paso temporal. Esto permite:
\begin{itemize}
	\item \textbf{Generalización temporal}: El mismo modelo funciona para secuencias de cualquier longitud
	\item \textbf{Eficiencia paramétrica}: $O(d_h^2 + d_h d_x)$ parámetros independientes de la longitud de secuencia
	\item \textbf{Inductive bias temporal}: Asume que las mismas reglas de transición se aplican en todo momento
\end{itemize}

La \autoref{fig:rnn_unrolled} muestra cómo una RNN con ciclo (compacta) se ``desenvuelve'' en el tiempo para formar un \gls{DAG} que permite aplicar backpropagation.

\begin{figure}[h]
\centering
\begin{tikzpicture}[
    cell/.style={rectangle, draw, rounded corners, minimum width=1.2cm, minimum height=0.8cm, fill=gdlblue!15},
    input/.style={circle, draw, minimum size=0.5cm, fill=gdlgreen!20},
    output/.style={circle, draw, minimum size=0.5cm, fill=gdlred!20},
    hidden/.style={circle, draw, minimum size=0.5cm, fill=gdlyellow!20},
    arrow/.style={->, thick, >=stealth}
]

% --- RNN compacta (izquierda) ---
\node[font=\bfseries\small] at (-4.0, 3.0) {RNN compacta};
\node[font=\scriptsize] at (-4.0, 2.6) {(grafo cíclico)};

\node[cell] (rnn) at (-4.0, 0.5) {RNN};
\node[input] (xc) at (-4.0, -0.8) {$\mathbf{x}$};
\node[output] (yc) at (-4.0, 1.8) {$\mathbf{y}$};

\draw[arrow] (xc) -- (rnn);
\draw[arrow] (rnn) -- (yc);
% Ciclo de recurrencia: bucle claro a la derecha
\draw[arrow] (rnn.east) to[out=40, in=-40, distance=1.1cm] node[pos=0.25, above, font=\scriptsize] {$\mathbf{h}$} (rnn.east);

% --- Flecha de transformación ---
\draw[->, very thick, dashed] (-2.3, 0.5) -- (-1.3, 0.5) node[midway, above, font=\scriptsize] {unroll};

% --- RNN desplegada (derecha) ---
\node[font=\bfseries\small] at (3.5, 3.0) {RNN desplegada en $T$ pasos};
\node[font=\scriptsize] at (3.5, 2.6) {(DAG para BPTT)};

% Paso t-1
\node[cell] (r1) at (1.0, 0.5) {RNN};
\node[input] (x1) at (1.0, -0.8) {$\mathbf{x}_{t-1}$};
\node[output] (y1) at (1.0, 1.8) {$\mathbf{y}_{t-1}$};
\node[hidden] (h0) at (-0.5, 0.5) {};
\node[font=\scriptsize] at (-0.5, 0.1) {$\mathbf{h}_{t-2}$};

\draw[arrow] (x1) -- (r1);
\draw[arrow] (r1) -- (y1);
\draw[arrow] (h0) -- (r1);

% Paso t
\node[cell] (r2) at (3.5, 0.5) {RNN};
\node[input] (x2) at (3.5, -0.8) {$\mathbf{x}_t$};
\node[output] (y2) at (3.5, 1.8) {$\mathbf{y}_t$};

\draw[arrow] (x2) -- (r2);
\draw[arrow] (r2) -- (y2);
\draw[arrow] (r1) -- node[above, font=\scriptsize] {$\mathbf{h}_{t-1}$} (r2);

% Paso t+1
\node[cell] (r3) at (6, 0.5) {RNN};
\node[input] (x3) at (6, -0.8) {$\mathbf{x}_{t+1}$};
\node[output] (y3) at (6, 1.8) {$\mathbf{y}_{t+1}$};
\node[hidden] (hf) at (8, 0.5) {};

\draw[arrow] (x3) -- (r3);
\draw[arrow] (r3) -- (y3);
\draw[arrow] (r2) -- node[above, font=\scriptsize] {$\mathbf{h}_t$} (r3);
\draw[arrow] (r3) -- node[above, font=\scriptsize] {$\mathbf{h}_{t+1}$} (hf);

% Nota sobre parámetros compartidos
\draw[decorate, decoration={brace, amplitude=5pt, mirror}] (0.5, -1.5) -- (6.5, -1.5);
\node[font=\scriptsize, align=center] at (3.5, -2.0) {Mismos parámetros $\boldsymbol{\theta}$ en todos los pasos};

\end{tikzpicture}
\caption[RNN compacta vs desplegada en el tiempo]{Izquierda: Representación compacta de una RNN como grafo cíclico. Derecha: La misma RNN ``desplegada'' (unrolled) en el tiempo, mostrando cómo el estado oculto $\mathbf{h}_t$ se propaga entre pasos temporales. Esta representación como DAG permite aplicar Backpropagation Through Time (BPTT).}
\label{fig:rnn_unrolled}
\end{figure}

Las \gls{RNN} como la que acabamos de describir, que únicamente tiene una conexión recurrente sin ninguna complejidad extra, sufren de varias limitaciones:
\begin{itemize}
	\item \textbf{Memoria limitada}: Dificultad para capturar dependencias a muy largo plazo
	\item \textbf{Sesgo reciente}: Tendencia a dar más peso a información reciente
	\item \textbf{Paralelización limitada}: El procesamiento secuencial impide paralelización temporal
	\item \textbf{Saturación de activación}: Las funciones de activación sigmoidales causan saturación
\end{itemize}

Además este tipo de redes necesitan de una variación del algoritmo de aprendizaje que se usa normalmente para entrenar las redes de neuronas (ver~\ref{sec:learning_ann}): \gls{BPTT}~\citep{werbos1990bptt}

\subsubsection{LSTMs}\label{subsubsec:lstm}

Las redes Long Short-Term Memory (\gls{LSTM}), introducidas por Hochreiter y Schmidhuber~\citep{hochreiter1997long}, representan una evolución crucial de las \gls{RNN} que aborda específicamente el problema del gradiente desvaneciente y permite el aprendizaje de dependencias a largo plazo. Las \glspl{LSTM} son variantes más sofisticadas de \glspl{RNN} que procesan geometría temporal mediante mecanismos de compuertas que controlan el flujo de información. La arquitectura original constaba de dos compuertas (entrada y salida); la compuerta de olvido (\textit{forget gate}) fue añadida posteriormente por \citet{gers2000learning}.

\begin{definition}[LSTM Cell]\label{def:lstm_cell}
	Una \textit{celda \gls{LSTM}} mantiene dos estados: un estado celular $\mathbf{C}_t$ y un estado oculto $\mathbf{h}_t$, controlados por tres compuertas (gates):
	\begin{itemize}
		\item \textbf{Forget gate}: $\mathbf{f}_t = \sigma(W_f \cdot [\mathbf{h}_{t-1}, \mathbf{x}_t] + \mathbf{b}_f)$
		\item \textbf{Input gate}: $\mathbf{i}_t = \sigma(W_i \cdot [\mathbf{h}_{t-1}, \mathbf{x}_t] + \mathbf{b}_i)$
		\item \textbf{Output gate}: $\mathbf{o}_t = \sigma(W_o \cdot [\mathbf{h}_{t-1}, \mathbf{x}_t] + \mathbf{b}_o)$
	\end{itemize}
	El estado celular $\mathbf{C}_t$ se actualiza mediante:
	\begin{gather*}
		\tilde{\mathbf{C}}_t = \tanh(W_C \cdot [\mathbf{h}_{t-1}, \mathbf{x}_t] + \mathbf{b}_C)\\
		\mathbf{C}_t = \mathbf{f}_t \odot \mathbf{C}_{t-1} + \mathbf{i}_t \odot \tilde{\mathbf{C}}_t\\
		\mathbf{h}_t = \mathbf{o}_t \odot \tanh(\mathbf{C}_t)
	\end{gather*}
	donde $\odot$ denota el producto elemento a elemento.
\end{definition}

Cada compuerta tiene una función específica:

\begin{itemize}
	\item \textbf{Forget gate} ($\mathbf{f}_t$): Decide qué información descartar del estado celular anterior
	\begin{itemize}
		\item $f_{t,i} \approx 0$: Olvida completamente la información del componente $i$
		\item $f_{t,i} \approx 1$: Retiene completamente la información del componente $i$
	\end{itemize}

	\item \textbf{Input gate} ($\mathbf{i}_t$): Controla qué nueva información almacenar en el estado celular
	\begin{itemize}
		\item Trabaja junto con $\tilde{\mathbf{C}}_t$ para determinar qué valores actualizar
		\item Permite escritura selectiva en la memoria
	\end{itemize}

	\item \textbf{Output gate} ($\mathbf{o}_t$): Controla qué partes del estado celular exponer como salida
	\begin{itemize}
		\item Filtra el estado interno para producir el estado oculto
		\item Permite lectura selectiva de la memoria
	\end{itemize}
\end{itemize}

La clave está en que el forget gate puede aprender a mantenerse cerca de 1 para información importante, creando \textit{caminos de gradiente} que persisten a través de largos períodos temporales.

La \autoref{fig:lstm_arquitectura} ilustra la arquitectura interna de una celda LSTM, mostrando el flujo de información a través de las tres compuertas y el estado celular.

\begin{figure}[htbp]
\centering
\begin{tikzpicture}[
    % Estilos mejorados para mejor legibilidad
    gate/.style={
        rectangle, draw, rounded corners=4pt,
        minimum width=1.6cm, minimum height=1cm,
        fill=#1!25, font=\small\bfseries,
        line width=0.6pt
    },
    op/.style={
        rectangle, draw, rounded corners=3pt,
        minimum width=1.3cm, minimum height=0.9cm,
        fill=gdlyellow!30, font=\normalsize\bfseries,
        line width=0.6pt
    },
    state/.style={
        rectangle, draw, rounded corners=3pt,
        minimum width=1.4cm, minimum height=0.8cm,
        fill=gdlgray!15, font=\small,
        line width=0.6pt
    },
    conn/.style={->, >=stealth, draw=gdlgray!65, line width=0.7pt},
    main/.style={->, >=stealth, line width=0.8pt}
]

% ===== ESTADO CELULAR (línea superior) =====
\draw[->, >=stealth, line width=1.5pt, gdlblue!60] (-1.5,5) -- (13,5);
\node[above, font=\small, gdlblue!70] at (5.25,5.5) {Estado celular $\mathbf{C}_t$};
\node[left, font=\small] at (-1.5,5) {$\mathbf{C}_{t-1}$};

% ===== REGIÓN 1: FORGET GATE =====
\node[gate=gdlred] (fg) at (1.5,2) {$\sigma$};
\node[below=0.15cm of fg, font=\scriptsize, gdlred!60!black] {Forget Gate};
\node[op] (fgmult) at (1.5,5) {$\odot$};

% ===== REGIÓN 2: INPUT GATE + CANDIDATO =====
\node[gate=gdlgreen] (ig) at (4.5,2) {$\sigma$};
\node[below=0.15cm of ig, font=\scriptsize, gdlgreen!40!black] {Input Gate};

\node[gate=gdlblue] (cand) at (7,2) {tanh};
\node[below=0.15cm of cand, font=\scriptsize, gdlblue!50!black] {Candidato $\tilde{\mathbf{C}}_t$};

\node[op] (igmult) at (5.75,3.5) {$\odot$};
\node[op] (add) at (5.75,5) {$\oplus$};

% ===== REGIÓN 3: OUTPUT GATE =====
\node[gate=gdlpurple] (og) at (9.5,2) {$\sigma$};
\node[below=0.15cm of og, font=\scriptsize, gdlpurple!60!black] {Output Gate};

\node[gate=gdlblue] (tanho) at (9.5,4) {tanh};
\node[op] (ogmult) at (11,3) {$\odot$};

% ===== ENTRADAS =====
\node[state] (xt) at (5.25,-0.8) {$\mathbf{x}_t$};
\node[state] (ht1) at (-0.8,0.5) {$\mathbf{h}_{t-1}$};

% ===== SALIDA =====
\node[state, fill=gdlorange!15] (ht) at (14,3) {$\mathbf{h}_t$};

% ===== CONEXIONES DE ENTRADA (más separadas y tenues) =====
% Concatenación visual [h_{t-1}, x_t]
\coordinate (concat) at (5.25, 0.5);

% h_{t-1} hacia punto de concatenación
\draw[conn] (ht1) -- (concat);

% x_t hacia punto de concatenación
\draw[conn] (xt) -- (concat);

% Desde concatenación hacia cada compuerta (rutas separadas)
\draw[conn] (concat) -- (1.5, 0.5) -- (fg);
\draw[conn] (concat) -- (4.5, 0.5) -- (ig);
\draw[conn] (concat) -- (7, 0.5) -- (cand);
\draw[conn] (concat) -- (9.5, 0.5) -- (og);

% Etiqueta de concatenación
\node[font=\scriptsize] at (5.25, 0.15) {$[\mathbf{h}_{t-1}, \mathbf{x}_t]$};

% ===== FLUJO INTERNO =====
% Forget gate → multiplicación con estado
\draw[main] (fg) -- (fgmult);

% Input gate y candidato → multiplicación
\draw[main, gdlgreen!60!black] (ig) |- (igmult);
\draw[main, gdlgreen!60!black] (cand) |- (igmult);

% Multiplicación input → suma con estado
\draw[main] (igmult) -- (add);

% Estado anterior pasa por forget
\draw[main, gdlblue!50] (fgmult) -- (add);

% Nuevo estado hacia output tanh
\draw[main, gdlblue!50] (8.2,5) -- (8.2,4) -- (tanho.west);

% Output gate y tanh → multiplicación final
\draw[main, gdlpurple!60!black] (og) -| (ogmult);
\draw[main, gdlgreen!50!black] (tanho) -| (ogmult);

% Multiplicación final → salida h_t
\draw[main] (ogmult) -- (ht);

% ===== RECURRENCIA (hacia siguiente paso) =====
\draw[conn, dashed, line width=0.8pt] (ht.north) -- (14,6.3) -- (-0.8,6.3) -- (-0.8,0.9)
    node[pos=0.5, above, font=\scriptsize] {$\mathbf{h}_t$ recurrente};

% ===== SEPARADORES DE REGIONES (opcional, muy sutil) =====
\draw[gdlgray!20, dashed, line width=0.3pt] (3,-.5) -- (3,5.5);
\draw[gdlgray!20, dashed, line width=0.3pt] (8.25,-.5) -- (8.25,5.5);

\end{tikzpicture}
\caption[Arquitectura de celda LSTM]{Arquitectura interna de una celda LSTM. El estado celular $\mathbf{C}_t$ (línea superior azul) fluye horizontalmente, siendo modulado por las tres compuertas: la \textit{forget gate} (roja) decide qué información descartar, la \textit{input gate} (verde) junto con el candidato $\tilde{\mathbf{C}}_t$ determinan qué nueva información almacenar, y la \textit{output gate} (púrpura) controla qué parte del estado celular exponer como salida $\mathbf{h}_t$. Los operadores $\odot$ representan multiplicación elemento a elemento y $\oplus$ la suma.}
\label{fig:lstm_arquitectura}
\end{figure}

\textbf{Gated Recurrent Unit (GRU)}: Propuesta por Cho et al.~\citep{cho2014learning}, es una simplificación del \gls{LSTM} que combina forget e input gates en una sola compuerta de actualización. Los estudios empíricos han mostrado que el rendimiento de \gls{GRU} es comparable al de \gls{LSTM} en muchas tareas~\citep{chung2014empirical}:
\begin{align}
\mathbf{z}_t &= \sigma(W_z \cdot [\mathbf{h}_{t-1}, \mathbf{x}_t] + \mathbf{b}_z) \\
\mathbf{r}_t &= \sigma(W_r \cdot [\mathbf{h}_{t-1}, \mathbf{x}_t] + \mathbf{b}_r) \\
\tilde{\mathbf{h}}_t &= \tanh(W_h \cdot [\mathbf{r}_t \odot \mathbf{h}_{t-1}, \mathbf{x}_t] + \mathbf{b}_h) \\
\mathbf{h}_t &= (1 - \mathbf{z}_t) \odot \mathbf{h}_{t-1} + \mathbf{z}_t \odot \tilde{\mathbf{h}}_t
\end{align}

Las \gls{LSTM} han sido fundamentales para avances en traducción automática, reconocimiento de voz y modelado de lenguaje, estableciendo las bases para arquitecturas más modernas como los \glspl{Transformer} que abordan la limitación de paralelización temporal.

\subsubsection{Residuales}

Las conexiones residuales, popularizadas por He et al.~\citep{he2016deep} en las ResNets, representan un avance arquitectónico fundamental que permitió el entrenamiento efectivo de redes neuronales extremadamente profundas. Su impacto trasciende las redes convolucionales, influenciando prácticamente todas las arquitecturas modernas de deep learning. No son una capa en sentido estricto, sino una forma de conectar capas existentes.

\begin{definition}[Conexión residual]\label{def:conexion_residual}
	Una conexión residual (o skip connection) es un tipo de arquitectura donde la salida de una capa se suma directamente a la salida de una capa posterior. Para un bloque de capas $\mathcal{F}$:
	\[
	\mathbf{y} = \mathcal{F}(\mathbf{x}) + \mathbf{x}
	\]
	donde $\mathbf{x}$ es la entrada del bloque y $\mathcal{F}(\mathbf{x})$ representa la transformación aprendida.
\end{definition}

La \autoref{fig:skip_connection} compara un bloque tradicional con un bloque residual, mostrando cómo la conexión directa permite que el gradiente fluya sin atenuación.

\begin{figure}[h]
\centering
\begin{tikzpicture}[
    node distance=1cm,
    block/.style={rectangle, draw, minimum width=1.8cm, minimum height=0.7cm, fill=gdlblue!15},
    sum/.style={circle, draw, minimum size=0.5cm, fill=gdlyellow!20},
    arrow/.style={->, thick, >=stealth},
    skip/.style={->, thick, >=stealth, dashed, color=gdlred!70}
]

% Bloque tradicional (izquierda)
\node[font=\bfseries] at (-2, 2.5) {Sin residual};
\node[block] (f1) at (-2, 1) {$\mathcal{F}_1$};
\node[block] (f2) at (-2, -0.5) {$\mathcal{F}_2$};
\draw[arrow] (-2, 2) -- (f1);
\draw[arrow] (f1) -- (f2);
\draw[arrow] (f2) -- (-2, -1.5);
\node[font=\scriptsize] at (-2, 2.2) {$\mathbf{x}$};
\node[font=\scriptsize] at (-2, -1.7) {$\mathbf{y}$};

% Bloque residual (derecha)
\node[font=\bfseries] at (2, 2.5) {Con residual};
\node[block] (r1) at (2, 1) {$\mathcal{F}_1$};
\node[block] (r2) at (2, -0.5) {$\mathcal{F}_2$};
\node[sum] (plus) at (2, -1.5) {$+$};

\draw[arrow] (2, 2) -- (r1);
\draw[arrow] (r1) -- (r2);
\draw[arrow] (r2) -- (plus);
\draw[arrow] (plus) -- (2, -2.3);
\node[font=\scriptsize] at (2, 2.2) {$\mathbf{x}$};
\node[font=\scriptsize] at (2, -2.5) {$\mathbf{y} = \mathcal{F}(\mathbf{x}) + \mathbf{x}$};

% Skip connection
\draw[skip] (2, 1.8) -- (3.5, 1.8) -- (3.5, -1.5) -- (plus);
\node[font=\scriptsize, color=gdlred!70] at (3.8, 0) {skip};

% Gradiente
\node[font=\footnotesize, align=center] at (-2, -2.5) {Gradiente:\\se atenúa};
\node[font=\footnotesize, align=center] at (2, -3.3) {Gradiente:\\fluye directo};

\end{tikzpicture}
\caption[Comparación de bloques sin y con conexión residual]{Comparación de bloques sin (izquierda) y con (derecha) conexión residual. La skip connection (línea punteada roja) permite que el gradiente fluya directamente durante la retropropagación, evitando el desvanecimiento en redes profundas.}
\label{fig:skip_connection}
\end{figure}

Antes de las ResNets, era empíricamente observado que redes más profundas no necesariamente generalizaban mejor, incluso cuando el error de entrenamiento era menor. El uso de las conexiones residuales reducía el problema del desvanecimiento del gradiente.

\section{Redes de neuronas como grafos}\label{sec:redes_grafos}

En las secciones anteriores hemos definido las capas neuronales y su composición de forma algebraica, como funciones que transforman señales. Sin embargo, existe una perspectiva complementaria igualmente importante: la representación de redes neuronales como \textit{grafos}, donde los nodos representan unidades de cómputo y las aristas representan el flujo de información.

Esta perspectiva de grafos es fundamental por varias razones:
\begin{itemize}
    \item Proporciona un lenguaje matemático preciso para describir la \textit{arquitectura} de una red
    \item Permite distinguir entre diferentes familias de redes (feedforward, recurrentes, etc.) mediante propiedades topológicas
    \item Es la base de la \textit{diferenciación automática} utilizada en frameworks modernos de deep learning
    \item Conecta naturalmente con el campo de \gls{GDL}, donde la propia estructura del grafo codifica información sobre el dominio de los datos
\end{itemize}

\subsection{Grafos}

Hasta ahora hemos presentado las redes neuronales como composiciones de capas de forma algo informal. Para formalizar rigurosamente su estructura, recurrimos a la teoría de grafos, que proporciona el lenguaje matemático preciso para describir las relaciones de conectividad entre neuronas y capas.

\begin{definition}[\Gls{grafo}]\label{def:grafo_neuronal}
Un \textit{\gls{grafo}} es una estructura matemática $\mathcal{G} = (V, E)$ que consta de:
\begin{itemize}
    \item Un conjunto $V$ de \textit{vértices} (también llamados nodos)
    \item Un conjunto $E$ de \textit{aristas} (también llamadas enlaces o conexiones), donde cada arista conecta un par de vértices
\end{itemize}
Formalmente, para un grafo no dirigido, $E \subseteq \{\{u, v\} : u, v \in V, u \neq v\}$, mientras que para un grafo dirigido, $E \subseteq V \times V$\footnote{Nótese que la definición de grafo dirigido permite auto-lazos $(v, v)$, lo cual es relevante para redes recurrentes donde una neurona puede realimentar su propia salida.}.
\end{definition}

Los grafos proporcionan una abstracción natural para modelar relaciones entre objetos. En el contexto de redes neuronales, los vértices representan neuronas (o capas, según el nivel de abstracción) y las aristas representan conexiones ponderadas que transmiten señales.

\begin{definition}[\Gls{grafo_dirigido}]\label{def:grafo_dirigido}
Un \textit{\gls{grafo_dirigido}} es un grafo $\mathcal{G} = (V, \vec{E})$ donde cada arista tiene una dirección específica. Una arista dirigida $(u, v) \in \vec{E}$ indica que hay una conexión del vértice $u$ (origen) al vértice $v$ (destino). La dirección importa: $(u, v) \neq (v, u)$.
\end{definition}

\begin{definition}[Grafo no dirigido]\label{def:grafo_no_dirigido}
Un \textit{grafo no dirigido} es un grafo $\mathcal{G} = (V, E)$ donde las aristas son simétricas: si $(u, v) \in E$, entonces $(v, u) \in E$. Equivalentemente, $E$ puede verse como un conjunto de pares no ordenados $\{\{u, v\} : u, v \in V\}$.
\end{definition}

En una \gls{ANN}, las aristas dirigidas representan el flujo de información: la salida de una neurona $u$ se convierte en entrada para la neurona $v$. Este flujo direccional es fundamental para entender cómo se propagan las señales durante la inferencia (forward pass) y cómo se retropropagan los gradientes durante el entrenamiento (backward pass).

\begin{definition}[Camino en un grafo dirigido]\label{def:camino}
Dado un \gls{grafo_dirigido} $\mathcal{G} = (V, \vec{E})$, un \textit{camino} de longitud $k$ es una secuencia de vértices $v_0, v_1, \ldots, v_k$ tal que $(v_i, v_{i+1}) \in \vec{E}$ para todo $i = 0, \ldots, k-1$.
\end{definition}

\begin{definition}[Ciclo en un grafo dirigido]\label{def:ciclo}
Dado un \gls{grafo_dirigido} $\mathcal{G} = (V, \vec{E})$, un \textit{ciclo} es un camino $v_0, v_1, \ldots, v_k$ donde $v_0 = v_k$, $k \geq 1$, y los vértices $v_0, v_1, \ldots, v_{k-1}$ son distintos dos a dos.
\end{definition}

Un ciclo en una red neuronal indica retroalimentación: la salida de una neurona (o capa) puede, directa o indirectamente, influir en su propia entrada. Esto es fundamental en redes recurrentes que procesan secuencias temporales.

\begin{definition}[\Gls{grafo_aciclico_dirigido} (DAG)]\label{def:dag}
Un \textit{\gls{grafo_aciclico_dirigido}} (DAG, por sus siglas en inglés Directed Acyclic Graph) es un \gls{grafo_dirigido} que no contiene ciclos.
\end{definition}

Los \glspl{DAG} tienen propiedades computacionales especialmente útiles:

\begin{theorem}[Ordenamiento topológico]\label{thm:ordenamiento_topologico}
Todo grafo acíclico dirigido $\mathcal{G} = (V, \vec{E})$ admite al menos un \textit{ordenamiento topológico}: una numeración de sus vértices $v_1, v_2, \ldots, v_{|V|}$ tal que si $(v_i, v_j) \in \vec{E}$, entonces $i < j$.
\end{theorem}

\begin{proof}
Procedemos por inducción en el número de vértices $n = |V|$.

\textbf{Caso base} ($n = 1$): Un grafo con un solo vértice y sin aristas (por ser acíclico) tiene trivialmente un ordenamiento topológico.

\textbf{Paso inductivo}: Supongamos que todo \gls{DAG} con $n$ vértices tiene un ordenamiento topológico. Consideremos un \gls{DAG} $\mathcal{G} = (V, \vec{E})$ con $n+1$ vértices.

Afirmamos que $\mathcal{G}$ debe tener al menos un vértice $v$ sin aristas entrantes (grado de entrada cero). En efecto, si todos los vértices tuvieran al menos una arista entrante, podríamos construir un camino infinito retrocediendo indefinidamente, lo que por el principio del palomar (hay finitos vértices) implicaría un ciclo, contradiciendo que $\mathcal{G}$ es acíclico.

Seleccionamos un vértice $v \in V$ con grado de entrada cero y lo colocamos primero en el ordenamiento. Consideremos el grafo reducido $\mathcal{G}' = (V \setminus \{v\}, \vec{E}')$ donde $\vec{E}'$ es $\vec{E}$ restringido a $V \setminus \{v\}$. Este grafo tiene $n$ vértices y sigue siendo acíclico (eliminar un vértice no introduce ciclos). Por hipótesis de inducción, $\mathcal{G}'$ tiene un ordenamiento topológico. Concatenando $v$ al inicio de este ordenamiento obtenemos un ordenamiento topológico de $\mathcal{G}$.
\end{proof}

La existencia de un ordenamiento topológico garantiza que podemos computar las salidas de todas las neuronas en un orden secuencial sin necesidad de resolver sistemas de ecuaciones implícitas, lo que hace que las redes feedforward sean computacionalmente eficientes.

\subsection{Redes neuronales como grafos de neuronas}

Una \gls{ANN} en su forma más general puede representarse como un \gls{grafo_dirigido} $\mathcal{G}_{\text{neuronas}} = (V_N, E_N)$ donde:
\begin{itemize}
    \item $V_N$: conjunto de neuronas de la red
    \item $E_N$: conjunto de conexiones ponderadas. Una arista $(n_i, n_j) \in E_N$ con peso $w_{ij} \in \mathbb{R}$ indica que la salida de la neurona $n_i$ se multiplica por $w_{ij}$ y contribuye a la entrada de la neurona $n_j$
\end{itemize}

\subsubsection[Redes feedforward: DAGs de neuronas]{Redes feedforward: \glspl{DAG} de neuronas}

Las redes neuronales feedforward o de propagación hacia adelante son aquellas cuyo grafo de neuronas es un \gls{DAG}. La ausencia de ciclos garantiza que:

\begin{enumerate}
    \item \textbf{Computación determinística en un paso}: Dada una entrada, podemos calcular la salida de cada neurona exactamente una vez siguiendo un ordenamiento topológico. Si $\sigma_i$ denota la ordenación topológica de la neurona $i$, entonces:
    \[
    y_i = \sigma_i\left(\sum_{(j,i) \in E_N} w_{ji} y_j + b_i\right)
    \]
    donde $y_j$ ya ha sido computado para todos los predecesores $j$ de $i$.

    \item \textbf{Paralelización eficiente}: Todas las neuronas en el mismo ``nivel'' del \gls{DAG} (misma distancia desde la entrada) pueden computarse en paralelo, lo que permite aprovechar arquitecturas de hardware modernas (GPUs, TPUs).

    \item \textbf{Backpropagation eficiente}: El cálculo de gradientes mediante la regla de la cadena puede realizarse en un único pase hacia atrás siguiendo el orden topológico inverso. Veremos más sobre el aprendizaje en \glspl{ANN} en la sección \ref{sec:learning_ann}.
\end{enumerate}

\subsubsection{Redes recurrentes: grafos cíclicos}

Las \glspl{RNN} y otras arquitecturas con retroalimentación se modelan mediante grafos dirigidos \textit{con ciclos}. Un ciclo representa memoria o estado interno: la salida en el paso temporal $t$ depende no solo de la entrada actual, sino también de estados internos de pasos anteriores.

Matemáticamente, esto se traduce en ecuaciones de la forma:
\[
h_t = f(h_{t-1}, x_t; \theta)
\]
donde $h_t$ es el estado oculto en el tiempo $t$, $x_t$ es la entrada actual, y $\theta$ son los parámetros. El ciclo implícito $(h_{t-1} \to h_t \to h_{t+1} \to \cdots)$ permite que la red capture dependencias temporales.

A pesar de que los grafos cíclicos (redes recurrentes) tienen mayor expresividad teórica para modelar dinámicas temporales, las arquitecturas modernas tienden a favorecer \glspl{DAG} por razones prácticas:

\begin{enumerate}
    \item \textbf{Eficiencia computacional}: Los \glspl{DAG} permiten paralelización masiva, crucial para entrenar modelos con millones o miles de millones de parámetros.

    \item \textbf{Estabilidad numérica}: Las redes recurrentes sufren de problemas de gradientes que explotan o se desvanecen a lo largo de secuencias largas, dificultando su entrenamiento.

    \item \textbf{Simplicidad conceptual}: Los \glspl{DAG} son más fáciles de analizar teóricamente y depurar en la práctica.
\end{enumerate}

No obstante, mecanismos como la atención y las arquitecturas \gls{Transformer} (véase el \autoref{ch:attention}) han logrado capturar dependencias de largo alcance manteniendo la estructura de \gls{DAG}, combinando lo mejor de ambos mundos.

\subsection{Capas como grafos de transformaciones}\label{subsec:capas_grafos_transformaciones}

Además de entender las \gls{ANN} como grafos donde los \textit{nodos son neuronas individuales}, existe una perspectiva complementaria de mayor nivel de abstracción que es fundamental para entender arquitecturas modernas: pensar en las redes como \textit{grafos de capas}, donde cada nodo representa una capa completa (potencialmente con miles de neuronas) y las aristas representan el flujo de datos entre capas.

\begin{definition}[\Gls{grafo_de_capas}]\label{def:grafo_de_capas}
Dada una red neuronal compuesta por $k$ capas $L_1, L_2, \ldots, L_k$, su \textit{\gls{grafo_de_capas}} es un \gls{grafo_dirigido} $\mathcal{G}_{\mathcal{N}} = (V_L, E_L, V_{\text{in}}, V_{\text{out}})$ donde:
\begin{itemize}
    \item $V_L = \{L_1, L_2, \ldots, L_k\}$: el conjunto de capas
    \item $(L_i, L_j) \in E_L$ si y solo si la salida de la capa $L_i$ es (parte de) la entrada de la capa $L_j$
    \item $V_{\text{in}} \subseteq V_L$: las \textit{capas de entrada}, aquellas con grado de entrada cero (no reciben salida de otras capas)
    \item $V_{\text{out}} \subseteq V_L$: las \textit{capas de salida}, aquellas con grado de salida cero (su salida no alimenta otras capas)
\end{itemize}
\end{definition}

Este nivel de abstracción simplifica la visualización de arquitecturas complejas y facilita el análisis del flujo de información a través de la red. Además, refleja naturalmente la estructura algebraica: cada capa $L_i$ se puede descomponer como una transformación lineal (representable matricialmente) seguida de una función de activación no lineal, $L_i(x) = \sigma(W_i x + b_i)$, y las aristas del grafo corresponden a estas composiciones de transformaciones entre espacios vectoriales.

\subsubsection[Arquitecturas feedforward como DAGs de capas]{Arquitecturas feedforward como \glspl{DAG} de capas}

Las arquitecturas feedforward clásicas (como perceptrones multicapa, \glspl{CNN} convolucionales estándar) tienen grafos de capas que son \glspl{DAG} con una estructura particularmente simple: un camino lineal.

\begin{example}[Red neuronal feedforward]
Una red neuronal feedforward con $k$ capas se representa como:
\[
\mathcal{N}(x) = L_k \circ L_{k-1} \circ \cdots \circ L_2 \circ L_1(x)
\]
Su grafo de capas es el camino dirigido:
\[
L_1 \to L_2 \to L_3 \to \cdots \to L_{k-1} \to L_k
\]
Este es el \gls{DAG} más simple posible: completamente lineal, sin ramificaciones ni conexiones saltadas (skip connections).
\end{example}

\subsubsection[Arquitecturas modernas con DAGs no lineales]{Arquitecturas modernas con \glspl{DAG} no lineales}

Las arquitecturas modernas han abandonado la restricción de grafos lineales, explorando \glspl{DAG} más complejos que permiten flujos de información más ricos:

Las conexiones residuales, introducidas en ResNet~\citep{he2016deep}, agregan aristas que ``saltan'' capas:
\[
h^{(\ell+1)} = h^{(\ell)} + F(h^{(\ell)}; \theta_\ell)
\]
donde $F$ es una función aprendida. En el grafo de capas, esto crea una arista directa $L_\ell \to L_{\ell+1}$ que coexiste con el camino a través de las capas intermedias, resultando en un \gls{DAG} no lineal.

Redes como Inception~\citep{szegedy2015going} procesan la entrada a través de múltiples caminos paralelos que luego se concatenan o suman:
\[
h^{(\ell+1)} = \text{Concatenar}(F_1(h^{(\ell)}), F_2(h^{(\ell)}), F_3(h^{(\ell)}))
\]
El grafo de capas tiene una estructura de árbol invertido: un nodo (entrada) se ramifica en múltiples caminos que posteriormente se fusionan.

\subsubsection{Redes recurrentes como grafos cíclicos de capas}

Las \glspl{RNN} y \glspl{LSTM} introducidas en la Sección~\ref{subsubsec:rnn}, por el contrario, tienen grafos de capas con ciclos:

\begin{example}[Red neuronal recurrente básica]
Una \gls{RNN} procesa una secuencia $x_1, x_2, \ldots, x_T$ manteniendo un estado oculto $h_t$ que se actualiza recursivamente:
\[
h_t = f(W_{hh} h_{t-1} + W_{xh} x_t + b)
\]
El grafo de capas contiene un ciclo auto-dirigido: la capa recurrente tiene una arista hacia sí misma, representando la dependencia temporal $h_{t-1} \to h_t$.
\end{example}

\subsubsection{Conexión entre perspectivas de grafo de neuronas y grafo de capas}

Las dos perspectivas (grafo de neuronas y grafo de capas) están relacionadas jerárquicamente:
\begin{itemize}
    \item El \textbf{grafo de neuronas} $\mathcal{G}_{\text{neuronas}} = (V_N, E_N)$ proporciona la descripción microscópica completa de la red
    \item El \textbf{grafo de capas} $\mathcal{G}_{\mathcal{N}} = (V_L, E_L)$ es una abstracción macroscópica donde cada nodo $L_i \in V_L$ representa un subgrafo completo de neuronas
\end{itemize}

Formalmente, existe una función de proyección $\pi: V_N \to V_L$ que asigna cada neurona a su capa correspondiente. Las aristas en $E_L$ resumen la conectividad entre grupos de neuronas:
\[
(L_i, L_j) \in E_L \iff \exists n \in \pi^{-1}(L_i), m \in \pi^{-1}(L_j) : (n, m) \in E_N
\]

Esta abstracción de dos niveles es fundamental: permite razonar sobre arquitecturas globales (grafos de capas) sin perderse en los detalles de millones de conexiones individuales (grafo de neuronas), al tiempo que se preserva la capacidad de computar gradientes exactos a nivel de cada parámetro individual.

\section{Aprendizaje en redes de neuronas}\label{sec:learning_ann}

El aprendizaje o entrenamiento de una red neuronal es un proceso iterativo que implica la optimización de sus parámetros \(\theta\) o pesos para minimizar un error. A la \gls{ANN} se le proporcionan conjuntos de datos de entrenamiento que contienen ejemplos de la salida esperada de la \gls{ANN} dado unos datos de entrada. Con el error en cada paso se ejecuta un algoritmo de optimización de los parámetros \(\theta\), en particular se utiliza el descenso del gradiente o alguna de sus variedades como el descenso de gradiente estocástico (\gls{SGD}).

\subsection{Loss function}

El error de la \gls{ANN} es una medida de la diferencia entre la salida esperada y la salida de la red. Esta medida del error se realiza con una función denominada función de pérdida (\emph{loss function}) o \textit{función de coste}.

\begin{definition}[Función de pérdida]\label{def:funcion_perdida}
	Una \textit{función de pérdida} es una función \(\mathcal{L}: \mathcal{Y} \times \mathcal{Y} \to \mathbb{R}_{\geq 0}\) que mide la discrepancia entre la predicción \(\hat{y} \in \mathcal{Y}\) y el valor verdadero \(y \in \mathcal{Y}\), satisfaciendo:
	\begin{enumerate}
		\item \textbf{No negatividad}: \(\mathcal{L}(y, \hat{y}) \geq 0\) para todo \(y, \hat{y} \in \mathcal{Y}\)
		\item \textbf{Consistencia}: \(\mathcal{L}(y, y) = 0\) para todo \(y \in \mathcal{Y}\) (la predicción perfecta no incurre pérdida)
		\item \textbf{Diferenciabilidad}: \(\mathcal{L}\) es diferenciable respecto a \(\hat{y}\) para permitir optimización basada en gradientes
	\end{enumerate}
\end{definition}

\begin{remark}
Algunas funciones de pérdida satisfacen la condición más fuerte \(\mathcal{L}(y, \hat{y}) = 0 \iff y = \hat{y}\) (por ejemplo, MSE y MAE), mientras que otras alcanzan el mínimo cero en una región más amplia: la pérdida hinge (\autoref{def:hinge_loss}) satisface \(\mathcal{L}(y, \hat{y}) = 0\) cuando \(y \hat{y} \geq 1\), no solo cuando \(y = \hat{y}\).
\end{remark}

La función de pérdida se puede utilizar en más de una capa de la \gls{ANN}, se pueden combinar varias, e incluso se puede utilizar con los parámetros \(\theta\) de la \gls{ANN} para evitar que crezcan de forma desmedida\footnote{Aplicar una función de pérdida sobre los pesos de una red es un proceso que se llama regularización y sirve para evitar o disminuir errores acumulativos debidos a la capacidad de representación de los ordenadores y también para forzar una mayor generalización del aprendizaje.}. En el contexto de este libro solo consideraremos una única función de pérdida aplicada sobre la última capa de la \gls{ANN}, que constará de más de una salida.


\subsubsection{Funciones de pérdida comunes}

A continuación presentamos las funciones de pérdida más utilizadas en la práctica, organizadas según el tipo de tarea.

\paragraph{Regresión}

Los problemas de \textit{regresión} consisten en predecir un valor numérico continuo a partir de las características de entrada. Matemáticamente, buscamos aproximar una función \(f: \mathcal{X} \to \mathbb{R}^d\) donde el espacio de salida es un subconjunto de \(\mathbb{R}^d\). Ejemplos típicos incluyen la predicción de precios de viviendas, la estimación de temperaturas futuras, o la regresión de coordenadas espaciales. Las funciones de pérdida para regresión miden la distancia entre el valor predicho \(\hat{y}\) y el valor real \(y\) en el espacio euclídeo.

\begin{definition}[Error cuadrático medio (MSE)]\label{def:mse}
	El \textit{error cuadrático medio} (\gls{MSE}) para un conjunto de \(n\) predicciones es:
	\[
		\mathcal{L}_{\text{MSE}}(y, \hat{y}) = \frac{1}{n} \sum_{i=1}^{n} (y_i - \hat{y}_i)^2
	\]
	Para una sola predicción: \(\mathcal{L}(y, \hat{y}) = (y - \hat{y})^2\).
\end{definition}

El \gls{MSE} penaliza fuertemente los errores grandes debido al término cuadrático, lo que lo hace sensible a valores atípicos (\textit{outliers}). Su gradiente es \(\nabla_{\hat{y}} \mathcal{L} = -2(y - \hat{y})\), proporcional al error.

\begin{definition}[Error absoluto medio (MAE)]\label{def:mae}
	El \textit{error absoluto medio} (\gls{MAE}) es:
	\[
		\mathcal{L}_{\text{MAE}}(y, \hat{y}) = \frac{1}{n} \sum_{i=1}^{n} |y_i - \hat{y}_i|
	\]
\end{definition}

El \gls{MAE} es más robusto a \textit{outliers} que el \gls{MSE}, pero presenta un problema: no es diferenciable en \(\hat{y} = y\). En la práctica, se usa el subgradiente \(\text{sign}(y - \hat{y})\).

\begin{definition}[Pérdida de Huber]\label{def:huber}
	La \textit{pérdida de Huber} combina las ventajas del \gls{MSE} y \gls{MAE}:
	\[
		\mathcal{L}_{\delta}(y, \hat{y}) = \begin{cases}
			\frac{1}{2}(y - \hat{y})^2 & \text{si } |y - \hat{y}| \leq \delta \\
			\delta |y - \hat{y}| - \frac{1}{2}\delta^2 & \text{en otro caso}
		\end{cases}
	\]
	donde \(\delta > 0\) es un hiperparámetro que controla la transición.
\end{definition}

La pérdida de Huber es cuadrática para errores pequeños (proporcionando gradientes suaves) y lineal para errores grandes (robustez a \textit{outliers}). Es diferenciable en todo punto, incluyendo la transición en \(|y - \hat{y}| = \delta\).

\paragraph{Clasificación}

Los problemas de \textit{clasificación} consisten en asignar cada entrada a una de \(K\) categorías discretas predefinidas. Matemáticamente, buscamos aproximar una función \(f: \mathcal{X} \to \{1, \ldots, K\}\), aunque en la práctica la red produce un vector de probabilidades \(\hat{y} \in \Delta^{K-1}\) (el símplex de dimensión \(K-1\)). Cuando \(K = 2\) hablamos de \textit{clasificación binaria} (por ejemplo, detectar correo \textit{spam}); cuando \(K > 2\), de \textit{clasificación multiclase} (por ejemplo, reconocer dígitos manuscritos del 0 al 9, o clasificar imágenes en categorías). Las funciones de pérdida para clasificación penalizan las discrepancias entre la distribución predicha y la clase real.

\begin{definition}[Entropía cruzada]\label{def:cross_entropy}
	La \textit{entropía cruzada} mide la divergencia entre la distribución real de clases \(y\) y la distribución estimada \(\hat{y}\):
	\begin{equation}
		\mathcal{L}(y, \hat{y}) = -\sum_{i=1}^{K} y_i \log(\hat{y}_i)
		\label{eq:cross_entropy}
	\end{equation}
	donde $K$ es el número de clases, \(y_i\) es la probabilidad real y \(\hat{y}_i\) la probabilidad estimada por la red.
\end{definition}

Para que los valores de la última capa sean probabilidades, se aplica la función softmax (\autoref{def:softmax}).

\begin{definition}[Pérdida logística binaria]\label{def:binary_cross_entropy}
	Para clasificación binaria con \(y \in \{0, 1\}\) y \(\hat{y} \in (0, 1)\):
	\[
		\mathcal{L}(y, \hat{y}) = -y \log(\hat{y}) - (1-y) \log(1 - \hat{y})
	\]
	Esta es la entropía cruzada especializada al caso binario, también conocida como \textit{log loss}.
\end{definition}

\begin{definition}[Pérdida hinge]\label{def:hinge_loss}
	La \textit{pérdida hinge}, usada en máquinas de vectores de soporte (SVM), para \(y \in \{-1, +1\}\):
	\[
		\mathcal{L}(y, \hat{y}) = \max(0, 1 - y \cdot \hat{y})
	\]
\end{definition}

La pérdida hinge no penaliza predicciones correctas con margen mayor que 1, lo que induce soluciones con margen máximo. No es diferenciable en \(y \cdot \hat{y} = 1\).

\paragraph{Para distribuciones de probabilidad}

Además de los problemas estándar de regresión y clasificación, algunas aplicaciones requieren que la red produzca una distribución de probabilidad completa como salida, no solo un punto o una clase. Esto es común en modelos generativos, donde la red aprende a generar datos que siguen una distribución objetivo, y en destilación de conocimiento, donde una red estudiante aprende a imitar las distribuciones de probabilidad de una red maestra. Las siguientes funciones de pérdida miden la divergencia o distancia entre distribuciones de probabilidad.

\begin{definition}[Divergencia de Kullback-Leibler]\label{def:kl_divergence}
	La \textit{divergencia KL} entre dos distribuciones \(P\) y \(Q\) sobre un espacio discreto \(\mathcal{X}\) es:
	\[
		D_{\text{KL}}(P \| Q) = \sum_{x \in \mathcal{X}} P(x) \log \frac{P(x)}{Q(x)}
	\]
	Para distribuciones continuas, la suma se reemplaza por una integral.
\end{definition}

La divergencia \gls{KL} mide cuánta información se pierde al usar \(Q\) para aproximar \(P\). Nótese que \(D_{\text{KL}}(P \| Q) \neq D_{\text{KL}}(Q \| P)\) en general, por lo que no es simétrica. La entropía cruzada se relaciona con \gls{KL} mediante: \(H(P, Q) = H(P) + D_{\text{KL}}(P \| Q)\), donde \(H(P)\) es la entropía de \(P\) y \(H(P, Q)\) la entropía cruzada.

La divergencia \gls{KL} es un caso especial de una familia más general de divergencias:

\begin{definition}[\(f\)-divergencia]\label{def:f_divergence}
	Sea \(f: (0, \infty) \to \mathbb{R}\) una función convexa con \(f(1) = 0\). La \textit{\(f\)-divergencia} entre distribuciones \(P\) y \(Q\) es:
	\[
		D_f(P \| Q) = \sum_{x \in \mathcal{X}} Q(x) \, f\left(\frac{P(x)}{Q(x)}\right)
	\]
	Para distribuciones continuas con densidades \(p\) y \(q\): \(D_f(P \| Q) = \int q(x) \, f\left(\frac{p(x)}{q(x)}\right) dx\).
\end{definition}

Las \(f\)-divergencias unifican múltiples medidas de discrepancia entre distribuciones:
\begin{itemize}
	\item \textbf{KL divergence}: \(f(t) = t \log t\) produce \(D_{\text{KL}}(P \| Q)\)
	\item \textbf{KL inversa}: \(f(t) = -\log t\) produce \(D_{\text{KL}}(Q \| P)\)
	\item \textbf{Jensen-Shannon}: \(f(t) = t \log t - (t+1) \log \frac{t+1}{2}\) produce una divergencia simétrica
	\item \textbf{Variación total}: \(f(t) = \frac{1}{2}|t - 1|\) produce \(\text{TV}(P, Q) = \frac{1}{2}\sum_x |P(x) - Q(x)|\)
	\item \textbf{Chi-cuadrado}: \(f(t) = (t-1)^2\) produce la divergencia \(\chi^2\)
\end{itemize}

La propiedad clave de las \(f\)-divergencias es que \(D_f(P \| Q) \geq 0\) con igualdad si y solo si \(P = Q\), lo que las hace adecuadas como funciones de pérdida para comparar distribuciones.

\paragraph{Para aprendizaje de representaciones}

El \textit{aprendizaje de representaciones} busca construir una función \(f: \mathcal{X} \to \mathbb{R}^d\) que mapee datos de entrada a un espacio de \textit{embeddings} donde la geometría refleje la semántica del problema. Un caso particular es el \textit{aprendizaje métrico} (\textit{metric learning}), donde el objetivo es que la distancia en el espacio de embeddings codifique similitud: elementos semánticamente similares (misma clase, mismo objeto, mismo hablante) deben estar cerca, mientras que elementos disimilares deben estar separados. Aplicaciones típicas incluyen reconocimiento facial (verificar si dos imágenes son de la misma persona), búsqueda de imágenes similares, y sistemas de recomendación. Las funciones de pérdida para aprendizaje métrico operan sobre pares o tripletes de ejemplos, no sobre ejemplos individuales.

\begin{definition}[Pérdida contrastiva]\label{def:contrastive_loss}
	La \textit{pérdida contrastiva}~\citep{hadsell2006dimensionality} para un par de embeddings \((f(x_1), f(x_2))\) con etiqueta \(y \in \{0, 1\}\) (0 = similares, 1 = disimilares):
	\[
		\mathcal{L}(y, x_1, x_2) = (1-y) \frac{1}{2} d^2 + y \frac{1}{2} \max(0, m - d)^2
	\]
	donde \(d = \|f(x_1) - f(x_2)\|_2\) es la distancia euclidiana y \(m > 0\) es el margen.
\end{definition}

La pérdida contrastiva penaliza pares similares que están lejos y pares disimilares que están demasiado cerca (dentro del margen \(m\)).

\begin{definition}[Pérdida triplet]\label{def:triplet_loss}
	La \textit{pérdida triplet}~\citep{schroff2015facenet} para un trío (anchor \(a\), positivo \(p\), negativo \(n\)):
	\[
		\mathcal{L}(a, p, n) = \max(0, \|f(a) - f(p)\|_2 - \|f(a) - f(n)\|_2 + m)
	\]
	donde \(m > 0\) es el margen que asegura separación entre clases.
\end{definition}

La pérdida triplet fuerza que el anchor esté más cerca del positivo que del negativo por un margen \(m\). Es ampliamente utilizada en reconocimiento facial y verificación de similitud.

\subsubsection{Funciones de pérdida vs. métricas}

Para comprender por qué muchas funciones de pérdida no son métricas, primero establecemos la definición formal de métrica.

\begin{definition}[Métrica]\label{def:metrica}
	Una \textit{métrica} (o \textit{distancia}) en un conjunto \(X\) es una función \(d: X \times X \to \mathbb{R}_{\geq 0}\) que satisface para todo \(x, y, z \in X\):
	\begin{enumerate}
		\item \textbf{No negatividad}: \(d(x, y) \geq 0\)
		\item \textbf{Identidad de los indiscernibles}: \(d(x, y) = 0 \iff x = y\)
		\item \textbf{Simetría}: \(d(x, y) = d(y, x)\)
		\item \textbf{Desigualdad triangular}: \(d(x, z) \leq d(x, y) + d(y, z)\)
	\end{enumerate}
\end{definition}

Es importante destacar que muchas funciones de pérdida comunes \textbf{no son métricas} en el sentido de la \autoref{def:metrica}. Por ejemplo:
\begin{itemize}
	\item La \textbf{entropía cruzada} no satisface la simetría: \(\mathcal{L}(y, \hat{y}) \neq \mathcal{L}(\hat{y}, y)\)
	\item El \textbf{error cuadrático medio} (MSE) sí define una métrica cuando se toma la raíz cuadrada
	\item La \textbf{divergencia Kullback-Leibler} no satisface ni simetría ni desigualdad triangular
\end{itemize}

Lo crucial para el aprendizaje automático no es que \(\mathcal{L}\) sea una métrica, sino que:
\begin{itemize}
	\item Proporcione gradientes informativos para la optimización
	\item Penalice más los errores grandes que los pequeños
	\item Sea computacionalmente eficiente
	\item Se alinee con el objetivo de la tarea (clasificación, regresión, etc.)
\end{itemize}

Tanto las métricas como las funciones de pérdida miden ``distancia'' entre valores, pero tienen propósitos fundamentalmente distintos:
\begin{itemize}
	\item \textbf{Métricas}: Definen una topología sobre el espacio, permitiendo hablar de convergencia, continuidad y completitud. Son esenciales en análisis matemático.
	\item \textbf{Funciones de pérdida}: Diseñadas para guiar la optimización basada en gradientes, priorizan la calidad de la señal de error sobre propiedades topológicas.
\end{itemize}

\subsection{Descenso del gradiente}\label{subsec:descenso_gradiente}

Una vez definida la función de pérdida como medida cuantitativa del error de la red, el siguiente paso es determinar cómo ajustar los parámetros para minimizarla. El método más extendido para entrenar redes neuronales es el aprendizaje supervisado mediante retropropagación del error\footnote{Existen también métodos de entrenamiento no supervisados y auto-supervisados que quedan fuera del alcance de este libro. No obstante, la idea subyacente es común: disponer de alguna función de pérdida que proporcione una señal sobre el error de la red. Por ejemplo, las \acrlong{GAN} (\gls{GAN})~\citep{goodfellow2014generative} emplean una función de pérdida adversarial en la que dos redes compiten entre sí.}. Este procedimiento puede resumirse en los siguientes pasos:

\begin{enumerate}
	\item \textbf{Inicialización de parámetros}: Los pesos y sesgos de la red se inicializan típicamente con valores aleatorios pequeños.
	\item \textbf{Propagación hacia adelante (forward pass)}: Para una entrada dada, la información se propaga desde la capa de entrada, pasando por las capas ocultas, hasta la capa de salida. Esto produce una predicción.
	\item \textbf{Cálculo de la pérdida}: Se compara la predicción con la salida verdadera y se calcula un valor de pérdida.
	\item \textbf{Retropropagación (backpropagation~\citep{backpropagation})}: Se calcula el gradiente de la función de pérdida con respecto a cada uno de los pesos de la red utilizando la regla de la cadena. Este gradiente indica la dirección en la que se deben ajustar los pesos para reducir la pérdida. Actualmente, se utilizan técnicas de diferenciación automática~\citep{Neidinger2010IntroductionTA} para calcular estos gradientes de manera eficiente.
	\item \textbf{Actualización de pesos}: Los pesos se ajustan utilizando métodos de optimización como el descenso por gradiente estocástico (\gls{SGD}) o variantes más avanzadas (\gls{Adam}, RMSProp, Adagrad). Este ajuste se realiza muchas veces, iterando sobre el conjunto de datos, hasta que la red converge a un conjunto de parámetros que minimizan la pérdida.
\end{enumerate}

El resultado de este proceso es una red que, idealmente, generaliza bien a datos no vistos. Esto significa que la red no sólo memoriza las muestras de entrenamiento, sino que aprende patrones subyacentes que pueden aplicarse a nuevas entradas.

En general no es factible aplicar una iteración del algoritmo anterior sobre todo el conjunto de datos de entrenamiento, sino que se realiza en lotes o batches. Se coge un subconjunto pequeño de los datos de entrenamiento y se aplica una iteración. Y así se va repitiendo hasta un que se alcanza un criterio de parada predefinido (un error suficientemente pequeño o un número de iteraciones suficientes). Este algoritmo se denomina \gls{SGD}.

\subsubsection{Formalización matemática}

Formalizamos ahora los algoritmos de optimización fundamentales~\citep{robbins1951stochastic, bottou2018optimization}.

\begin{definition}[Descenso del gradiente]\label{def:gradient_descent}
	Sea \(\mathcal{L}: \mathbb{R}^p \to \mathbb{R}\) una función de pérdida diferenciable. El \textit{descenso del gradiente} es el algoritmo iterativo definido por:
	\[
		\theta_{t+1} = \theta_t - \eta \nabla_\theta \mathcal{L}(\theta_t)
	\]
	donde \(\eta > 0\) es la \textit{tasa de aprendizaje} (o \textit{learning rate}) y \(\nabla_\theta \mathcal{L}(\theta_t)\) es el gradiente de la pérdida evaluado en \(\theta_t\).
\end{definition}

El algoritmo de descenso del gradiente sigue la dirección de máximo decrecimiento de la función de pérdida. Sin embargo, calcular el gradiente exacto sobre todos los datos de entrenamiento es computacionalmente costoso para conjuntos grandes.

La \autoref{fig:loss_landscapes} ilustra la diferencia fundamental entre superficies de pérdida convexas (con un único mínimo global) y no convexas (típicas en deep learning, con múltiples mínimos locales). Esta distinción tiene implicaciones cruciales para la convergencia de los algoritmos de optimización.

\begin{figure}[h]
\centering
\includegraphics[width=0.95\textwidth]{figures/loss_landscapes}
\caption[Superficies de pérdida convexa vs no convexa]{Comparación de superficies de pérdida. Izquierda: Función convexa con un único mínimo global donde cualquier trayectoria de descenso del gradiente converge al óptimo. Derecha: Función no convexa típica de redes neuronales profundas, con múltiples mínimos locales. Estudios recientes sugieren que en redes sobreparametrizadas, la mayoría de mínimos locales tienen valores de pérdida similares al global.}
\label{fig:loss_landscapes}
\end{figure}

\begin{definition}[Descenso del gradiente estocástico (SGD)]\label{def:sgd}
	Sea \(\{(x_i, y_i)\}_{i=1}^N\) un conjunto de datos y \(\mathcal{L}(\theta) = \frac{1}{N}\sum_{i=1}^N \ell(f_\theta(x_i), y_i)\) la pérdida empírica. El \textit{descenso del gradiente estocástico} aproxima el gradiente usando un subconjunto aleatorio (\textit{mini-batch}) \(\mathcal{B} \subset \{1,\ldots,N\}\):
\[
	\theta_{t+1} = \theta_t - \eta \cdot \frac{1}{|\mathcal{B}|}\sum_{i \in \mathcal{B}} \nabla_\theta \ell(f_\theta(x_i), y_i)
\]
donde \(|\mathcal{B}|\) es el tamaño del mini-batch.
\end{definition}

El gradiente estimado por \gls{SGD} es un estimador \textit{no sesgado} del gradiente real:
\[
	\mathbb{E}_{\mathcal{B}}\left[\frac{1}{|\mathcal{B}|}\sum_{i \in \mathcal{B}} \nabla_\theta \ell(f_\theta(x_i), y_i)\right] = \nabla_\theta \mathcal{L}(\theta)
\]
Aunque tiene varianza mayor que el gradiente exacto, esta varianza puede actuar como regularización implícita, ayudando a escapar de mínimos locales subóptimos.

\subsubsection{Retropropagación (Backpropagation)}\label{subsubsec:backpropagation}

El algoritmo de retropropagación~\citep{backpropagation} es el método estándar para calcular los gradientes necesarios en el descenso del gradiente. Su eficiencia radica en explotar la estructura de grafo de la red neuronal.

\begin{definition}[Grafo computacional]\label{def:computational_graph}
	Un \textit{grafo computacional} es un grafo acíclico dirigido \(G = (V, E)\) donde:
	\begin{itemize}
		\item Cada nodo \(v \in V\) representa una variable o una operación diferenciable \(f_v\)
		\item Los nodos de entrada (sin aristas entrantes) corresponden a parámetros \(\theta\) y datos de entrada \(x\)
		\item Cada arista \((u, v) \in E\) indica que \(v\) depende de \(u\)
		\item El nodo de salida (único, sin aristas salientes) produce el valor de la pérdida \(\mathcal{L}\)
	\end{itemize}
\end{definition}

El grafo computacional permite descomponer el cálculo del gradiente en operaciones locales mediante la regla de la cadena, como se ilustra en la \autoref{fig:computational_graph}.

\begin{figure}[h]
\centering
\begin{tikzpicture}[
    node distance=1.2cm,
    inputnode/.style={circle, draw, thick, minimum size=0.8cm, fill=gdlblue!20},
    opnode/.style={circle, draw, thick, minimum size=0.8cm, fill=gdlgreen!20},
    lossnode/.style={circle, draw, thick, minimum size=0.8cm, fill=gdlred!20},
    arrow/.style={->, thick, >=stealth}
]

% Nodos de entrada
\node[inputnode] (x) at (0,0) {$x$};
\node[inputnode] (w) at (0,1.5) {$w$};
\node[inputnode] (b) at (0,-1.5) {$b$};

% Operaciones intermedias
\node[opnode] (mult) at (2,0.75) {$\times$};
\node[opnode] (add) at (4,0) {$+$};
\node[opnode] (sigma) at (6,0) {$\sigma$};
\node[opnode] (diff) at (8,0) {$-$};
\node[opnode] (sq) at (10,0) {$(\cdot)^2$};

% Nodo de pérdida
\node[lossnode] (L) at (12,0) {$\mathcal{L}$};

% Target
\node[inputnode] (y) at (8,1.5) {$y$};

% Conexiones
\draw[arrow] (x) -- (mult);
\draw[arrow] (w) -- (mult);
\draw[arrow] (mult) -- node[above, font=\scriptsize] {$wx$} (add);
\draw[arrow] (b) -- (add);
\draw[arrow] (add) -- node[above, font=\scriptsize] {$z$} (sigma);
\draw[arrow] (sigma) -- node[above, font=\scriptsize] {$\hat{y}$} (diff);
\draw[arrow] (y) -- (diff);
\draw[arrow] (diff) -- node[above, font=\scriptsize] {$e$} (sq);
\draw[arrow] (sq) -- node[above, font=\scriptsize] {$\frac{1}{2}$} (L);

% Etiquetas de tipos de nodos
\node[font=\footnotesize, gdlgray] at (0,-2.3) {Entradas};
\node[font=\footnotesize, gdlgray] at (6,-1.5) {Operaciones};
\node[font=\footnotesize, gdlgray] at (12,-1) {Pérdida};

% Forward y Backward
\draw[dashed, gdlgray, ->] (0.5,-2.8) -- (11.5,-2.8) node[midway, below, font=\scriptsize] {Forward pass};
\draw[dashed, gdlgray, <-] (0.5,-3.3) -- (11.5,-3.3) node[midway, below, font=\scriptsize] {Backward pass};

\end{tikzpicture}
\caption[Grafo computacional para una neurona]{Grafo computacional para una neurona simple con pérdida MSE. Los nodos azules representan entradas (parámetros $w$, $b$ y dato $x$), los verdes representan operaciones diferenciables ($\times$, $+$, $\sigma$, $-$, $(\cdot)^2$), y el rojo representa la pérdida final $\mathcal{L}$. Las aristas indican dependencias: $z = wx + b$, $\hat{y} = \sigma(z)$, $\mathcal{L} = \frac{1}{2}(y - \hat{y})^2$.}
\label{fig:computational_graph}
\end{figure}

Los tres grafos definidos en esta sección representan la misma red a diferentes niveles de abstracción:
\begin{itemize}
	\item El \textbf{grafo de capas} (\autoref{def:grafo_de_capas}) es la vista macroscópica: cada nodo es una capa completa ($L_1, L_2, \ldots$)
	\item El \textbf{grafo de neuronas} $\mathcal{G}_{\text{neuronas}}$ es la vista mesoscópica: cada nodo es una neurona individual con sus conexiones ponderadas
	\item El \textbf{grafo computacional} (\autoref{def:computational_graph}) es la vista microscópica: cada nodo es una operación elemental ($\times$, $+$, $\sigma$, etc.)
\end{itemize}
El grafo computacional ``despliega'' cada neurona del grafo de neuronas en sus operaciones constituyentes (multiplicación por peso, suma del sesgo, aplicación de la activación), permitiendo aplicar la regla de la cadena de forma sistemática durante la retropropagación.

\begin{theorem}[Regla de la cadena en grafos computacionales]\label{thm:backprop}
	Sea \(G = (V, E)\) un grafo computacional con nodo de salida \(\mathcal{L}\). Para cada nodo intermedio \(v\) con sucesores \(\{u_1, \ldots, u_k\}\) en el grafo (es decir, nodos que dependen de \(v\)):
	\[
		\frac{\partial \mathcal{L}}{\partial v} = \sum_{i=1}^{k} \frac{\partial \mathcal{L}}{\partial u_i} \cdot \frac{\partial u_i}{\partial v}
	\]
	El algoritmo de backpropagation calcula todos los gradientes \(\frac{\partial \mathcal{L}}{\partial v}\) para \(v \in V\) en tiempo \(O(|V| + |E|)\), es decir, lineal en el tamaño del grafo.
\end{theorem}

\begin{proof}
	Demostramos la corrección del algoritmo y analizamos su complejidad.

	\textbf{Existencia de ordenación topológica.}
	Como \(G = (V, E)\) es un grafo acíclico dirigido (DAG), existe una ordenación topológica \(v_1, v_2, \ldots, v_n\) de los nodos tal que si \((v_i, v_j) \in E\), entonces \(i < j\). Esta ordenación se puede calcular en tiempo \(O(|V| + |E|)\) mediante algoritmos estándar (por ejemplo, DFS o eliminación iterativa de fuentes).

	\textbf{Forward pass.}
	Recorremos los nodos en orden topológico. Para cada nodo \(v_i\):
	\begin{itemize}
		\item Si \(v_i\) es un nodo de entrada (parámetro o dato), su valor ya está definido.
		\item Si \(v_i\) es un nodo de operación con predecesores \(\{u_1, \ldots, u_k\}\), calculamos \(v_i = f_{v_i}(u_1, \ldots, u_k)\).
	\end{itemize}
	Por la propiedad de la ordenación topológica, todos los predecesores de \(v_i\) tienen índices menores y ya han sido evaluados. El coste de evaluar cada nodo es \(O(\text{grado de entrada})\), sumando en total \(O(|V| + |E|)\).

	\textbf{Backward pass: Corrección.}
	Para demostrar la fórmula de la regla de la cadena, sea \(v\) un nodo intermedio y sean \(u_1, \ldots, u_k\) sus sucesores directos (nodos que dependen de \(v\)). La pérdida \(\mathcal{L}\) depende de \(v\) únicamente a través de estos sucesores. Por la regla de la cadena multivariable:
	\[
		\frac{\partial \mathcal{L}}{\partial v} = \sum_{i=1}^{k} \frac{\partial \mathcal{L}}{\partial u_i} \cdot \frac{\partial u_i}{\partial v}
	\]
	donde \(\frac{\partial u_i}{\partial v}\) es la derivada parcial local de la operación en \(u_i\) respecto a su entrada \(v\).

	\textbf{Backward pass: Algoritmo.}
	Recorremos los nodos en orden topológico inverso \(v_n, v_{n-1}, \ldots, v_1\):
	\begin{enumerate}
		\item \textbf{Inicialización}: Para el nodo de salida \(\mathcal{L}\), establecemos \(\frac{\partial \mathcal{L}}{\partial \mathcal{L}} = 1\).
		\item \textbf{Propagación}: Para cada nodo \(v_i\), ya tenemos calculados los gradientes \(\frac{\partial \mathcal{L}}{\partial u_j}\) de todos sus sucesores \(u_j\) (tienen índices mayores). Calculamos:
		\[
			\frac{\partial \mathcal{L}}{\partial v_i} = \sum_{(v_i, u_j) \in E} \frac{\partial \mathcal{L}}{\partial u_j} \cdot \frac{\partial u_j}{\partial v_i}
		\]
	\end{enumerate}

	\textbf{Análisis de complejidad.}
	En el backward pass, cada arista \((v_i, u_j) \in E\) contribuye exactamente una vez a la suma del gradiente de \(v_i\). El coste de calcular cada derivada parcial local \(\frac{\partial u_j}{\partial v_i}\) es \(O(1)\) para operaciones elementales. Por tanto:
	\begin{itemize}
		\item Cada nodo se visita una vez: contribución \(O(|V|)\)
		\item Cada arista se procesa una vez: contribución \(O(|E|)\)
	\end{itemize}

	La complejidad total del algoritmo completo (forward + backward) es \(O(|V| + |E|)\).
\end{proof}

\paragraph{Diferenciación automática}
En la práctica, los frameworks modernos de deep learning (PyTorch, TensorFlow, JAX) implementan diferenciación automática~\citep{Neidinger2010IntroductionTA} de forma transparente: durante el forward pass, construyen dinámicamente el grafo computacional registrando cada operación; durante el backward pass, recorren este grafo en orden topológico inverso aplicando la regla de la cadena. Esto libera al usuario de calcular derivadas manualmente, permitiéndole concentrarse en el diseño de arquitecturas.

\subsubsection{Ejemplo numérico: XOR}

Para ilustrar concretamente el proceso de aprendizaje, consideremos el problema XOR con una red neuronal mínima de arquitectura 2-2-1 (2 entradas, 2 neuronas ocultas, 1 salida).

\textbf{Datos de entrenamiento:}
\begin{center}
	\begin{tabular}{cc|c}
		\(x_1\) & \(x_2\) & \(y\) \\
		\hline
		0       & 0       & 0     \\
		0       & 1       & 1     \\
		1       & 0       & 1     \\
		1       & 1       & 0     \\
	\end{tabular}
\end{center}

\textbf{Arquitectura:}
\begin{itemize}
	\item Capa de entrada: \(\mathbf{x} = [x_1, x_2]^T\)
	\item Capa oculta: \(\mathbf{h} = \sigma(W^{(1)}\mathbf{x} + \mathbf{b}^{(1)})\) con \(W^{(1)} \in \mathbb{R}^{2 \times 2}\), \(\mathbf{b}^{(1)} \in \mathbb{R}^2\)
	\item Capa de salida: \(\hat{y} = \sigma(W^{(2)}\mathbf{h} + b^{(2)})\) con \(W^{(2)} \in \mathbb{R}^{1 \times 2}\), \(b^{(2)} \in \mathbb{R}\)
	\item Función de activación: \(\sigma(z) = \frac{1}{1 + e^{-z}}\) (sigmoide logística, véase la \autoref{tab:funciones_activacion})
\end{itemize}

\textbf{Ejemplo de propagación hacia adelante:}

Supongamos pesos iniciales:
\begin{gather*}
    W^{(1)} = \begin{pmatrix}
		          0.5 & -0.3 \\
		          0.2 & 0.4
	\end{pmatrix}, \quad
	\mathbf{b}^{(1)} = \begin{pmatrix}
		                   -0.1 \\
		                   0.3
	\end{pmatrix}\\
    W^{(2)} = \begin{pmatrix}
		          0.7 & -0.5
	\end{pmatrix}, \quad
	b^{(2)} = 0.2\\
\end{gather*}

Para la entrada \(\mathbf{x} = [1, 0]^T\):

1. Capa oculta:
\begin{gather*}
    \mathbf{z}^{(1)} = W^{(1)}\mathbf{x} + \mathbf{b}^{(1)} = \begin{pmatrix}
		                                                          0.5 \cdot 1 + (-0.3) \cdot 0 - 0.1 \\
		                                                          0.2 \cdot 1 + 0.4 \cdot 0 + 0.3
	\end{pmatrix} = \begin{pmatrix}
		                0.4 \\
		                0.5
	\end{pmatrix}\\
    \mathbf{h} = \sigma(\mathbf{z}^{(1)}) = \begin{pmatrix}
		                                        \sigma(0.4) \\
		                                        \sigma(0.5)
	\end{pmatrix} = \begin{pmatrix}
		                0.599 \\
		                0.622
	\end{pmatrix}\\
\end{gather*}

2. Capa de salida:
\begin{gather*}
    z^{(2)} = W^{(2)}\mathbf{h} + b^{(2)} = 0.7 \cdot 0.599 + (-0.5) \cdot 0.622 + 0.2 = 0.308\\
    \hat{y} = \sigma(0.308) = 0.576\\
\end{gather*}

3. Error (usando \gls{MSE}):
\[
	\mathcal{L} = \frac{1}{2}(y - \hat{y})^2 = \frac{1}{2}(1 - 0.576)^2 = 0.090
\]

\textbf{Retropropagación:}

Para calcular los gradientes de manera sistemática, definimos la \textit{señal de error} (o delta) de cada capa. Para la capa de salida:
\[
	\delta^{(2)} = \frac{\partial \mathcal{L}}{\partial z^{(2)}} = -(y - \hat{y}) \cdot \sigma'(z^{(2)}) = -(y - \hat{y}) \cdot \hat{y}(1 - \hat{y}) \in \mathbb{R}
\]
donde usamos que la derivada de la sigmoide es \(\sigma'(z) = \sigma(z)(1-\sigma(z))\).

El gradiente respecto a los pesos de la capa de salida es entonces:
\[
	\frac{\partial \mathcal{L}}{\partial W^{(2)}} = \delta^{(2)} \cdot \mathbf{h}^T \in \mathbb{R}^{1 \times 2}
\]

Para la capa oculta, la señal de error se propaga hacia atrás:
\[
	\boldsymbol{\delta}^{(1)} = \left( W^{(2)T} \cdot \delta^{(2)} \right) \odot \sigma'(\mathbf{z}^{(1)}) = \left( W^{(2)T} \cdot \delta^{(2)} \right) \odot \left[ \mathbf{h} \odot (1 - \mathbf{h}) \right] \in \mathbb{R}^{2}
\]
donde \(\odot\) denota el producto de Hadamard (elemento a elemento).

El gradiente respecto a los pesos de la capa oculta es:
\[
	\frac{\partial \mathcal{L}}{\partial W^{(1)}} = \boldsymbol{\delta}^{(1)} \cdot \mathbf{x}^T \in \mathbb{R}^{2 \times 2}
\]
Este es un producto exterior entre el vector de señales de error y el vector de entrada.

\textbf{Cálculos numéricos:}

Calculamos los valores intermedios paso a paso. Primero, el error de señal de la capa de salida:
\begin{gather*}
	\delta^{(2)} = -(y - \hat{y}) \cdot \hat{y}(1 - \hat{y}) = -(1 - 0.576) \cdot 0.576 \cdot (1 - 0.576) = -0.424 \cdot 0.244 = -0.103
\end{gather*}

Para los gradientes de la capa de salida:
\begin{gather*}
	\frac{\partial \mathcal{L}}{\partial W^{(2)}} = \delta^{(2)} \cdot \mathbf{h}^T = -0.103 \cdot \begin{pmatrix} 0.599 & 0.622 \end{pmatrix} = \begin{pmatrix} -0.062 & -0.064 \end{pmatrix} \\[0.3em]
	\frac{\partial \mathcal{L}}{\partial b^{(2)}} = \delta^{(2)} = -0.103
\end{gather*}

Para la capa oculta, propagamos el error hacia atrás:
\begin{gather*}
	\boldsymbol{\delta}^{(1)} = \delta^{(2)} \cdot W^{(2)T} \odot \mathbf{h} \odot (1 - \mathbf{h}) \\
	= -0.103 \cdot \begin{pmatrix} 0.7 \\ -0.5 \end{pmatrix} \odot \begin{pmatrix} 0.599 \cdot 0.401 \\ 0.622 \cdot 0.378 \end{pmatrix}
	= \begin{pmatrix} -0.072 \\ 0.052 \end{pmatrix} \odot \begin{pmatrix} 0.240 \\ 0.235 \end{pmatrix}
	= \begin{pmatrix} -0.017 \\ 0.012 \end{pmatrix}
\end{gather*}

\begin{gather*}
	\frac{\partial \mathcal{L}}{\partial W^{(1)}} = \boldsymbol{\delta}^{(1)} \cdot \mathbf{x}^T = \begin{pmatrix} -0.017 \\ 0.012 \end{pmatrix} \cdot \begin{pmatrix} 1 & 0 \end{pmatrix} = \begin{pmatrix} -0.017 & 0 \\ 0.012 & 0 \end{pmatrix}
\end{gather*}

\textbf{Actualización de pesos:}

Con tasa de aprendizaje \(\eta = 0.5\), actualizamos los pesos de la capa de salida:
\begin{gather*}
	W^{(2)}_{\text{nuevo}} = W^{(2)} - \eta \cdot \frac{\partial \mathcal{L}}{\partial W^{(2)}} = \begin{pmatrix} 0.7 & -0.5 \end{pmatrix} - 0.5 \cdot \begin{pmatrix} -0.062 & -0.064 \end{pmatrix} = \begin{pmatrix} 0.731 & -0.468 \end{pmatrix} \\[0.3em]
	b^{(2)}_{\text{nuevo}} = b^{(2)} - \eta \cdot \frac{\partial \mathcal{L}}{\partial b^{(2)}} = 0.2 - 0.5 \cdot (-0.103) = 0.252
\end{gather*}

\textbf{Verificación:}

Realizamos un nuevo forward pass con los pesos actualizados para verificar que la pérdida ha disminuido:
\begin{gather*}
	z^{(2)}_{\text{nuevo}} = W^{(2)}_{\text{nuevo}} \cdot \mathbf{h} + b^{(2)}_{\text{nuevo}} = 0.731 \cdot 0.599 + (-0.468) \cdot 0.622 + 0.252 = 0.399 \\
	\hat{y}_{\text{nuevo}} = \sigma(0.399) = 0.598 \\
	\mathcal{L}_{\text{nuevo}} = \frac{1}{2}(1 - 0.598)^2 = 0.081
\end{gather*}

Efectivamente, \(\mathcal{L}_{\text{nuevo}} = 0.081 < 0.090 = \mathcal{L}_{\text{anterior}}\), confirmando que el descenso del gradiente reduce la pérdida. Este proceso se repite iterativamente hasta que la red aprende a clasificar correctamente todas las entradas del problema XOR.

\begin{figure}[h]
	\centering
	\begin{tikzpicture}[
		node distance=2cm,
		neuron/.style={circle, draw, thick, minimum size=1cm},
		input/.style={circle, draw, thick, minimum size=0.9cm, fill=gdlblue!20},
		hidden/.style={circle, draw, thick, minimum size=1.0cm, fill=gdlgreen!20},
		output/.style={circle, draw, thick, minimum size=1.0cm, fill=gdlred!20},
		bias/.style={circle, draw, thick, minimum size=0.7cm, fill=gdlgray!20},
		weight/.style={midway, above, sloped, font=\small},
		activation/.style={below=0.5cm, font=\small}
	]

% Input layer
		\node[input] (x1) at (0,2) {\(x_1\)};
		\node[input] (x2) at (0,0) {\(x_2\)};
		\node[bias] (b0) at (2,-2) {\scriptsize\(+1\)};

% Hidden layer (separada más de la entrada)
		\node[hidden] (h1) at (4,2) {\(h_1\)};
		\node[hidden] (h2) at (4,0) {\(h_2\)};
		\node[bias] (b1) at (6,-2) {\scriptsize\(+1\)};

% Output layer
		\node[output] (y) at (7.5,0) {\(\hat{y}\)};

% Connections with weights
		\draw[->, thick] (x1) -- (h1) node[weight, pos=0.2] {\(0.5\)};
		\draw[->, thick] (x1) -- (h2) node[weight, pos=0.2] {\(0.2\)};
		\draw[->, thick] (x2) -- (h1) node[weight, pos=0.2] {\(-0.3\)};
		\draw[->, thick] (x2) -- (h2) node[weight, pos=0.2] {\(0.4\)};
		\draw[->, thick] (h1) -- (y) node[weight, pos=0.2] {\(0.7\)};
		\draw[->, thick] (h2) -- (y) node[weight, pos=0.2] {\(-0.5\)};

% Bias connections (pesos al inicio de la flecha)
		\draw[->, dashed] (b0) -- (h1) node[weight, pos=0.2] {\(-0.1\)};
		\draw[->, dashed] (b0) -- (h2) node[weight, pos=0.2] {\(0.3\)};
		\draw[->, dashed] (b1) -- (y) node[weight, pos=0.2] {\(0.2\)};

% Layer labels
		\node[above=0.3cm of x1] {Entrada};
		\node[above=0.3cm of h1] {Capa oculta};
		\node[above=0.3cm of y] {Salida};

% Activation values for example x=[1,0]
		\node[activation] at (x1) {\(1\)};
		\node[activation] at (x2) {\(0\)};
		\node[activation] at (h1) {\(0.599\)};
		\node[activation] at (h2) {\(0.622\)};
		\node[activation] at (y) {\(0.576\)};

	\end{tikzpicture}
	\caption[Red neuronal para el problema XOR]{Red neuronal para el problema XOR con arquitectura 2-2-1. Se muestran los pesos sinápticos del ejemplo y las activaciones para la entrada \([1,0]^T\). Los nodos de sesgo (+1) se muestran con conexiones discontinuas.}
	\label{fig:xor_network}
\end{figure}

\subsubsection{El problema del desvanecimiento del gradiente}\label{subsubsec:vanishing_gradient}

Una consecuencia directa de la regla de la cadena es el \textit{problema del desvanecimiento del gradiente} (\textit{vanishing gradient problem}). Según el \autoref{thm:backprop}, el gradiente de la pérdida respecto a los parámetros de las primeras capas es proporcional al producto de las derivadas de las funciones de activación a lo largo de toda la red:
\[
\frac{\partial \mathcal{L}}{\partial \theta^{(1)}} \propto \prod_{\ell=1}^{L} \sigma'(z^{(\ell)})
\]

Si cada derivada $|\sigma'(z^{(\ell)})|$ es menor que 1 (como ocurre con funciones saturantes como la sigmoide, donde $|\sigma'(z)| \leq 0.25$), este producto decrece exponencialmente con la profundidad $L$. Como resultado, los gradientes de las primeras capas se vuelven numéricamente cero, impidiendo su actualización durante el entrenamiento.

Este fenómeno fue uno de los principales obstáculos para entrenar redes profundas en la década de 1990. La solución vino con el uso de funciones de activación no saturantes como \gls{ReLU}, donde $\sigma'(z) = 1$ para $z > 0$, permitiendo que los gradientes fluyan sin atenuación (véase la \autoref{prop:gradiente_no_saturante}).

\begin{figure}[htbp]
\centering
\begin{tikzpicture}[
    layer/.style={rectangle, draw, thick, minimum width=1cm, minimum height=0.6cm, font=\small},
    gradbar/.style={rectangle, minimum height=0.25cm},
    arrow/.style={->, thick, >=stealth}
]

% === PANEL IZQUIERDO: Sigmoide ===
\begin{scope}[xshift=-4cm]
    % Título
    \node[font=\bfseries] at (0.5, 2.8) {Sigmoide (saturante)};

    % Capas con gradiente decreciente (numeración invertida: L_5 arriba, L_1 abajo)
    \foreach \i/\layer/\grad/\col in {1/5/1.0/100, 2/4/0.65/75, 3/3/0.35/50, 4/2/0.15/30, 5/1/0.05/15} {
        \node[layer, fill=gdlred!\col!white] (s\i) at (0.5, 2-\i*0.9) {$L_\layer$};
        % Barra de gradiente a la izquierda
        \fill[gdlblue!80] (-1.2, 1.85-\i*0.9) rectangle (-1.2+\grad, 2.15-\i*0.9);
        \draw[gdlgray, thin] (-1.2, 1.85-\i*0.9) rectangle (-0.2, 2.15-\i*0.9);
    }

    % Etiqueta eje gradiente
    \node[font=\scriptsize, rotate=90] at (-1.6, 0) {Magnitud $|\nabla|$};

    % Flecha backprop (de arriba hacia abajo: desde la salida hacia la entrada)
    \draw[arrow, color=gdlgray!70, line width=1.5pt] (1.3, 1.6) -- (1.3, -2.3);
    \node[font=\scriptsize, color=gdlgray!70, rotate=90] at (1.7, -0.3) {backprop};

    % Etiqueta resultado
    \node[font=\small, align=center, color=gdlred!70!black] at (0.5, -3.25) {Gradiente $\to 0$};
\end{scope}

% === PANEL DERECHO: ReLU ===
\begin{scope}[xshift=4cm]
    % Título
    \node[font=\bfseries] at (0.5, 2.8) {ReLU (no saturante)};

    % Capas con gradiente estable (numeración invertida: L_5 arriba, L_1 abajo)
    \foreach \i/\layer/\grad in {1/5/1.0, 2/4/0.95, 3/3/0.90, 4/2/0.85, 5/1/0.80} {
        \node[layer, fill=gdlgreen!50!white] (r\i) at (0.5, 2-\i*0.9) {$L_\layer$};
        % Barra de gradiente a la izquierda
        \fill[gdlblue!80] (-1.2, 1.85-\i*0.9) rectangle (-1.2+\grad, 2.15-\i*0.9);
        \draw[gdlgray, thin] (-1.2, 1.85-\i*0.9) rectangle (-0.2, 2.15-\i*0.9);
    }

    % Etiqueta eje gradiente
    \node[font=\scriptsize, rotate=90] at (-1.6, 0) {Magnitud $|\nabla|$};

    % Flecha backprop (de arriba hacia abajo: desde la salida hacia la entrada)
    \draw[arrow, color=gdlgray!70, line width=1.5pt] (1.3, 1.6) -- (1.3, -2.3);
    \node[font=\scriptsize, color=gdlgray!70, rotate=90] at (1.7, -0.3) {backprop};

    % Etiqueta resultado
    \node[font=\small, align=center, color=gdlgreen!50!black] at (0.5, -3.25) {Gradiente estable};
\end{scope}

% === ECUACIÓN CENTRAL ===
\node[font=\small, align=center, draw=gdlgray!30, rounded corners, fill=gdlgray!5, inner sep=6pt] at (0, -1.2) {
    $\displaystyle\frac{\partial \mathcal{L}}{\partial \theta_1} \propto \prod_{\ell=1}^{L} \sigma'(z^{(\ell)})$
};

\end{tikzpicture}
\caption[Problema del desvanecimiento del gradiente]{Comparación del flujo de gradientes durante la retropropagación. \textbf{Izquierda}: Con funciones saturantes como la sigmoide, donde $|\sigma'(z)| \leq 0.25$, el producto de derivadas decrece exponencialmente. \textbf{Derecha}: Con ReLU, donde $\sigma'(z) = 1$ para $z > 0$, el gradiente se mantiene estable permitiendo entrenar redes más profundas.}
\label{fig:vanishing_gradient}
\end{figure}

\subsection{Teorema de aproximación universal}\label{sec:aproximacion_universal}

El teorema de aproximación universal garantiza que cualquier función continua se puede aproximar arbitrariamente bien por una \gls{ANN} con una única capa oculta. Su demostración requiere herramientas de teoría de la medida y de análisis funcional: los conceptos de medida y $\sigma$-álgebra fundamentan los teoremas de representación (Riesz) y separación (Hahn-Banach) que constituyen el núcleo de la prueba. Presentamos primero los preliminares necesarios y después el enunciado con su demostración completa. Para un tratamiento en profundidad de estos fundamentos, véase \citet{Folland1999} o \citet{Rudin1976}.

\paragraph{Preliminares de teoría de la medida}

\begin{definition}[$\sigma$-álgebra]\label{def:sigma_algebra}
	Sea $X$ un conjunto no vacío. Una \textit{$\sigma$-álgebra} (o $\sigma$-campo) sobre $X$ es una colección $\mathcal{F}$ de subconjuntos de $X$ que satisface:
	\begin{enumerate}
		\item $X \in \mathcal{F}$
		\item Si $A \in \mathcal{F}$, entonces $A^c = X \setminus A \in \mathcal{F}$ (cerrada bajo complementos)
		\item Si $\{A_n\}_{n=1}^\infty \subseteq \mathcal{F}$, entonces $\bigcup_{n=1}^\infty A_n \in \mathcal{F}$ (cerrada bajo uniones numerables)
	\end{enumerate}
	Los elementos de $\mathcal{F}$ se denominan \textit{conjuntos medibles}.
\end{definition}

\begin{definition}[Espacio medible]\label{def:espacio_medible}
	Un \textit{espacio medible} es un par $(X, \mathcal{F})$ donde $X$ es un conjunto no vacío y $\mathcal{F}$ es una $\sigma$-álgebra sobre $X$.
\end{definition}

\begin{definition}[$\sigma$-álgebra de Borel]\label{def:sigma_algebra_borel}
	Sea $(X, \tau)$ un espacio topológico. La \textit{$\sigma$-álgebra de Borel} $\mathcal{B}(X)$ es la $\sigma$-álgebra más pequeña que contiene todos los conjuntos abiertos de $X$. Sus elementos se denominan \textit{conjuntos de Borel} o \textit{borelianos}.
\end{definition}

En particular, para $\mathbb{R}^n$, la $\sigma$-álgebra de Borel $\mathcal{B}(\mathbb{R}^n)$ es generada por los rectángulos abiertos, o equivalentemente, por los semiespacios de la forma $\{x \in \mathbb{R}^n : w^T x + b > 0\}$.

\begin{definition}[Medida]\label{def:medida}
	Sea $(X, \mathcal{F})$ un espacio medible. Una \textit{medida} es una función $\mu: \mathcal{F} \to [0, +\infty]$ que satisface:
	\begin{enumerate}
		\item $\mu(\emptyset) = 0$
		\item \textbf{$\sigma$-aditividad}: Para cualquier sucesión $\{A_n\}_{n=1}^\infty$ de conjuntos disjuntos en $\mathcal{F}$:
		\[
			\mu\left(\bigcup_{n=1}^\infty A_n\right) = \sum_{n=1}^\infty \mu(A_n)
		\]
	\end{enumerate}
\end{definition}

\begin{definition}[Espacio de medida]\label{def:espacio_medida}
	Un \textit{espacio de medida} es una terna $(X, \mathcal{F}, \mu)$ donde $X$ es un conjunto no vacío, $\mathcal{F}$ es una $\sigma$-álgebra sobre $X$, y $\mu: \mathcal{F} \to [0, +\infty]$ es una medida. Si $\mu(X) < \infty$, decimos que $\mu$ es una \textit{medida finita}. Si $\mu(X) = 1$, decimos que $\mu$ es una \textit{medida de probabilidad}.
\end{definition}

\begin{definition}[Medida finita con signo]\label{def:medida_con_signo}
	Sea $(X, \mathcal{F})$ un espacio medible. Una \textit{medida finita con signo} es una función $\mu: \mathcal{F} \to \mathbb{R}$ que satisface $\mu(\emptyset) = 0$ y es $\sigma$-aditiva: para cualquier sucesión $\{A_n\}_{n=1}^\infty$ de conjuntos disjuntos en $\mathcal{F}$,
	\[
		\mu\left(\bigcup_{n=1}^\infty A_n\right) = \sum_{n=1}^\infty \mu(A_n)
	\]
	donde la serie converge absolutamente. La \textit{variación total} de $\mu$ es la medida $|\mu|$ definida por
	\[
		|\mu|(A) = \sup \left\{ \sum_{i=1}^n |\mu(A_i)| : \{A_i\}_{i=1}^n \text{ partición finita de } A \right\}
	\]
\end{definition}

\begin{definition}[Medida de Lebesgue]\label{def:medida_lebesgue}
	La \textit{medida de Lebesgue} $\lambda$ en $\mathbb{R}^n$ es la única medida en $(\mathbb{R}^n, \mathcal{B}(\mathbb{R}^n))$ que extiende la noción de volumen $n$-dimensional y satisface:
	\begin{itemize}
		\item Para un rectángulo $R = [a_1, b_1] \times \cdots \times [a_n, b_n]$: $\lambda(R) = \prod_{i=1}^n (b_i - a_i)$
		\item Es invariante bajo traslaciones: $\lambda(A + x) = \lambda(A)$ para todo $x \in \mathbb{R}^n$
	\end{itemize}
\end{definition}

Un conjunto $A$ tiene \textit{medida de Lebesgue nula} (o medida cero) si $\lambda(A) = 0$. Ejemplos incluyen: puntos aislados, conjuntos finitos, conjuntos numerables como $\mathbb{Q}$, y el conjunto de Cantor.

\begin{definition}[Función medible]\label{def:funcion_medible}
	Sean $(X, \mathcal{F})$ y $(Y, \mathcal{G})$ espacios medibles. Una función $f: X \to Y$ es \textit{medible} (o $\mathcal{F}$-$\mathcal{G}$-medible) si la preimagen de todo conjunto medible es medible:
	\[
		f^{-1}(B) = \{x \in X : f(x) \in B\} \in \mathcal{F} \quad \forall B \in \mathcal{G}
	\]
\end{definition}

Cuando $Y = \mathbb{R}$ con la $\sigma$-álgebra de Borel, es suficiente verificar que $\{x : f(x) > \alpha\} \in \mathcal{F}$ para todo $\alpha \in \mathbb{R}$. En particular, toda función continua es medible respecto a las $\sigma$-álgebras de Borel.

\begin{definition}[Integral de Lebesgue]\label{def:integral_lebesgue}
	Sea $(X, \mathcal{F}, \mu)$ un espacio de medida. La \textit{integral de Lebesgue} se construye en tres pasos:
	\begin{enumerate}
		\item \textbf{Funciones simples}: Para $s = \sum_{i=1}^n a_i \chi_{A_i}$ con $a_i \geq 0$ y $A_i \in \mathcal{F}$ disjuntos, donde $\chi_{A}$ es la función característica de $A$, definimos
		\[
			\int_X s \, d\mu = \sum_{i=1}^n a_i \mu(A_i)
		\]
		\item \textbf{Funciones no negativas}: Para $f: X \to [0, +\infty]$ medible,
		\[
			\int_X f \, d\mu = \sup\left\{\int_X s \, d\mu : s \text{ simple, } 0 \leq s \leq f\right\}
		\]
		\item \textbf{Funciones generales}: Para $f: X \to \mathbb{R}$ medible, escribimos $f = f^+ - f^-$ con $f^+ = \max(f, 0)$ y $f^- = \max(-f, 0)$. Decimos que $f$ es \textit{integrable} (o $f \in L^1(\mu)$) si $\int_X f^+ \, d\mu < \infty$ y $\int_X f^- \, d\mu < \infty$, y definimos
		\[
			\int_X f \, d\mu = \int_X f^+ \, d\mu - \int_X f^- \, d\mu
		\]
	\end{enumerate}
\end{definition}

\begin{definition}[Casi por doquier]\label{def:casi_por_doquier}
	Sea $(X, \mathcal{F}, \mu)$ un espacio de medida. Una propiedad $P(x)$ se cumple \textit{casi por doquier} (c.p.d.) o \textit{para casi todo $x$} (abreviado a.e. del inglés \textit{almost everywhere}) si el conjunto de puntos donde no se cumple tiene medida cero:
	\[
		\mu(\{x \in X : \neg P(x)\}) = 0
	\]
\end{definition}

En el contexto de la medida de Lebesgue en $\mathbb{R}^n$, ``casi por doquier'' significa ``excepto en un conjunto de medida de Lebesgue nula''.

\paragraph{Preliminares de análisis funcional}

\begin{definition}[Espacio vectorial]\label{def:espacio_vectorial}
	Un \textit{espacio vectorial} sobre $\mathbb{R}$ es un conjunto $V$ equipado con dos operaciones, suma $+: V \times V \to V$ y producto por escalares $\cdot: \mathbb{R} \times V \to V$, que satisfacen los siguientes axiomas para todo $u, v, w \in V$ y $\alpha, \beta \in \mathbb{R}$:
	\begin{enumerate}
		\item \textbf{Asociatividad de la suma:} $(u + v) + w = u + (v + w)$
		\item \textbf{Conmutatividad de la suma:} $u + v = v + u$
		\item \textbf{Elemento neutro:} existe $0 \in V$ tal que $v + 0 = v$
		\item \textbf{Elemento opuesto:} para cada $v$ existe $-v \in V$ tal que $v + (-v) = 0$
		\item \textbf{Compatibilidad del producto:} $\alpha(\beta v) = (\alpha\beta)v$
		\item \textbf{Elemento unidad:} $1 \cdot v = v$
		\item \textbf{Distributividad respecto a la suma vectorial:} $\alpha(u + v) = \alpha u + \alpha v$
		\item \textbf{Distributividad respecto a la suma escalar:} $(\alpha + \beta)v = \alpha v + \beta v$
	\end{enumerate}
\end{definition}

\begin{definition}[Espacio vectorial normado]\label{def:espacio_normado}
	Un \textit{espacio vectorial normado} es un par $(X, \|\cdot\|)$ donde $X$ es un espacio vectorial real y $\|\cdot\|: X \to [0, +\infty)$ es una \textit{norma}, es decir, una función que satisface:
	\begin{enumerate}
		\item $\|x\| = 0$ si y solo si $x = 0$
		\item $\|\alpha x\| = |\alpha| \|x\|$ para todo $\alpha \in \mathbb{R}$
		\item $\|x + y\| \leq \|x\| + \|y\|$ (desigualdad triangular)
	\end{enumerate}
\end{definition}

Ejemplos fundamentales incluyen $\mathbb{R}^n$ con la norma euclidiana y los espacios de funciones $L^p(\Omega)$.

\begin{definition}[Recubrimiento]\label{def:recubrimiento}
	Sea $X$ un espacio topológico y $K \subseteq X$. Un \textit{recubrimiento} de $K$ es una familia de conjuntos $\{U_\alpha\}_{\alpha \in A}$ tal que $K \subseteq \bigcup_{\alpha \in A} U_\alpha$. Si cada $U_\alpha$ es un conjunto abierto de $X$, se dice que es un \textit{recubrimiento abierto}. Un \textit{subrecubrimiento} es una subfamilia que sigue siendo un recubrimiento de $K$.
\end{definition}

\begin{definition}[Espacio compacto]\label{def:espacio_compacto}
	Un subconjunto $K$ de un espacio topológico es \textit{compacto} si todo recubrimiento abierto de $K$ admite un subrecubrimiento finito.
\end{definition}

Intuitivamente, la compacidad generaliza las propiedades de los conjuntos cerrados y acotados en $\mathbb{R}^n$ a espacios topológicos generales.

\begin{theorem}[Heine-Borel]\label{thm:heine_borel}
	Un subconjunto $K \subset \mathbb{R}^n$ es compacto si y solo si es cerrado y acotado \citep[Teorema~2.41]{Rudin1976}.
\end{theorem}

En particular, el hipercubo unitario $I_m = [0,1]^m$ es compacto, lo cual es esencial para el teorema de aproximación universal.

\begin{definition}[Funcional lineal]\label{def:funcional_lineal}
	Sea $X$ un espacio vectorial real. Un \textit{funcional lineal} es una aplicación $f: X \to \mathbb{R}$ que satisface:
	\[
		f(\alpha x + \beta y) = \alpha f(x) + \beta f(y) \quad \forall x, y \in X, \, \forall \alpha, \beta \in \mathbb{R}
	\]
	Es decir, un funcional lineal es un elemento del espacio dual algebraico $X^*$.
\end{definition}

\begin{definition}[Funcional lineal continuo]\label{def:funcional_continuo}
	Sea $(X, \|\cdot\|)$ un espacio vectorial normado. Un funcional lineal $f: X \to \mathbb{R}$ es \textit{continuo} (o \textit{acotado}) si existe una constante $M > 0$ tal que:
	\[
		|f(x)| \leq M\|x\| \quad \forall x \in X
	\]
\end{definition}

El conjunto de todos los funcionales lineales continuos en $X$ forma el \textit{espacio dual} $X'$, que es un espacio de Banach con la norma del operador $\|f\| = \sup_{\|x\| \leq 1} |f(x)|$.

\begin{definition}[Espacio $C(K)$]\label{def:espacio_CK}
	Sea $K$ un espacio compacto. El espacio $C(K)$ es el conjunto de todas las funciones continuas $f: K \to \mathbb{R}$ equipado con la \textit{norma supremo} (o norma uniforme):
	\[
		\|f\|_\infty = \sup_{x \in K} |f(x)|
	\]
\end{definition}

Con esta norma, $C(K)$ es un espacio de Banach (espacio vectorial normado completo). En el contexto del teorema de aproximación universal, trabajamos con $C(I_m)$, el espacio de funciones continuas sobre el hipercubo unitario.

\begin{theorem}[Convergencia dominada de Lebesgue]\label{thm:convergencia_dominada}
	Sea \((X, \mathcal{F}, \mu)\) un espacio de medida y \(\{f_n\}_{n=1}^\infty\) una sucesión de funciones medibles \(f_n: X \to \mathbb{R}\) tal que:
	\begin{enumerate}
		\item \(f_n(x) \to f(x)\) para casi todo \(x \in X\)
		\item Existe una función integrable \(g: X \to \mathbb{R}\) tal que \(|f_n(x)| \leq g(x)\) para casi todo \(x \in X\) y todo \(n \in \mathbb{N}\)
	\end{enumerate}
	Entonces \(f\) es integrable y:
	\[
		\lim_{n \to \infty} \int_X f_n \, d\mu = \int_X f \, d\mu = \int_X \lim_{n \to \infty} f_n \, d\mu
	\]
	\citep[Teorema~2.24]{Folland1999}.
\end{theorem}

Este teorema es fundamental para intercambiar límites e integrales, lo cual será necesario para demostrar que las funciones sigmoidales son discriminatorias.

\begin{definition}[Conjunto convexo]\label{def:conjunto_convexo}
	Un subconjunto $C$ de un espacio vectorial $V$ es \textit{convexo} si para todo $x, y \in C$ y $\lambda \in [0,1]$, se tiene que $\lambda x + (1-\lambda)y \in C$. Equivalentemente, $C$ contiene el segmento que une cualquier par de sus puntos.
\end{definition}

\begin{theorem}[Hahn-Banach, forma geométrica]\label{thm:hahn_banach}
	Sean \(X\) un espacio vectorial topológico real, \(A\) y \(B\) subconjuntos convexos disjuntos de \(X\) con \(A\) abierto. Entonces existe un funcional lineal continuo \(f: X \to \mathbb{R}\) y una constante \(\alpha \in \mathbb{R}\) tales que:
	\[
		f(a) < \alpha \leq f(b) \quad \forall a \in A, b \in B
	\]
	\citep[Teorema~5.5]{Folland1999}.
\end{theorem}

\begin{corollary}[Caracterización de densidad]\label{cor:caracterizacion_densidad}
	Sea \(X\) un espacio vectorial normado y \(Y \subset X\) un subespacio. Entonces \(Y\) es denso en \(X\) si y solo si el único funcional lineal continuo que se anula en \(Y\) es el funcional cero.
\end{corollary}

\begin{proof}
	\textbf{(\(\Rightarrow\))} Supongamos que $Y$ es denso en $X$. Sea $f: X \to \mathbb{R}$ un funcional lineal continuo tal que $f(y) = 0$ para todo $y \in Y$. Para cualquier $x \in X$, existe una sucesión $\{y_n\} \subset Y$ con $y_n \to x$. Por la continuidad de $f$:
	\[
		f(x) = f\!\left(\lim_{n \to \infty} y_n\right) = \lim_{n \to \infty} f(y_n) = 0
	\]
	Por tanto $f \equiv 0$.

	\textbf{(\(\Leftarrow\))} Supongamos que $Y$ no es denso en $X$. Entonces $\overline{Y} \neq X$, y existe $x_0 \in X \setminus \overline{Y}$. Como $\overline{Y}$ es un subespacio cerrado, tomamos una bola abierta $B_r(0) \subset X$ y definimos $A = \overline{Y} + B_r(0)$, que es abierto y convexo, con $A \cap \{x_0\} = \emptyset$ para $r$ suficientemente pequeño (ya que $x_0 \notin \overline{Y}$). Por el \autoref{thm:hahn_banach}, existe un funcional lineal $f$ y $\alpha \in \mathbb{R}$ tales que:
	\[
		f(a) < \alpha \leq f(x_0) \quad \forall a \in A
	\]
	Como $\overline{Y}$ es un subespacio, $0 \in \overline{Y} \subset A$, luego $f(0) = 0 < \alpha$. Además, para cualquier $y \in Y$ y $t \in \mathbb{R}$, se tiene $ty \in \overline{Y} \subset A$, de modo que $tf(y) = f(ty) < \alpha$ para todo $t$. Esto solo es posible si $f(y) = 0$. Sin embargo, $f(x_0) \geq \alpha > 0$, por lo que $f \not\equiv 0$.
\end{proof}

\begin{theorem}[Representación de Riesz para \(C(K)\)]\label{thm:riesz_representacion}
	Sea \(K\) un espacio compacto y \(C(K)\) el espacio de funciones continuas reales en \(K\) con la norma supremo. Para todo funcional lineal continuo \(L: C(K) \to \mathbb{R}\), existe una única medida de Borel regular finita con signo \(\mu\) en \(K\) tal que:
	\[
		L(f) = \int_K f \, d\mu \quad \forall f \in C(K)
	\]
	Además, \(\|L\| = |\mu|(K)\), donde \(|\mu|\) es la variación total de \(\mu\) \citep[Teorema~7.17]{Folland1999}.
\end{theorem}

\begin{theorem}[Stone-Weierstrass]\label{thm:stone_weierstrass}
	Sea \(K\) un espacio compacto y \(\mathcal{A} \subset C(K)\) una subálgebra que separa puntos de \(K\) y no se anula en ningún punto. Entonces \(\mathcal{A}\) es densa en \(C(K)\) con respecto a la norma supremo \citep[Teorema~7.32]{Rudin1976}.
\end{theorem}

En particular, si \(\mathcal{A}\) contiene las funciones constantes y separa puntos, entonces es densa en \(C(K)\).

Existen dos enfoques principales para demostrar el teorema de aproximación universal:
\begin{itemize}
	\item \textbf{Enfoque de \citet{cybenko1989}}: Usa el corolario de Hahn-Banach junto con el teorema de representación de Riesz. Demuestra que las funciones sigmoidales son ``discriminatorias'' (se anulan en una medida solo si la medida es cero). Esta es la demostración que presentamos a continuación.
	\item \textbf{Enfoque de \citet{hornik1989}}: Usa el teorema de Stone-Weierstrass, demostrando que las redes neuronales forman una subálgebra que separa puntos.
\end{itemize}
Incluimos Stone-Weierstrass en los preliminares por su importancia histórica y su conexión con el enfoque alternativo.

\paragraph{Enunciado y demostración}

El teorema de aproximación universal es una justificación teórica fundamental para el uso de redes neuronales para aproximar funciones, demostrado independientemente por \citet{cybenko1989} y \citet{hornik1989}:

\begin{definition}[Función sigmoidal]\label{def:funcion_sigmoidal}
	Una función \(\phi: \mathbb{R} \to \mathbb{R}\) es \textit{sigmoidal} si es continua, acotada y satisface:
	\[
		\lim_{t \to -\infty} \phi(t) = 0 \quad \text{y} \quad \lim_{t \to +\infty} \phi(t) = 1
	\]
	Ejemplos clásicos incluyen la función logística \(\sigma(t) = \frac{1}{1+e^{-t}}\) y la tangente hiperbólica reescalada, ambas descritas en detalle en la \autoref{subsubsec:funciones_activacion_capas}.
\end{definition}

\begin{theorem}[Aproximación Universal - Cybenko, 1989]\label{thm:aproximacion_universal}
	Sea \(\phi: \mathbb{R} \to \mathbb{R}\) una función sigmoidal. Sea \(I_m = [0,1]^m\) el hipercubo unitario \(m\)-dimensional. Entonces, el conjunto de funciones de la forma:
	\[
		F(x) = \sum_{i=1}^{N} v_i \phi(w_i^T x + b_i)
	\]
	con \(N \in \mathbb{N}\), \(v_i, b_i \in \mathbb{R}\) y \(w_i \in \mathbb{R}^m\), es denso en \(C(I_m)\) con respecto a la norma uniforme.

	Es decir, para cualquier \(f \in C(I_m)\) y cualquier \(\epsilon > 0\), existen parámetros tales que:
	\[
		\sup_{x \in I_m} |F(x) - f(x)| < \epsilon
	\]
\end{theorem}

\textbf{Por qué el hipercubo unitario.}
La restricción a \(I_m = [0,1]^m\) no es una limitación real:
\begin{itemize}
	\item Para cualquier compacto \(K \subset \mathbb{R}^m\), por el teorema de Heine-Borel (\autoref{thm:heine_borel}), \(K\) está acotado y por tanto contenido en un hipercubo \([-M, M]^m\). Mediante una transformación afín se puede trasladar el problema a \(I_m\), por lo que el resultado se extiende a \(C(K)\) para cualquier compacto \(K \subset \mathbb{R}^m\)
	\item La compacidad es esencial: garantiza la continuidad uniforme de \(f\) y la existencia de aproximaciones
	\item En la práctica, los datos siempre están acotados, por lo que esta restricción es natural
\end{itemize}

\textbf{Extensiones del teorema.}
El teorema se ha extendido significativamente desde 1989:
\begin{itemize}
	\item \textbf{Funciones de activación más generales}: El teorema vale para cualquier función no polinomial \citep{leshno1993multilayer}
	\item \textbf{ReLU}: Para la función \gls{ReLU} (véase \autoref{subsubsec:funciones_activacion_capas}), el resultado de Leshno et al. implica que una capa oculta es suficiente con anchura arbitraria. Sin embargo, con anchura acotada, se requiere profundidad suficiente para lograr aproximación universal \citep{lu2017expressive,hanin2019universal}
	\item \textbf{Aproximación en \(L^p\)}: Resultados similares valen para otras normas \citep{hornik1991approximation}
	\item \textbf{Tasas de aproximación}: Se conocen cotas sobre el número de neuronas \(N\) necesarias en función de \(\epsilon\) y la regularidad de \(f\) \citep{barron1993universal}
\end{itemize}

\begin{proof}
	La demostración sigue el enfoque de \citet{cybenko1989} usando análisis funcional. Procedemos en varios pasos:

	La estrategia es indirecta: en lugar de construir explícitamente la aproximación, demostraremos que ninguna función continua puede ``escapar'' de ser aproximada. Para ello, usamos dualidad en análisis funcional. Si las redes neuronales no fueran densas en \(C(I_m)\), existiría un funcional lineal continuo no nulo que se anula en todas ellas (corolario de Hahn-Banach). Traducido a medidas mediante el teorema de Riesz, esto significaría que existe una medida no nula ``invisible'' para todas las neuronas. Demostrar que tal medida no puede existir (es decir, que \(\phi\) es \textit{discriminatoria}) es el corazón de la prueba.

	\textbf{Paso 1: Funciones discriminatorias.}
	Definimos el conjunto de funciones sigmoidales discriminatorias. Una función \(\phi\) es discriminatoria si para cualquier medida finita con signo \(\mu\) en \(I_m\):
	\[
		\int_{I_m} \phi(w^T x + b) \, d\mu(x) = 0 \quad \forall w \in \mathbb{R}^m, b \in \mathbb{R} \implies \mu = 0
	\]

	En otras palabras, la familia de funciones \(\{\phi(w^T x + b)\}_{w,b}\) es suficientemente rica como para ``detectar'' cualquier medida no nula: no existe ninguna distribución de masa no trivial que sea invisible para todas las neuronas simultáneamente. El término ``discriminatoria'' refleja esta capacidad de discriminar entre la medida cero y cualquier otra medida, de forma análoga a cómo las funciones test en la teoría de distribuciones distinguen distribuciones distintas.

	\textbf{Paso 2: Las funciones sigmoidales son discriminatorias.}
	Sea \(\phi\) una función sigmoidal (continua, acotada, \(\lim_{t \to -\infty} \phi(t) = 0\), \(\lim_{t \to +\infty} \phi(t) = 1\)). Supongamos que \(\mu\) es una medida tal que \(\int_{I_m} \phi(w^T x + b) \, d\mu(x) = 0\) para todo \(w, b\). Debemos demostrar que \(\mu = 0\).

	Para cualquier \(w \in \mathbb{R}^m\), \(b \in \mathbb{R}\) y parámetro auxiliar \(\alpha \in \mathbb{R}\), definimos la familia de funciones:
	\[
		\psi_{\lambda}(x) = \phi(\lambda(w^T x + b) + \alpha)
	\]

	La idea es que al multiplicar el argumento de \(\phi\) por \(\lambda\) grande, la transición suave de la sigmoide se ``comprime'' hasta convertirse en un salto abrupto en el hiperplano \(w^T x + b = 0\). Geométricamente, cada hiperplano divide \(I_m\) en dos semiespacios, y el límite \(\lambda \to \infty\) convierte la neurona en un indicador del semiespacio positivo.

	Cuando \(\lambda \to \infty\), tenemos convergencia puntual a una función escalonada:
	\[
		\psi_{\lambda}(x) \to \begin{cases}
			                        1 & \text{si } w^T x + b > 0 \\
			                        0 & \text{si } w^T x + b < 0 \\
			                        \phi(\alpha) & \text{si } w^T x + b = 0
		\end{cases}
	\]

	Como \(|\psi_{\lambda}(x)| \leq \sup_{t \in \mathbb{R}} |\phi(t)| < \infty\) uniformemente en \(\lambda\) y \(x\), y esta cota es integrable sobre el compacto \(I_m\), podemos aplicar el \textbf{teorema de convergencia dominada de Lebesgue}:
	\[
		\int_{I_m} \psi_{\lambda}(x) \, d\mu(x) \to \mu(H^+_{w,b}) + \phi(\alpha) \cdot \mu(H_{w,b})
	\]
	donde \(H^+_{w,b} = \{x \in I_m : w^T x + b > 0\}\) es el semiespacio abierto y \(H_{w,b} = \{x \in I_m : w^T x + b = 0\}\) es el hiperplano.

	Por hipótesis, el lado izquierdo es cero para todo \(\lambda\). Por continuidad del límite, obtenemos:
	\[
		\mu(H^+_{w,b}) + \phi(\alpha) \cdot \mu(H_{w,b}) = 0 \quad \text{para todo } \alpha \in \mathbb{R}
	\]

	Para ver por qué, observemos que la ecuación \(\mu(H^+_{w,b}) + \phi(\alpha) \cdot \mu(H_{w,b}) = 0\) es una función afín de \(\phi(\alpha)\). Si \(\mu(H_{w,b}) \neq 0\), esta ecuación determinaría un único valor de \(\phi(\alpha)\); pero como \(\phi\) toma un rango continuo de valores (de 0 a 1), la ecuación no podría satisfacerse para todo \(\alpha\). Así, necesariamente \(\mu(H_{w,b}) = 0\) y, sustituyendo, \(\mu(H^+_{w,b}) = 0\).

	Como esta ecuación vale para todo \(\alpha\) y \(\phi\) no es constante (ya que \(\phi(\alpha) \to 0\) cuando \(\alpha \to -\infty\) y \(\phi(\alpha) \to 1\) cuando \(\alpha \to +\infty\)), necesariamente \(\mu(H_{w,b}) = 0\) y \(\mu(H^+_{w,b}) = 0\).

	Esto vale para todo \(w\) y \(b\), por lo que \(\mu\) se anula en todos los semiespacios abiertos. Como los semiespacios generan la \(\sigma\)-álgebra de Borel y las medidas finitas con signo están determinadas por sus valores en una clase generadora estable bajo intersecciones, concluimos que \(\mu = 0\).

	\textbf{Paso 3: Densidad en \(C(I_m)\).}
	Sea \(\mathcal{S}\) el conjunto de funciones de la forma:
	\[
		G(x) = \sum_{i=1}^{N} v_i \phi(w_i^T x + b_i)
	\]

	La conexión entre la propiedad discriminatoria y la densidad se establece mediante dos teoremas clásicos de análisis funcional. El corolario de Hahn-Banach (\autoref{cor:caracterizacion_densidad}) reduce el problema de densidad a verificar que ningún funcional no nulo se anula en \(\mathcal{S}\). El teorema de Riesz (\autoref{thm:riesz_representacion}) traduce funcionales a medidas, llevando el problema exactamente a la condición de la definición de función discriminatoria.

	Por el Paso 2, \(\phi\) es discriminatoria. Aplicamos el \textbf{corolario del teorema de Hahn-Banach}: el conjunto \(\mathcal{S}\) es denso en \(C(I_m)\) si y solo si el único funcional lineal continuo que se anula en \(\mathcal{S}\) es el funcional cero.

	Sea \(L\) un funcional lineal acotado en \(C(I_m)\) tal que \(L(G) = 0\) para todo \(G \in \mathcal{S}\). Por el \textbf{teorema de representación de Riesz para \(C(K)\)}, existe una única medida finita con signo \(\mu\) tal que:
	\[
		L(h) = \int_{I_m} h(x) \, d\mu(x) \quad \forall h \in C(I_m)
	\]

	Como \(L(G) = 0\) para todo \(G \in \mathcal{S}\), en particular:
	\[
		\int_{I_m} \phi(w^T x + b) \, d\mu(x) = 0 \quad \forall w \in \mathbb{R}^m, b \in \mathbb{R}
	\]

	Por la definición de función discriminatoria y el Paso 2, esto implica \(\mu = 0\), luego \(L = 0\).

	\textbf{Paso 4: Aproximación uniforme.}
	Por el paso anterior, \(\mathcal{S}\) es denso en \(C(I_m)\) con respecto a la norma supremo. Por definición de densidad, para cualquier \(f \in C(I_m)\) y \(\epsilon > 0\), existe \(F \in \mathcal{S}\) tal que:
	\[
		\|F - f\|_\infty = \sup_{x \in I_m} |F(x) - f(x)| < \epsilon
	\]

	Esto completa la demostración. \(\square\)
\end{proof}

Este teorema garantiza que una red neuronal con una sola capa oculta puede aproximar cualquier función continua con precisión arbitraria. Sin embargo, es un resultado de \textbf{existencia}, no de \textbf{construcción}: demuestra que los pesos adecuados existen, pero no dice cómo encontrarlos. Además, el número de neuronas necesarias puede ser prohibitivamente grande, y las redes profundas (múltiples capas) pueden representar funciones complejas de manera más eficiente, permitiendo aprender representaciones jerárquicas de los datos.

\subsection{Teoría de convergencia}\label{subsec:teoria_convergencia}

El teorema de aproximación universal garantiza la existencia de una configuración de pesos que aproxima cualquier función objetivo. En la \autoref{subsec:descenso_gradiente} hemos presentado el descenso del gradiente y la retropropagación como el mecanismo para buscar esos pesos. Ahora formalizamos las condiciones bajo las cuales este proceso converge efectivamente y establecemos sus tasas de convergencia. Para un tratamiento en profundidad, véase \citet{bottou2018optimization}.

\subsubsection{Condiciones de regularidad}

Los teoremas de convergencia del descenso del gradiente no se cumplen para funciones de pérdida arbitrarias: requieren ciertas hipótesis de regularidad. En esta sección formalizamos las dos condiciones fundamentales: la \textit{$L$-suavidad}, que controla cuánto puede variar el gradiente y permite acotar el progreso por iteración, y la \textit{convexidad}, que garantiza la ausencia de mínimos locales espurios. Estas propiedades proporcionan el marco analítico necesario para los resultados de convergencia que presentamos a continuación.

\begin{definition}[Función $L$-suave]\label{def:L_smooth}
	Una función diferenciable \(f: \mathbb{R}^n \to \mathbb{R}\) es \textit{\(L\)-suave} (o tiene gradiente \(L\)-Lipschitz) si existe una constante \(L > 0\) tal que:
	\[
		\|\nabla f(x) - \nabla f(y)\| \leq L\|x - y\| \quad \forall x, y \in \mathbb{R}^n
	\]
\end{definition}

Esta condición es análoga a la continuidad de Lipschitz (\autoref{prop:lipschitz_continuidad}), pero aplicada al gradiente en lugar de a la función misma.

\begin{proposition}[Caracterización equivalente de $L$-suavidad]\label{prop:L_smooth_equiv}
	Sea \(f: \mathbb{R}^n \to \mathbb{R}\) diferenciable. Las siguientes afirmaciones son equivalentes:
	\begin{enumerate}
		\item \(f\) es \(L\)-suave (gradiente \(L\)-Lipschitz)
		\item Para todo \(x, y \in \mathbb{R}^n\): \(f(y) \leq f(x) + \langle \nabla f(x), y - x \rangle + \frac{L}{2}\|y - x\|^2\)
	\end{enumerate}
\end{proposition}

\begin{proof}
	\textbf{(1) \(\Rightarrow\) (2):} Definimos \(g(t) = f(x + t(y-x))\) para \(t \in [0,1]\). Por el teorema fundamental del cálculo:
	\[
		f(y) - f(x) = g(1) - g(0) = \int_0^1 g'(t) \, dt = \int_0^1 \langle \nabla f(x + t(y-x)), y-x \rangle \, dt
	\]
	Restando \(\langle \nabla f(x), y-x \rangle\):
	\begin{align*}
		f(y) - f(x) - \langle \nabla f(x), y-x \rangle &= \int_0^1 \langle \nabla f(x + t(y-x)) - \nabla f(x), y-x \rangle \, dt \\
		&\leq \int_0^1 \|\nabla f(x + t(y-x)) - \nabla f(x)\| \cdot \|y-x\| \, dt \\
		&\leq \int_0^1 L \cdot t\|y-x\| \cdot \|y-x\| \, dt = \frac{L}{2}\|y-x\|^2
	\end{align*}

	\textbf{(2) \(\Rightarrow\) (1):} Aplicando (2) dos veces (intercambiando \(x\) e \(y\)) y sumando:
	\begin{align*}
		f(y) &\leq f(x) + \langle \nabla f(x), y-x \rangle + \frac{L}{2}\|y-x\|^2 \\
		f(x) &\leq f(y) + \langle \nabla f(y), x-y \rangle + \frac{L}{2}\|x-y\|^2
	\end{align*}
	Sumando: \(0 \leq \langle \nabla f(x) - \nabla f(y), y-x \rangle + L\|y-x\|^2\), es decir:
	\[
		\langle \nabla f(x) - \nabla f(y), x-y \rangle \leq L\|x-y\|^2
	\]
	Por Cauchy-Schwarz: \(\|\nabla f(x) - \nabla f(y)\| \cdot \|x-y\| \geq \langle \nabla f(x) - \nabla f(y), x-y \rangle\), pero esta desigualdad por sí sola no basta. Para completar rigurosamente, se requiere que \(f\) sea convexa: en ese caso, la desigualdad anterior implica co-coercividad del gradiente, que a su vez da la condición de Lipschitz. Véase \citet{Nesterov2018}, Teorema~2.1.5, para los detalles.
\end{proof}

\begin{remark}
	La dirección (1) \(\Rightarrow\) (2) vale para toda función diferenciable, sin necesidad de convexidad. La dirección (2) \(\Rightarrow\) (1), en cambio, requiere que \(f\) sea convexa \citep[Teorema~2.1.5]{Nesterov2018}. Dado que en los teoremas de convergencia posteriores (\autoref{thm:gd_convergence_convex}, \autoref{thm:gd_nonconvex}) solo se utiliza la dirección (1) \(\Rightarrow\) (2), esta restricción no afecta a los resultados del resto del capítulo.
\end{remark}

La constante \(L\) controla cuánto puede variar el gradiente localmente. Intuitivamente, funciones con \(L\) pequeño son ``más planas'' y permiten pasos de optimización más grandes.

La tasa \(\eta = 1/L\) es óptima en el siguiente sentido. La cota de \(L\)-suavidad
(\autoref{prop:L_smooth_equiv}) garantiza que para todo par de puntos:
\[
	f(y) \leq f(x) + \langle \nabla f(x), y - x \rangle + \frac{L}{2}\|y - x\|^2
\]
Sustituyendo el paso de descenso del gradiente con tasa genérica \(\eta > 0\),
es decir, \(x = \theta_t\) e \(y = \theta_{t+1} = \theta_t - \eta\nabla f(\theta_t)\),
de modo que \(y - x = -\eta\nabla f(\theta_t)\):
\begin{align*}
	f(\theta_{t+1})
		&\leq f(\theta_t) + \langle \nabla f(\theta_t),\, -\eta\nabla f(\theta_t) \rangle
			+ \frac{L}{2}\|\!-\eta\nabla f(\theta_t)\|^2 \\
		&= f(\theta_t) - \eta\|\nabla f(\theta_t)\|^2
			+ \frac{L\eta^2}{2}\|\nabla f(\theta_t)\|^2 \\
		&= f(\theta_t)
			- \underbrace{\left(\eta - \frac{L\eta^2}{2}\right)}_{g(\eta)}
			\|\nabla f(\theta_t)\|^2
\end{align*}
El coeficiente \(g(\eta) = \eta - \frac{L\eta^2}{2}\) es una parábola cóncava que alcanza su máximo cuando \(g'(\eta) = 1 - L\eta = 0\), es decir, \(\eta^* = 1/L\), con valor \(g(1/L) = 1/(2L)\). Además, \(g(\eta) > 0\) requiere \(\eta < 2/L\): para tasas \(\eta \geq 2/L\) la cota deja de garantizar descenso y el algoritmo puede divergir. Para \(0 < \eta < 1/L\) el algoritmo converge, pero con decrecimiento garantizado subóptimo.

\begin{definition}[Función convexa]\label{def:convex_function}
	Una función \(f: \mathbb{R}^n \to \mathbb{R}\) es \textit{convexa} si para todo \(x, y \in \mathbb{R}^n\) y \(\lambda \in [0,1]\):
	\[
		f(\lambda x + (1-\lambda)y) \leq \lambda f(x) + (1-\lambda)f(y)
	\]
\end{definition}

Esta definición está íntimamente relacionada con la noción de conjunto convexo (\autoref{def:conjunto_convexo}): una función es convexa si y solo si su epígrafo $\{(x,y) \in \mathbb{R}^{n+1} : y \geq f(x)\}$ es un conjunto convexo.

\begin{proposition}[Caracterización de convexidad para funciones diferenciables]\label{prop:convex_equiv}
	Sea \(f: \mathbb{R}^n \to \mathbb{R}\) diferenciable. Entonces \(f\) es convexa si y solo si:
	\[
		f(y) \geq f(x) + \langle \nabla f(x), y - x \rangle \quad \forall x, y \in \mathbb{R}^n
	\]
	Es decir, el grafo de \(f\) está por encima de todos sus hiperplanos tangentes.
\end{proposition}

\begin{proof}
	\textbf{(\(\Rightarrow\)):} Supongamos \(f\) convexa. Para cualquier \(x, y\) y \(\lambda \in (0,1]\):
	\[
		f(x + \lambda(y-x)) \leq (1-\lambda)f(x) + \lambda f(y)
	\]
	Reordenando: \(\frac{f(x + \lambda(y-x)) - f(x)}{\lambda} \leq f(y) - f(x)\).

	Tomando \(\lambda \to 0^+\), el lado izquierdo converge a la derivada direccional \(\langle \nabla f(x), y-x \rangle\), obteniendo \(\langle \nabla f(x), y-x \rangle \leq f(y) - f(x)\).

	\textbf{(\(\Leftarrow\)):} Supongamos que vale la condición del gradiente. Sean \(x, y \in \mathbb{R}^n\) y \(\lambda \in [0,1]\). Definimos \(z = \lambda x + (1-\lambda)y\). Aplicando la hipótesis:
	\begin{align*}
		f(x) &\geq f(z) + \langle \nabla f(z), x - z \rangle = f(z) + (1-\lambda)\langle \nabla f(z), x - y \rangle \\
		f(y) &\geq f(z) + \langle \nabla f(z), y - z \rangle = f(z) - \lambda\langle \nabla f(z), x - y \rangle
	\end{align*}
	Multiplicando la primera por \(\lambda\) y la segunda por \((1-\lambda)\) y sumando:
	\[
		\lambda f(x) + (1-\lambda)f(y) \geq f(z) = f(\lambda x + (1-\lambda)y)
	\]
\end{proof}

\subsubsection{Convergencia del descenso del gradiente}

\begin{theorem}[Convergencia de \gls{GD} para funciones convexas]\label{thm:gd_convergence_convex}
	Sea \(f: \mathbb{R}^n \to \mathbb{R}\) convexa y \(L\)-suave con mínimo global \(f^* = f(\theta^*)\). Si aplicamos descenso del gradiente con tasa de aprendizaje \(\eta = 1/L\), entonces:
	\[
		f(\theta_T) - f^* \leq \frac{L\|\theta_0 - \theta^*\|^2}{2T}
	\]
	Es decir, el descenso del gradiente converge con tasa \(O(1/T)\) al óptimo.
\end{theorem}

\begin{proof}
	Procedemos en varios pasos.

	\textbf{Paso 1: Desigualdad de descenso.}
	Por la \autoref{prop:L_smooth_equiv}, para cualquier par de puntos:
	\[
		f(y) \leq f(x) + \langle \nabla f(x), y - x \rangle + \frac{L}{2}\|y - x\|^2
	\]
	Aplicando esta desigualdad con \(x = \theta_t\) y \(y = \theta_{t+1} = \theta_t - \frac{1}{L}\nabla f(\theta_t)\):
	\begin{align*}
		f(\theta_{t+1}) &\leq f(\theta_t) + \left\langle \nabla f(\theta_t), -\frac{1}{L}\nabla f(\theta_t) \right\rangle + \frac{L}{2}\left\|-\frac{1}{L}\nabla f(\theta_t)\right\|^2 \\
		&= f(\theta_t) - \frac{1}{L}\|\nabla f(\theta_t)\|^2 + \frac{1}{2L}\|\nabla f(\theta_t)\|^2 \\
		&= f(\theta_t) - \frac{1}{2L}\|\nabla f(\theta_t)\|^2
	\end{align*}
	Esta desigualdad garantiza que el valor de la función decrece en cada iteración.

	\textbf{Paso 2: Cota sobre la distancia al óptimo.}
	Analizamos cómo evoluciona la distancia \(\|\theta_{t+1} - \theta^*\|^2\):
	\begin{align*}
		\|\theta_{t+1} - \theta^*\|^2 &= \left\|\theta_t - \frac{1}{L}\nabla f(\theta_t) - \theta^*\right\|^2 \\
		&= \|\theta_t - \theta^*\|^2 - \frac{2}{L}\langle \nabla f(\theta_t), \theta_t - \theta^* \rangle + \frac{1}{L^2}\|\nabla f(\theta_t)\|^2
	\end{align*}

	\textbf{Paso 3: Uso de la convexidad.}
	Por la \autoref{prop:convex_equiv}:
	\[
		f(\theta^*) \geq f(\theta_t) + \langle \nabla f(\theta_t), \theta^* - \theta_t \rangle
	\]
	Reordenando:
	\[
		\langle \nabla f(\theta_t), \theta_t - \theta^* \rangle \geq f(\theta_t) - f^*
	\]

	\textbf{Paso 4: Combinación de las cotas.}
	Sustituyendo la desigualdad del Paso 3 en la del Paso 2:
	\[
		\|\theta_{t+1} - \theta^*\|^2 \leq \|\theta_t - \theta^*\|^2 - \frac{2}{L}(f(\theta_t) - f^*) + \frac{1}{L^2}\|\nabla f(\theta_t)\|^2
	\]

	Del Paso 1, tenemos \(\|\nabla f(\theta_t)\|^2 \leq 2L(f(\theta_t) - f(\theta_{t+1}))\). Sustituyendo:
	\begin{align*}
		\|\theta_{t+1} - \theta^*\|^2 &\leq \|\theta_t - \theta^*\|^2 - \frac{2}{L}(f(\theta_t) - f^*) + \frac{2}{L}(f(\theta_t) - f(\theta_{t+1})) \\
		&= \|\theta_t - \theta^*\|^2 - \frac{2}{L}(f(\theta_{t+1}) - f^*)
	\end{align*}

	Reordenando:
	\begin{equation}\label{eq:gd_telescopic}
		f(\theta_{t+1}) - f^* \leq \frac{L}{2}\left(\|\theta_t - \theta^*\|^2 - \|\theta_{t+1} - \theta^*\|^2\right)
	\end{equation}

	\textbf{Paso 5: Suma telescópica.}
	Sumando la desigualdad~\eqref{eq:gd_telescopic} para \(t = 0, 1, \ldots, T-1\):
	\[
		\sum_{t=0}^{T-1} (f(\theta_{t+1}) - f^*) \leq \frac{L}{2} \sum_{t=0}^{T-1} \left(\|\theta_t - \theta^*\|^2 - \|\theta_{t+1} - \theta^*\|^2\right)
	\]

	El lado derecho es una suma telescópica:
	\[
		\sum_{t=0}^{T-1} (f(\theta_{t+1}) - f^*) \leq \frac{L}{2}\left(\|\theta_0 - \theta^*\|^2 - \|\theta_T - \theta^*\|^2\right) \leq \frac{L}{2}\|\theta_0 - \theta^*\|^2
	\]

	\textbf{Paso 6: Conclusión.}
	Por el Paso 1, la sucesión \(\{f(\theta_t)\}\) es decreciente, por lo que \(f(\theta_T) \leq f(\theta_{t+1})\) para todo \(t < T\). Entonces:
	\[
		T(f(\theta_T) - f^*) \leq \sum_{t=0}^{T-1} (f(\theta_{t+1}) - f^*) \leq \frac{L}{2}\|\theta_0 - \theta^*\|^2
	\]

	Dividiendo por \(T\):
	\[
		f(\theta_T) - f^* \leq \frac{L\|\theta_0 - \theta^*\|^2}{2T}
	\]

	Esto demuestra la convergencia con tasa \(O(1/T)\).
\end{proof}

\begin{theorem}[Convergencia de SGD]\label{thm:sgd_convergence}
	Sea \(f\) convexa y \(L\)-suave, y sea \(g_t\) un estimador no sesgado del gradiente con varianza acotada:
	\[
		\mathbb{E}[g_t] = \nabla f(\theta_t), \quad \mathbb{E}[\|g_t - \nabla f(\theta_t)\|^2] \leq \sigma^2
	\]
	Supongamos además que los gradientes están uniformemente acotados: \(\|\nabla f(\theta_t)\| \leq G\) para todo \(t\). Con tasa de aprendizaje decreciente \(\eta_t = \eta_0/\sqrt{t}\), el SGD satisface:
	\[
		\mathbb{E}[f(\bar{\theta}_T)] - f^* = O\!\left(\frac{(\sigma^2 + G^2)\log T}{\sqrt{T}} + \frac{\|\theta_0 - \theta^*\|^2}{\eta_0\sqrt{T}}\right)
	\]
	donde \(\bar{\theta}_T = \frac{1}{T}\sum_{t=1}^T \theta_t\) es el promedio de los iterados.
\end{theorem}

\begin{proof}
	Presentamos un esquema de la demostración siguiendo~\citet{bottou2018optimization}.

	\textbf{Paso 1: Descomposición del gradiente estocástico.}
	Escribimos \(g_t = \nabla f(\theta_t) + \xi_t\) donde \(\xi_t = g_t - \nabla f(\theta_t)\) es el ruido estocástico con:
	\[
		\mathbb{E}[\xi_t | \theta_t] = 0, \quad \mathbb{E}[\|\xi_t\|^2 | \theta_t] \leq \sigma^2
	\]

	\textbf{Paso 2: Análisis del progreso esperado.}
	La actualización es \(\theta_{t+1} = \theta_t - \eta_t g_t\). Por convexidad:
	\[
		f(\theta^*) \geq f(\theta_t) + \langle \nabla f(\theta_t), \theta^* - \theta_t \rangle
	\]
	Reordenando: \(\langle \nabla f(\theta_t), \theta_t - \theta^* \rangle \geq f(\theta_t) - f^*\).

	Analizamos la distancia al óptimo:
	\begin{align*}
		\|\theta_{t+1} - \theta^*\|^2 &= \|\theta_t - \eta_t g_t - \theta^*\|^2 \\
		&= \|\theta_t - \theta^*\|^2 - 2\eta_t\langle g_t, \theta_t - \theta^* \rangle + \eta_t^2\|g_t\|^2
	\end{align*}

	Tomando esperanza condicional y usando que \(\mathbb{E}[g_t|\theta_t] = \nabla f(\theta_t)\):
	\[
		\mathbb{E}[\|\theta_{t+1} - \theta^*\|^2 | \theta_t] \leq \|\theta_t - \theta^*\|^2 - 2\eta_t(f(\theta_t) - f^*) + \eta_t^2(\|\nabla f(\theta_t)\|^2 + \sigma^2)
	\]

	\textbf{Paso 3: Suma telescópica y promediado.}
	Reordenamos la desigualdad del Paso~2 para aislar el término de progreso:
	\[
		2\eta_t(f(\theta_t) - f^*) \leq \|\theta_t - \theta^*\|^2 - \mathbb{E}[\|\theta_{t+1} - \theta^*\|^2 | \theta_t] + \eta_t^2(\|\nabla f(\theta_t)\|^2 + \sigma^2)
	\]
	Tomando esperanza total y sumando sobre \(t = 1, \ldots, T\):
	\[
		\sum_{t=1}^T 2\eta_t \mathbb{E}[f(\theta_t) - f^*] \leq \sum_{t=1}^T \Big(\mathbb{E}[\|\theta_t - \theta^*\|^2] - \mathbb{E}[\|\theta_{t+1} - \theta^*\|^2]\Big) + \sum_{t=1}^T \eta_t^2(\|\nabla f(\theta_t)\|^2 + \sigma^2)
	\]
	La primera suma es telescópica: los términos consecutivos se cancelan dejando únicamente los extremos,
	\[
		\sum_{t=1}^T \Big(\mathbb{E}[\|\theta_t - \theta^*\|^2] - \mathbb{E}[\|\theta_{t+1} - \theta^*\|^2]\Big) = \mathbb{E}[\|\theta_1 - \theta^*\|^2] - \mathbb{E}[\|\theta_{T+1} - \theta^*\|^2] \leq \|\theta_0 - \theta^*\|^2
	\]
	donde la última desigualdad usa que \(\mathbb{E}[\|\theta_1 - \theta^*\|^2] \leq \|\theta_0 - \theta^*\|^2\) (por el resultado del Paso~2 con \(t=0\) y \(f(\theta_0) \geq f^*\)) y que \(-\mathbb{E}[\|\theta_{T+1} - \theta^*\|^2] \leq 0\). Sustituyendo:
	\[
		\sum_{t=1}^T 2\eta_t \mathbb{E}[f(\theta_t) - f^*] \leq \|\theta_0 - \theta^*\|^2 + \sum_{t=1}^T \eta_t^2(\|\nabla f(\theta_t)\|^2 + \sigma^2)
	\]

	\textbf{Paso 4: Elección de tasa decreciente.}
	Elegimos \(\eta_t = \eta_0/\sqrt{t}\). Usando la hipótesis \(\|\nabla f(\theta_t)\| \leq G\) y sustituyendo en la desigualdad del Paso~3:
	\[
		\sum_{t=1}^T \frac{2\eta_0}{\sqrt{t}}\, \mathbb{E}[f(\theta_t) - f^*] \leq \|\theta_0 - \theta^*\|^2 + \sum_{t=1}^T \frac{\eta_0^2}{t}(G^2 + \sigma^2)
	\]

	Para acotar el lado derecho, usamos que la serie armónica satisface \(\sum_{t=1}^T \frac{1}{t} \leq 1 + \log T\):
	\[
		\sum_{t=1}^T \frac{\eta_0^2}{t}(G^2 + \sigma^2) \leq \eta_0^2(G^2 + \sigma^2)(1 + \log T)
	\]

	Para acotar el lado izquierdo por debajo, observamos que \(\eta_t = \eta_0/\sqrt{t}\) es decreciente, de modo que \(\eta_t \geq \eta_T = \eta_0/\sqrt{T}\) para todo \(t \leq T\). Como cada término \(\mathbb{E}[f(\theta_t) - f^*] \geq 0\) por convexidad:
	\[
		\sum_{t=1}^T \frac{2\eta_0}{\sqrt{t}}\, \mathbb{E}[f(\theta_t) - f^*] \geq \frac{2\eta_0}{\sqrt{T}} \sum_{t=1}^T \mathbb{E}[f(\theta_t) - f^*]
	\]

	Combinando ambas cotas y dividiendo por \(T\):
	\[
		\frac{1}{T}\sum_{t=1}^T \mathbb{E}[f(\theta_t) - f^*] \leq \frac{\|\theta_0 - \theta^*\|^2 + \eta_0^2(G^2 + \sigma^2)(1 + \log T)}{2\eta_0\sqrt{T}}
	\]

	Finalmente, por convexidad de \(f\) y la desigualdad de Jensen aplicada al promedio \(\bar{\theta}_T = \frac{1}{T}\sum_{t=1}^T \theta_t\):
	\[
		\mathbb{E}[f(\bar{\theta}_T)] \leq \frac{1}{T}\sum_{t=1}^T \mathbb{E}[f(\theta_t)]
	\]
	lo que, combinado con la cota anterior, da:
	\[
		\mathbb{E}[f(\bar{\theta}_T)] - f^* \leq \frac{\|\theta_0 - \theta^*\|^2 + \eta_0^2(\sigma^2 + G^2)(1 + \log T)}{2\eta_0\sqrt{T}}
	\]
	Esto establece la tasa \(O(\log(T)/\sqrt{T})\).
\end{proof}

La tasa \(O(\log(T)/\sqrt{T})\) de \gls{SGD} es más lenta que la tasa \(O(1/T)\) de \gls{GD} determinista, pero \gls{SGD} procesa \(T\) veces menos datos por iteración, lo que puede resultar en mejor eficiencia computacional total.

\begin{theorem}[Convergencia a puntos estacionarios (caso no convexo)]\label{thm:gd_nonconvex}
	Sea \(f: \mathbb{R}^n \to \mathbb{R}\) una función \(L\)-suave (no necesariamente convexa) y acotada inferiormente por \(f^*\). Con tasa de aprendizaje \(\eta \leq 1/L\), el descenso del gradiente satisface:
	\[
		\min_{t=0,\ldots,T-1} \|\nabla f(\theta_t)\|^2 \leq \frac{2(f(\theta_0) - f^*)}{\eta T}
	\]
	Es decir, \gls{GD} encuentra un punto \(\epsilon\)-estacionario (con \(\|\nabla f\| \leq \epsilon\)) en \(O(1/\epsilon^2)\) iteraciones.
\end{theorem}

\begin{proof}
	Utilizamos la caracterización de \(L\)-suavidad de la \autoref{def:L_smooth}.

	\textbf{Paso 1: Desigualdad de descenso.}
	Por \(L\)-suavidad, para la actualización \(\theta_{t+1} = \theta_t - \eta \nabla f(\theta_t)\):
	\begin{align*}
		f(\theta_{t+1}) &\leq f(\theta_t) + \langle \nabla f(\theta_t), \theta_{t+1} - \theta_t \rangle + \frac{L}{2}\|\theta_{t+1} - \theta_t\|^2 \\
		&= f(\theta_t) - \eta \|\nabla f(\theta_t)\|^2 + \frac{L\eta^2}{2}\|\nabla f(\theta_t)\|^2 \\
		&= f(\theta_t) - \eta\left(1 - \frac{L\eta}{2}\right)\|\nabla f(\theta_t)\|^2
	\end{align*}
	Si \(\eta \leq 1/L\), entonces \(1 - \frac{L\eta}{2} \geq \frac{1}{2}\), por lo que:
	\[
		f(\theta_{t+1}) \leq f(\theta_t) - \frac{\eta}{2}\|\nabla f(\theta_t)\|^2
	\]

	\textbf{Paso 2: Suma telescópica.}
	Sumando sobre \(t = 0, \ldots, T-1\):
	\[
		\sum_{t=0}^{T-1} \frac{\eta}{2}\|\nabla f(\theta_t)\|^2 \leq \sum_{t=0}^{T-1} \left(f(\theta_t) - f(\theta_{t+1})\right) = f(\theta_0) - f(\theta_T) \leq f(\theta_0) - f^*
	\]
	donde usamos que \(f(\theta_T) \geq f^*\) por la cota inferior.

	\textbf{Paso 3: Conclusión.}
	Dividiendo ambos lados por \(\eta/2\):
	\[
		\sum_{t=0}^{T-1} \|\nabla f(\theta_t)\|^2 \leq \frac{2(f(\theta_0) - f^*)}{\eta}
	\]
	Dado que cada término \(\|\nabla f(\theta_t)\|^2 \geq 0\), el valor mínimo de la colección no puede superar su promedio aritmético (de lo contrario, todos los términos superarían el promedio, y su suma excedería \(T\) veces el promedio, contradiciendo su definición):
	\[
		\min_{t=0,\ldots,T-1} \|\nabla f(\theta_t)\|^2 \leq \frac{1}{T}\sum_{t=0}^{T-1} \|\nabla f(\theta_t)\|^2
	\]
	Combinando con la cota anterior, dividida entre \(T\):
	\[
		\min_{t=0,\ldots,T-1} \|\nabla f(\theta_t)\|^2 \leq \frac{1}{T}\sum_{t=0}^{T-1} \|\nabla f(\theta_t)\|^2 \leq \frac{2(f(\theta_0) - f^*)}{\eta T}
	\]
\end{proof}

Las funciones de pérdida de redes neuronales son típicamente no convexas debido a las no linealidades. El \autoref{thm:gd_nonconvex} garantiza convergencia a puntos estacionarios, pero no a mínimos globales. Sin embargo, en la práctica la sobreparametrización de redes modernas crea muchos mínimos equivalentes, el ruido de \gls{SGD} ayuda a escapar de puntos de silla, y los mínimos locales encontrados suelen generalizar bien.

\subsection{Optimizadores avanzados}\label{subsec:optimizadores_avanzados}

Los optimizadores modernos mejoran el \gls{SGD} básico incorporando información histórica de los gradientes~\citep{kingma2014adam, duchi2011adagrad}. En esta sección, todas las operaciones sobre vectores de parámetros y gradientes se aplican elemento a elemento: utilizamos \(\odot\) para el producto de Hadamard y \(\oslash\) para la división elemento a elemento, mientras que las operaciones de raíz cuadrada y suma con escalares también se aplican componente a componente.

\subsubsection{SGD con Momentum}

\begin{definition}[SGD con Momentum]\label{def:sgd_momentum}
	El \textit{SGD con momentum}~\citep{polyak1964methods} introduce una variable de velocidad \(v_t\) que acumula gradientes pasados, con \(v_0 = 0\):
	\begin{align*}
		v_t &= \beta v_{t-1} + \nabla_\theta \mathcal{L}(\theta_{t-1}) \\
		\theta_t &= \theta_{t-1} - \eta v_t
	\end{align*}
	donde \(\beta \in [0,1)\) es el \textit{coeficiente de momentum} (típicamente \(\beta = 0.9\)).
\end{definition}

El momentum acumula gradientes pasados ponderados exponencialmente, suavizando las oscilaciones en direcciones de alta curvatura y acelerando la convergencia en direcciones consistentes. Nótese que la fórmula de momentum \textit{no} es una media móvil exponencial (EMA), ya que carece del factor \((1-\beta)\): mientras que una EMA normaliza para que los pesos sumen uno, el acumulador de momentum crece con factor \(1/(1-\beta)\) en régimen estacionario, y es la tasa de aprendizaje \(\eta\) la que absorbe esta escala.

\subsubsection{Métodos adaptativos}

\begin{definition}[Adagrad]\label{def:adagrad}
	El optimizador \textit{Adagrad}~\citep{duchi2011adagrad} adapta la tasa de aprendizaje por parámetro según la historia de gradientes, con \(G_0 = 0\):
	\begin{align*}
		G_t &= G_{t-1} + g_t \odot g_t \\
		\theta_t &= \theta_{t-1} - \eta \cdot g_t \oslash (\sqrt{G_t} + \epsilon)
	\end{align*}
	donde \(g_t = \nabla_\theta \mathcal{L}(\theta_{t-1})\) y \(\epsilon \approx 10^{-8}\) evita división por cero. Parámetros con gradientes históricamente grandes reciben tasas de aprendizaje más pequeñas.
\end{definition}

\begin{definition}[RMSProp]\label{def:rmsprop}
	\textit{RMSProp}~\citep{hinton2012rmsprop} corrige el decaimiento monótono de Adagrad usando una media móvil exponencial del cuadrado de los gradientes. Denotamos por \(v_t\) el acumulador del segundo momento (reutilizando el símbolo de la velocidad en momentum, siguiendo la convención habitual en la literatura), con \(v_0 = 0\):
	\begin{align*}
		v_t &= \beta v_{t-1} + (1-\beta) (g_t \odot g_t) \\
		\theta_t &= \theta_{t-1} - \eta \cdot g_t \oslash (\sqrt{v_t} + \epsilon)
	\end{align*}
	donde típicamente \(\beta = 0.9\). A diferencia de Adagrad, el denominador no crece indefinidamente, permitiendo aprendizaje sostenido en problemas no estacionarios.
\end{definition}

\begin{definition}[Adam]\label{def:adam}
	El optimizador \textit{Adam}~\citep{kingma2014adam} (Adaptive Moment Estimation) combina momentum con tasas adaptativas, manteniendo estimaciones de los dos primeros momentos del gradiente, con \(m_0 = v_0 = 0\):
	\begin{align*}
		m_t &= \beta_1 m_{t-1} + (1-\beta_1) g_t & \text{(primer momento)} \\
		v_t &= \beta_2 v_{t-1} + (1-\beta_2) (g_t \odot g_t) & \text{(segundo momento)} \\
		\hat{m}_t &= m_t / (1 - \beta_1^t) & \text{(corrección de sesgo)} \\
		\hat{v}_t &= v_t / (1 - \beta_2^t) & \\
		\theta_t &= \theta_{t-1} - \eta \cdot \hat{m}_t \oslash (\sqrt{\hat{v}_t} + \epsilon)
	\end{align*}
	donde típicamente \(\beta_1 = 0.9\), \(\beta_2 = 0.999\), \(\epsilon = 10^{-8}\).
\end{definition}

Las correcciones \(\hat{m}_t\) y \(\hat{v}_t\) compensan el sesgo hacia cero de las medias móviles exponenciales cuando se inicializan en cero; para \(t\) grande, \((1 - \beta^t) \approx 1\) y la corrección es negligible. Cabe aclarar que los términos ``primer momento'' y ``segundo momento'' en \gls{Adam} se refieren a \textit{medias móviles exponenciales} del gradiente y de su cuadrado, no a momentos probabilísticos en el sentido estricto. La terminología proviene de la analogía con la estimación de la media (\(\mathbb{E}[g_t]\)) y la varianza no centrada (\(\mathbb{E}[g_t^2]\)) del gradiente estocástico, pero \(m_t\) y \(v_t\) son estimadores sesgados que requieren corrección.

La actualización efectiva en \gls{Adam} para cada parámetro \(\theta_i\) es aproximadamente \(\Delta\theta_i \approx \eta \cdot \hat{m}_{t,i} / (\sqrt{\hat{v}_{t,i}} + \epsilon)\), lo que proporciona pasos más pequeños para parámetros con gradientes históricamente grandes y pasos relativamente mayores para parámetros con gradientes pequeños pero consistentes.

\section{Deep Learning}

\subsection{Del perceptrón multicapa al aprendizaje profundo}

El interés por los modelos neuronales comenzó cuando Rosenblatt demostró que el perceptrón podía aprender a clasificar patrones a partir de datos~\citep{rosenblatt1958perceptron}. Sin embargo, Minsky y Papert~\citep{minsky1969perceptrons} probaron que el perceptrón simple solo puede representar funciones linealmente separables, lo que reveló limitaciones fundamentales del modelo y provocó un prolongado periodo de reducción del interés y la financiación en este campo.

La introducción del algoritmo de retropropagación~\citep{backpropagation} proporcionó un método eficiente para entrenar redes multicapa mediante descenso del gradiente. No obstante, este avance resultó insuficiente por sí solo: las barreras computacionales y los problemas inherentes al entrenamiento de redes con muchas capas persistieron durante años.

En las secciones anteriores hemos visto herramientas matemáticas para entrenar redes neuronales. Sin embargo, durante décadas las redes profundas permanecieron impracticables debido al problema del desvanecimiento del gradiente analizado en la \autoref{subsubsec:vanishing_gradient}: a medida que se añaden capas, el producto de derivadas decrece exponencialmente, impidiendo que las primeras capas aprendan. La alternativa de incrementar la amplitud de la red (Ver el \autoref{thm:aproximacion_universal}) era computacionalmente demasiado costosa para realmente ver la utilidad de las \glspl{ANN} en problemas medianamente complejos.

El \gls{DL} o aprendizaje profundo surge como la convergencia de múltiples avances que superaron estas limitaciones:
\begin{itemize}
	\item Técnicas de inicialización que preservan la varianza de las activaciones a través de las capas.
	\item Funciones de activación como \gls{ReLU} (véase \autoref{subsubsec:funciones_activacion_capas}) que evitan la saturación y mantienen gradientes no nulos.
	\item Métodos de regularización como Dropout (\autoref{def:dropout}) que controlan el sobreajuste en redes de alta capacidad.
	\item Arquitecturas especializadas que reducen drásticamente el coste computacional explotando la estructura de los datos. Por ejemplo, las \glspl{CNN} aprovechan la localidad espacial para requerir órdenes de magnitud menos parámetros que las redes densas equivalentes.
\end{itemize}

La cristalización de estos avances y su impacto práctico se documentan en el \autoref{ch:cnn}, donde se analiza cómo las redes convolucionales materializaron la revolución del aprendizaje profundo.

Estas redes más profundas tienen mayor capacidad expresiva, pero requieren más datos y recursos computacionales. Una técnica habitual para suplir la escasez de datos es el \textit{aumento de datos}, que incrementa artificialmente el conjunto de entrenamiento aplicando transformaciones que preservan las etiquetas.

\begin{definition}[Aumento de datos]\label{def:data_augmentation}
	Sea $\{(\mathbf{x}_i, y_i)\}_{i=1}^N$ un conjunto de entrenamiento y $G$ un grupo de transformaciones que actúa sobre el espacio de entrada $\mathcal{X}$ preservando las etiquetas, es decir, tales que la etiqueta correcta de $g \cdot \mathbf{x}$ coincide con la de $\mathbf{x}$ para todo $g \in G$. El \textit{aumento de datos} (\emph{data augmentation}) consiste en reemplazar el conjunto original por el conjunto aumentado
	\[
		\{(g \cdot \mathbf{x}_i,\; y_i) : (\mathbf{x}_i, y_i) \in \{(\mathbf{x}_j, y_j)\}_{j=1}^N,\; g \in G\},
	\]
	cuyo cardinal es $|G| \cdot N$. En la práctica, se muestrea un subconjunto finito de $G$ en cada época de entrenamiento.
\end{definition}

El aumento de datos actúa como una forma de regularización implícita: fuerza al modelo a aprender representaciones invariantes respecto al grupo $G$. Sin embargo, como se analizará en el \autoref{subsec:equivarianza_invarianza}, esta estrategia es ineficiente comparada con las arquitecturas equivariantes del \gls{GDL}, que incorporan las simetrías directamente en su estructura.

A continuación analizamos las ventajas teóricas que justifican el uso de arquitecturas profundas.

\subsubsection{Ventajas de la profundidad}

Las redes profundas ofrecen varias ventajas fundamentales sobre las redes superficiales anchas, respaldadas tanto por resultados teóricos como empíricos:

\begin{enumerate}
	\item \textbf{Eficiencia de representación}:

	Existen funciones que exhiben una separación exponencial entre redes profundas y superficiales. Un ejemplo clásico es la función de paridad:

	\begin{example}[Función de paridad]
		La función de paridad \(f: \{0,1\}^n \to \{0,1\}\) definida como:
		\[
			f(x_1, \ldots, x_n) = \bigoplus_{i=1}^n x_i = \sum_{i=1}^n x_i \mod 2
		\]
		requiere \(O(2^n)\) neuronas en una red de 2 capas, pero solo \(O(n)\) neuronas en una red de \(O(\log n)\) capas.
	\end{example}

	\begin{theorem}[Separación exponencial~\citep{telgarsky2016benefitsdepthneuralnetworks}]
		Para cualquier entero \(k \geq 1\), existen funciones computables por redes de profundidad \(\Theta(k^3)\) con \(\Theta(1)\) neuronas por capa, que cualquier red de profundidad \(O(k)\) necesitaría \(\Omega(2^k)\) neuronas para aproximar con error constante.
	\end{theorem}

	\begin{proof}[Esbozo de la demostración]
		Definimos la función triangular base usando neuronas ReLU:
		\[
			\varphi(x) = 2\operatorname{ReLU}(x) - 4\operatorname{ReLU}\!\left(x - \tfrac{1}{2}\right) + 2\operatorname{ReLU}(x - 1),
		\]
		que forma una ``tienda'' en $[0,1]$ con un pico en $x = \tfrac{1}{2}$. Consideramos las composiciones iteradas $\varphi_1 = \varphi$, $\varphi_{k+1} = \varphi \circ \varphi_k$. La observación clave es que cada composición duplica el número de oscilaciones: $\varphi_k$ tiene exactamente $2^k$ oscilaciones en $[0,1]$, ya que $\varphi$ mapea $[0,1] \to [0,1]$ y cada arco de $\varphi_k$ genera dos arcos en $\varphi_{k+1}$.

		Puesto que cada composición requiere $O(1)$ neuronas adicionales, la función $\varphi_k$ es computable por una red de profundidad $O(k)$ y anchura constante. Sin embargo, una red con activación ReLU y $p$ neuronas totales produce una función lineal a trozos con a lo sumo $O(p^L)$ segmentos para profundidad $L$. Para una red de profundidad $O(k)$, aproximar $2^k$ oscilaciones requiere al menos $\Omega(2^k)$ neuronas. La construcción se extiende al caso $\Theta(k^3)$ vs $O(k)$ del enunciado mediante técnicas más refinadas de conteo de segmentos lineales. Véase \citet{telgarsky2016benefitsdepthneuralnetworks} para la demostración completa.
	\end{proof}

	La \autoref{fig:profundidad_vs_anchura} ilustra visualmente esta separación exponencial entre redes superficiales y profundas.

	\begin{figure}[h]
	\centering
	\begin{tikzpicture}[
	    neuron/.style={circle, draw, minimum size=0.3cm, fill=gdlblue!20},
	    arrow/.style={->, thick, >=stealth}
	]

	% Red superficial (izquierda)
	\node[font=\bfseries] at (-3, 3) {Red superficial (2 capas)};
	\node[font=\footnotesize] at (-3, 2.5) {$O(2^n)$ neuronas};

	% Entrada
	\foreach \i in {1,2,3} {
	    \node[neuron, fill=gdlgreen!30] (si\i) at (-4.5, 1.5-\i*0.5) {};
	}
	\node at (-4.5, -0.7) {$\vdots$};

	% Capa oculta (muchas neuronas)
	\foreach \i in {1,...,8} {
	    \node[neuron] (sh\i) at (-3, 2-\i*0.4) {};
	}
	\node at (-3, -1.5) {$\vdots$};
	\node[font=\scriptsize] at (-3, -2) {$2^n$};

	% Salida
	\node[neuron, fill=gdlred!30] (so) at (-1.5, 0) {};

	% Conexiones (simplificadas)
	\draw[gdlgray, thin] (si1) -- (sh1);
	\draw[gdlgray, thin] (si2) -- (sh4);
	\draw[gdlgray, thin] (si3) -- (sh7);
	\draw[gdlgray, thin] (sh1) -- (so);
	\draw[gdlgray, thin] (sh8) -- (so);

	% Red profunda (derecha)
	\node[font=\bfseries] at (3, 3) {Red profunda ($\log n$ capas)};
	\node[font=\footnotesize] at (3, 2.5) {$O(n)$ neuronas};

	% Capas
	\foreach \l in {0,...,4} {
	    \foreach \i in {1,2,3} {
	        \pgfmathsetmacro{\col}{ifthenelse(\l==0,"gdlgreen!30",ifthenelse(\l==4,"gdlred!30","gdlblue!20"))}
	        \node[neuron, fill=\col] (d\l\i) at (1.5+\l*0.8, 1-\i*0.6) {};
	    }
	}

	% Conexiones entre capas
	\foreach \l in {0,...,3} {
	    \pgfmathtruncatemacro{\nextl}{\l+1}
	    \foreach \i in {1,2,3} {
	        \foreach \j in {1,2,3} {
	            \draw[gdlgray, very thin] (d\l\i) -- (d\nextl\j);
	        }
	    }
	}

	% Etiquetas de complejidad
	\node[font=\small, align=center] at (-3, -3) {Complejidad:\\exponencial};
	\node[font=\small, align=center] at (3, -3) {Complejidad:\\polinomial};

	\end{tikzpicture}
	\caption[Profundidad vs anchura en redes neuronales]{Comparación entre red superficial ancha y red profunda estrecha para computar la misma función. La red superficial requiere exponencialmente más neuronas, mientras que la red profunda logra la misma expresividad con crecimiento polinomial.}
	\label{fig:profundidad_vs_anchura}
	\end{figure}

	\item \textbf{Aprendizaje jerárquico y composicionalidad}:

	Las redes profundas explotan la estructura composicional natural de muchos problemas:

	\begin{itemize}
		\item \textbf{Visión}: píxeles \(\to\) bordes \(\to\) texturas \(\to\) partes \(\to\) objetos
		\item \textbf{Lenguaje}: caracteres \(\to\) morfemas \(\to\) palabras \(\to\) frases \(\to\) semántica
		\item \textbf{Audio}: muestras \(\to\) fonemas \(\to\) palabras \(\to\) conceptos
	\end{itemize}

	Esta jerarquía se refleja matemáticamente en la composición de funciones:
	\[
		f(x) = f_L \circ f_{L-1} \circ \cdots \circ f_1(x)
	\]
	donde cada \(f_i\) aprende transformaciones incrementalmente más abstractas.

	\item \textbf{Reutilización de características y eficiencia paramétrica}:

	La profundidad permite compartir cómputos entre diferentes rutas de procesamiento. Consideremos una red que debe computar \(m\) funciones diferentes \(\{g_1, \ldots, g_m\}\):

	\begin{itemize}
		\item Red superficial: \(O(m \cdot n)\) parámetros (cada función independiente)
		\item Red profunda: \(O(n + m)\) parámetros (características compartidas + cabezas específicas)
	\end{itemize}

	\item \textbf{Sesgo inductivo implícito}:

	Las redes profundas tienen un sesgo implícito hacia funciones ``simples'' en el sentido de la complejidad de Kolmogorov:

	\begin{definition}[Complejidad de la función]\label{def:complejidad_funcion}
		Para una red neuronal \(f_\theta\) con arquitectura fija, definimos su complejidad como:
		\[
			\mathcal{C}(f_\theta) = \|\theta\|_0 + \sum_{l=1}^L \log(\text{fan-in}_l \cdot \text{fan-out}_l)
		\]
	\end{definition}

	Las redes profundas favorecen soluciones con menor \(\mathcal{C}(f_\theta)\), lo que corresponde a funciones más ``naturales'' según la distribución de datos del mundo real.

	\item \textbf{Optimización y geometría de la función de pérdida}:

	Contraintuitivamente, las redes más profundas pueden ser más fáciles de optimizar:

	\begin{itemize}
		\item Los mínimos locales en redes sobreparametrizadas tienden a tener error similar al global
		\item La profundidad induce una estructura jerárquica en la geometría de la función de pérdida
		\item Fenómeno de ``lazy training'': las redes profundas encuentran soluciones simples primero
	\end{itemize}
\end{enumerate}

\subsection{Arquitecturas modernas}

El desarrollo del deep learning ha producido arquitecturas especializadas, cada una diseñada para explotar estructuras específicas en los datos. Muchas de estas arquitecturas ya se han introducido en secciones anteriores de este capítulo; aquí ofrecemos una visión comparativa enfocada en sus características distintivas y en cómo se relacionan entre sí. En capítulos posteriores exploraremos en profundidad cómo las ideas geométricas transforman algunas de estas arquitecturas.

La \autoref{fig:arquitecturas_modernas} presenta una comparación visual de las cuatro arquitecturas principales del deep learning, mostrando sus características distintivas en términos de flujo de información y estructura de conectividad.

\begin{figure}[h]
\centering
\begin{tikzpicture}[scale=1.05, every node/.style={scale=1.0}]
% Estilos comunes
\tikzset{
    neuron/.style={circle, draw, minimum size=0.4cm, fill=#1!40},
    layer/.style={rectangle, draw, rounded corners, minimum width=0.8cm, minimum height=1.8cm, fill=#1!20},
    kernel/.style={rectangle, draw, minimum width=0.3cm, minimum height=0.3cm, fill=gdlorange!50},
    arrow/.style={->, thick, >=stealth},
    title/.style={font=\small\bfseries, above}
}

% === FEEDFORWARD ===
\begin{scope}[shift={(-6,0)}]
    \node[title] at (0.6,2.5) {Feedforward};
    % Capas
    \foreach \y in {0,0.6,1.2,1.8} {
        \foreach \x in {0,0.4,0.8,1.2} {
            \node[neuron=gdlblue] at (\x,\y) {};
        }
    }
    % Conexiones (algunas representativas)
    \foreach \x in {0,0.4,0.8,1.2} {
        \draw[gdlgray!50, thin] (\x,0.15) -- (0.2,0.45);
        \draw[gdlgray!50, thin] (\x,0.15) -- (0.6,0.45);
        \draw[gdlgray!50, thin] (\x,0.15) -- (1.0,0.45);
    }
    \draw[arrow, gdlblue!70] (0.6,-0.3) -- (0.6,0) node[below, pos=0, font=\tiny] {entrada};
    \draw[arrow, gdlblue!70] (0.6,1.95) -- (0.6,2.2) node[above, pos=1, font=\tiny] {salida};
    \node[below, font=\tiny, text=gdlgray] at (0.6,-0.6) {Flujo unidireccional};
\end{scope}

% === CNN ===
\begin{scope}[shift={(-2,0)}]
    \node[title] at (0.8,2.5) {CNN};
    % Feature maps (rectángulos apilados)
    \draw[fill=gdlgreen!20, draw=gdlgreen!60] (0,0) rectangle (1.2,1.6);
    \draw[fill=gdlgreen!30, draw=gdlgreen!60] (0.15,0.15) rectangle (1.05,1.35);
    \draw[fill=gdlgreen!40, draw=gdlgreen!60] (0.3,0.3) rectangle (0.9,1.1);
    % Kernel
    \draw[fill=gdlorange!50, draw=gdlorange!80] (0.4,0.5) rectangle (0.7,0.8);
    \node[font=\tiny] at (0.55,0.65) {$k$};
    % Flecha de convolución
    \draw[arrow, gdlorange!80, dashed] (0.55,0.9) -- (0.55,1.5);
    \draw[arrow, gdlgreen!70] (0.6,-0.3) -- (0.6,0) node[below, pos=0, font=\tiny] {imagen};
    \draw[arrow, gdlgreen!70] (0.6,1.7) -- (0.6,2.2) node[above, pos=1, font=\tiny] {features};
    \node[below, font=\tiny, text=gdlgray] at (0.6,-0.6) {Compartición pesos};
\end{scope}

% === RNN ===
\begin{scope}[shift={(2,0)}]
    \node[title] at (0.6,2.5) {RNN};
    % Celdas en secuencia
    \foreach \i/\x/\t in {0/0/$t{-}1$, 1/0.8/$t$, 2/1.6/$t{+}1$} {
        \node[layer=gdlred, minimum height=1cm, minimum width=0.5cm] (rnn\i) at (\x,0.8) {};
        \node[font=\tiny] at (\x,0.8) {\t};
    }
    % Conexiones recurrentes
    \draw[arrow, gdlred!70] (0.3,0.8) -- (0.5,0.8);
    \draw[arrow, gdlred!70] (1.1,0.8) -- (1.3,0.8);
    % Loop
    \draw[arrow, gdlred!70, thick] (rnn1.north) .. controls (0.8,1.6) and (1.2,1.6) .. (rnn2.north);
    \node[font=\tiny, gdlred!70] at (1.2,1.8) {$h_t$};
    % Entradas
    \foreach \x in {0,0.8,1.6} {
        \draw[arrow, gdlgray] (\x,0) -- (\x,0.3);
    }
    \node[below, font=\tiny] at (0.8,-0.3) {secuencia};
    \node[below, font=\tiny, text=gdlgray] at (0.8,-0.6) {Estado oculto};
\end{scope}

% === TRANSFORMER ===
\begin{scope}[shift={(6,0)}]
    \node[title] at (0.8,2.5) {Transformer};
    % Tokens
    \foreach \x in {0,0.5,1,1.5} {
        \node[neuron=gdlpurple, minimum size=0.3cm] (t\x) at (\x,0) {};
    }
    % Attention layer
    \draw[fill=gdlpurple!20, draw=gdlpurple!60] (-0.2,0.6) rectangle (1.7,1.2);
    \node[font=\tiny] at (0.75,0.9) {Atención};
    % Attention connections (all-to-all)
    \foreach \x in {0,0.5,1,1.5} {
        \draw[gdlpurple!40, thin] (\x,0.15) -- (0.75,0.55);
    }
    % Output
    \foreach \x in {0,0.5,1,1.5} {
        \node[neuron=gdlpurple, minimum size=0.3cm] at (\x,1.6) {};
        \draw[gdlpurple!40, thin] (0.75,1.25) -- (\x,1.45);
    }
    \draw[arrow, gdlpurple!70] (0.75,-0.5) -- (0.75,-0.15) node[below, pos=0, font=\tiny] {tokens};
    \draw[arrow, gdlpurple!70] (0.75,1.75) -- (0.75,2.2) node[above, pos=1, font=\tiny] {salida};
    \node[below, font=\tiny, text=gdlgray] at (0.75,-0.8) {Atención global};
\end{scope}

\end{tikzpicture}
\caption[Comparación de arquitecturas de deep learning]{Comparación de las cuatro arquitecturas principales de deep learning. \textbf{Feedforward}: capas completamente conectadas con flujo unidireccional. \textbf{CNN}: convoluciones con kernels compartidos que preservan estructura espacial. \textbf{RNN}: conexiones recurrentes que mantienen estado oculto temporal. \textbf{Transformer}: mecanismos de atención que conectan todos los elementos simultáneamente.}
\label{fig:arquitecturas_modernas}
\end{figure}

\paragraph{Redes feedforward.}
Las redes feedforward, basadas en capas densas (\autoref{def:capa_densa}), representan la arquitectura más básica donde la información fluye unidireccionalmente desde la entrada hasta la salida. Una red con $L$ capas se define como la composición $f(x) = f_L \circ f_{L-1} \circ \cdots \circ f_1(x)$, donde cada $f_i(x) = \sigma(W_i x + b_i)$ es una transformación afín seguida de una no-linealidad. La ausencia de ciclos garantiza que el cómputo sea determinista y eficiente.

\paragraph{Redes convolucionales.}
Las \glspl{CNN}, que estudiaremos en detalle en el \autoref{ch:cnn}, explotan la estructura local y la invarianza traslacional mediante operaciones de convolución. Para una entrada $x$ y un kernel $k$, la convolución discreta (\autoref{def:discrete_convolution}) se define como:
\[
	(x * k)[i] = \sum_j x[j] \cdot k[i-j].
\]
Las capas convolucionales preservan la localidad espacial y reducen drásticamente el número de parámetros mediante compartición de pesos, resultando en equivarianza traslacional: $T_a[f(x)] = f(T_a[x])$ donde $T_a$ es una traslación (\autoref{theorem:convolution_translation_equivariance}).

\paragraph{Redes recurrentes.}
Las \glspl{RNN} (\autoref{subsubsec:rnn}) procesan secuencias mediante un estado oculto $h_t$ que evoluciona según:
\[
	h_t = \sigma(W_{hh} h_{t-1} + W_{xh} x_t + b_h).
\]
Esta recurrencia permite capturar dependencias temporales, aunque sufre del problema del desvanecimiento del gradiente (\autoref{subsubsec:vanishing_gradient}). Variantes como \gls{LSTM} (\autoref{subsubsec:lstm}) y \gls{GRU} introducen mecanismos de compuerta que permiten el aprendizaje de dependencias a largo plazo.

\paragraph{Mecanismos de atención.}\label{par:atencion_arquitecturas}
Los mecanismos de atención, que formalizamos a continuación y extendemos en el \autoref{ch:attention}, permiten modelar dependencias de largo alcance sin recurrencia. La idea central es ponderar dinámicamente las contribuciones de diferentes partes de la entrada según su relevancia:
\[
	\text{Attention}(q, K, V) = \sum_{i=1}^n \alpha_i v_i, \quad \alpha_i = \frac{\exp(s(q, k_i))}{\sum_{j=1}^n \exp(s(q, k_j))}
\]
donde $s(q, k_i)$ es una función de puntuación que mide la relevancia entre la consulta $q$ y la clave $k_i$ (véase \autoref{def:generalized_attention}). Esta formulación es invariante a permutaciones de las entradas y generaliza operaciones como pooling y convolución. La complejidad cuadrática $O(n^2)$ en la longitud de secuencia ha motivado el desarrollo de variantes eficientes. Los mecanismos de atención forman la base de la arquitectura \gls{Transformer}.

\begin{definition}[Capa de Self-Attention]\label{def:self_attention_basica}
Sea $\{h_1, \ldots, h_n\} \subset \mathbb{R}^d$ una secuencia de vectores de entrada. Una \textit{capa de self-attention} con dimensión de clave $d_k$ se define mediante tres proyecciones lineales $Q, K, V \in \mathbb{R}^{d \times d_k}$ (query, key, value) y calcula:
\[
\mathrm{SelfAttention}(h_i) = \sum_{j=1}^n \alpha_{ij}\, V h_j, \quad \alpha_{ij} = \frac{\exp\!\left(\frac{(Qh_i)^T(Kh_j)}{\sqrt{d_k}}\right)}{\sum_{m=1}^n \exp\!\left(\frac{(Qh_i)^T(Kh_m)}{\sqrt{d_k}}\right)}
\]
donde los coeficientes $\alpha_{ij} \in (0,1)$ con $\sum_j \alpha_{ij} = 1$ se denominan \textit{pesos de atención}.
\end{definition}

\paragraph{Redes generativas adversarias.}
Las \glspl{GAN} consisten en dos redes en competencia: un generador $G: \mathcal{Z} \to \mathcal{X}$ y un discriminador $D: \mathcal{X} \to [0,1]$. El objetivo minimax es:
\[
	\min_G \max_D\; \mathbb{E}_{x \sim p_{\text{data}}}[\log D(x)] + \mathbb{E}_{z \sim p_z}[\log(1 - D(G(z)))].
\]
En el equilibrio de Nash, $G$ genera muestras indistinguibles de los datos reales.

\paragraph{Redes de grafos.}
Las \glspl{GNN} generalizan las operaciones neuronales a datos no euclidianos, un tema central del \autoref{ch:gdl}. La propagación de mensajes (\autoref{def:mpnn_unificador}) en un grafo $\mathcal{G} = (\mathcal{V}, \mathcal{E})$ se expresa como:
\[
	h_v^{(k+1)} = \sigma\!\left(W_{\text{self}}h_v^{(k)} + \sum_{u \in \mathcal{N}(v)} W_{\text{msg}}h_u^{(k)}\right),
\]
donde $\mathcal{N}(v)$ denota los vecinos del nodo $v$. Esta formulación preserva la permutación-invarianza del grafo.

\paragraph{Modelos de difusión.}
Los modelos de difusión aprenden a revertir un proceso de degradación gradual. Dado un proceso forward que añade ruido gaussiano:
\[
	q(x_t \mid x_{t-1}) = \mathcal{N}\!\left(x_t;\, \sqrt{1-\beta_t}\, x_{t-1},\, \beta_t \mathbf{I}\right),
\]
la red neuronal aprende el proceso inverso $p_\theta(x_{t-1} \mid x_t)$ para generar muestras desde ruido puro. La función objetivo está relacionada con la estimación del score $\nabla_x \log p(x)$.
